\documentclass[11pt,a4paper]{article}
\usepackage[T1]{fontenc}
\usepackage[utf8]{inputenc}
\usepackage{lmodern,amsmath,amssymb,amsthm,mathtools,microtype,mathrsfs}
\usepackage[margin=25mm,top=23mm,bottom=25mm]{geometry}
\usepackage{booktabs,array,longtable,enumitem,xcolor,fancyhdr}
\usepackage[colorlinks=true,linkcolor=blue!45!black,urlcolor=blue!45!black,citecolor=blue!45!black]{hyperref}
\hypersetup{pdftitle={Reading the Alpoge--Claude RH and Dirichlet-GRH proof: working version 0.46},pdfauthor={ChatGPT},pdfsubject={Partial human-readable reconstruction with explicit arithmetic inputs},pdfkeywords={Riemann hypothesis, Lean, theta lifts, prime sums, Arakelov slopes}}
\setlength{\parindent}{0pt}
\setlength{\parskip}{5pt plus 1pt}
\setlist{topsep=3pt,itemsep=3pt,leftmargin=*}
\setlength{\headheight}{14pt}
\pagestyle{fancy}\fancyhf{}
\fancyhead[L]{\small RH and Dirichlet GRH: a reconstruction}
\fancyhead[R]{\small Working version 0.46}
\fancyfoot[C]{\thepage}
\renewcommand{\headrulewidth}{0.3pt}
\numberwithin{equation}{section}
\newtheorem{theorem}{Theorem}[section]
\newtheorem{proposition}[theorem]{Proposition}
\newtheorem{lemma}[theorem]{Lemma}
\newtheorem{corollary}[theorem]{Corollary}
\theoremstyle{remark}\newtheorem{remark}[theorem]{Remark}
\newcommand{\Rea}{\operatorname{Re}}
\newcommand{\HS}{\mathrm{HS}}
\newcommand{\tr}{\operatorname{tr}}
\newcommand{\codim}{\operatorname{codim}}
\newcommand{\covol}{\operatorname{covol}}
\newcommand{\OK}{\mathcal O_K}
\newcommand{\src}[2]{\href{https://github.com/CocoSgt/riemann-hypothesis-lean/blob/7cbf7a0c32fca1badf05931879a37c3f8a2b03b4/src/#1}{\texttt{#2}}}
\newcommand{\status}[1]{\textbf{Status:} #1}
\begin{document}
\begin{center}
{\LARGE\bfseries RH and Dirichlet GRH: new reconstruction chapters}\\[7pt]
{\large Version 0.46 --- 23 November 2026}\\[5pt]
{\large Congruence integers, arithmetic rows, and local gain}
\end{center}
\textbf{Scope.} This edition contains the current audit and the new
chapters. It remains a partial human reconstruction. Earlier named
claims are in the cumulative manuscript and preserved component files.
The underlying formal proof is accepted as correct; the exact remaining
human reconstruction boundaries are stated below.
\setcounter{tocdepth}{1}
\tableofcontents
\clearpage
\begingroup\emergencystretch=3em
\section{Version 0.46: completed congruence and counting arguments, with the exact remaining inputs}
\label{sec:dependency-audit-v46}

\subsection{Current result and the two completed selected contracts}

The pinned formal source is accepted as fully checked and correct.
The task here is to reconstruct its mathematical argument in readable
prose, preserving each original contract or replacing an intermediate
route by a proved simplification with the same conclusion.
Version 0.46 adds 91 numbered claims to the 2677 preserved predecessor
claims, for 2768 claims in all. The mathematical dependency graph
records this human reconstruction; it is not a mechanically extracted
transitive Lean dependency audit.

The full \texttt{CNR3} contract is now reconstructed. For each fixed
imaginary quadratic field, its constants precede all varying weights,
levels, primes and finite indexed root-of-unity orbits. The nonzero
integer precedes every subsequent coefficient prime, closure, embedding,
Frobenius lift and pairwise-nonisomorphic admissible collection. The
proof includes the literal disjunction between the original rank-three
and rank-two-block predicates, and permits the coefficient prime 3.

The distinct class-B task \texttt{TASK\_\allowbreak PROP\_\allowbreak GK\_\allowbreak INSTALL} is also
completed: the original finite orbit, Frobenius and valuation argument
proves \(\mathrm{CNR3}\land\mathrm{OneChar}\Rightarrow
\mathrm{PropGK}\). The separate arithmetic producer
\texttt{TASK\_\allowbreak ONE\_\allowbreak CHAR\_\allowbreak PROOF} remains open. A completed implication
does not assert the existence or truth of an unconstructed premise.
The surrounding unconditional Case-I producer is therefore not declared
complete by this installation.

The selected sufficient route still has 121 nodes and 161 edges.
Its whole-obligation register decreases from 30 to 28: 20 class C and
8 class B. There are 21 shared RH/Dirichlet-GRH obligations, 7 additional
Dirichlet-GRH obligations, and no RH-only obligations in this selected
route. A claim count measures written inventory, not the fraction of
the proof completed. A single whole obligation may contain many
independent substantial producers.

\subsection{What the five new chapters establish}

\begin{enumerate}
\item RNC constructs the full congruence integer. It adelizes the actual
rank-two newform over \(\mathbb Q\), identifies both required weights,
uses the original local-flag disjunction to control the coefficient-prime
level, and builds finite full cusp catalogues by fixed arithmetic cells.
The integral-power Satake convention makes all complex conjugations
compatible. Purity and integral endpoint bounds give a nonzero integral
factor for each orbit-label pair. Their complete product supplies the
same integer for every later admissible collection. The original
Galois level is retained; the auxiliary automorphic levels are used
only to bound the catalogues.

\item PGK proves the exact finite installation with the original
\(\mathrm{OneChar}\) length theorem as an explicit premise. It keeps
the character kernel, fixed-field constants and all valuation factors,
with the integer selected before the coefficient datum. RNC supplies
the congruence premise; it does not supply the length premise.

\item ARO reconstructs actual reduction, the precise Haar-fiber
normalization, the concrete level subgroup and character, and full
\(\mathrm{EtOne}\) and event algebra on the original supplied adelic
\(\mathrm{SG.HidaSetting}\). Schur compression identifies bounded
intertwiners with the exact projector image. A finite label completion
preserves the original good rows and repeated labels, constructing
one \(\mathrm{AllRows}\), \(\mathrm{AllJoinA}\) and
\(\mathrm{PCountTie}\). The count installer retains the actual
\(\mathrm{CountIdempotentsClass}\) theorem. The prior local weight
and retained-defect consumers then apply to this same object.

\item ICL proves finite generation of the whole original integral cusp
lattice by nonzero integral-frame evaluations at CM points. It does
not assume that this lattice spans. Given the supplied original \(\mathrm{IsNab}\) connection, the
identity cover, a Neumann
inverse, absolutely convergent unfolding, integral harmonic span,
coefficient descent to the original field and one global clearing
scalar give the complete conditional spanning route. The original
principal-chart, level-system, raw-Gaussian, short-rule/dictionary and
ramified-jet producers remain five explicit input groups, alongside that connection. The exact
\(\mathrm{LatFull}\) and both-level \(\mathrm{FWTie}\) consequences
are proved under them.

\item GLG proves the literal \(\mathrm{RestStar}\) and
\(\mathrm{RemainingCountingStatement}\), including both prime
ranges, all ceiling budgets and the original error
\(b_o=b^{x^{-\epsilon/2}}\). It supplies served-index and coverage
bounds, integral local operators, coefficient-rational descent,
\(\mathrm{NPDesc}\), normalized local lengths and the full
five-conjunct \(\mathrm{LinkG}\) installer. The latter retains the
actual common compatibility, full admitted family, local first-growth
and mean-deficit inputs; their existence is not inferred.
\end{enumerate}

The arithmetic \(\mathrm{SG.HidaSetting}\) in ARO is the original
adelic row setting. It is distinct from the ordinary-Hecke setting
used in the completed version-0.45 Hida chapter. The earlier six-field,
same-ideal and two-Link support results are retained without changing
their statements or reclassifying their foundations.

\subsection{Foundations, scope, preservation and review}

RNC1 is one checked published application, class A: the polarized
automorphic Galois existence theorem and its precisely matched
rank-three/rank-two coefficient-prime compatibility. The weight
identification is staged before the Shin-regular specialization.
The primary theorem references and hypothesis matches are recorded
in the component's classical-input registry. The other new claims
are class-B human proofs in their printed scope. Classical CM,
abelian-scheme, automorphic, cohomological and period foundations
remain explicitly named where used.

Every new claim is independently reviewed by an agent other than its
author. The review manifests bind both the whole proof file and each
reviewed block. Exact cited source bytes are compared with the pinned
git blobs, and the cited source files are distributed with the archive.
All predecessor claim records and the complete historical proof body
are preserved. Mutable predecessor metadata is frozen in
\texttt{history/\allowbreak version\_\allowbreak 0\_\allowbreak 45}.

Draft-mode TeX compilation checks the two complete readers for syntax,
cross-references and missing characters. It does not constitute PDF
visual review or a replay of the formal proof. No Lean build, kernel
replay, new formal-correctness audit or completed human proof of RH/GRH
is claimed. The original proofs' correctness is not reopened.

The following register contains every remaining whole selected
obligation and its exact current contract. The new completed inputs
are recorded separately from the inputs still to be constructed.

\subsection{The complete current selected register}

\paragraph{1. GLOBAL\_\allowbreak DENOMINATOR — C — Actual LinkN local table and binding producer}

Closed(Des $\mathcal C$ 3) LinkN: the actual etaSum is bounded with exact dZ0 and same-host vvTerm, subject only to the literal ShaFinOdd and summability guards.

\emph{Shared objects and choice order.} one FamP; one class $\mathcal A$ and its target-compatible subclass; one coefficient field per family datum; one x and i at each application; original selected rows, words and levels; thresholds chosen before the common family; actual D.nd; actual etaJ and Aj; sha\_\allowbreak eq

\emph{Source.} \href{https://github.com/CocoSgt/riemann-hypothesis-lean/blob/7cbf7a0c32fca1badf05931879a37c3f8a2b03b4/src/Formalism/SharpRestriction/GlobalTheoremInputs.lean#L45-L48}{LinkN}.

\emph{Source.} \href{https://github.com/CocoSgt/riemann-hypothesis-lean/blob/7cbf7a0c32fca1badf05931879a37c3f8a2b03b4/src/Formalism/SharpRestriction/GlobalTheoremInputs.lean#L40-L44}{LinkN\_\allowbreak at}.

\emph{Source.} \href{https://github.com/CocoSgt/riemann-hypothesis-lean/blob/7cbf7a0c32fca1badf05931879a37c3f8a2b03b4/src/Development/AutomorphicArithmetic/CharacterBoxes/DivinNOfHypV.lean#L117-L134}{divinN\_\allowbreak of\_\allowbreak hypV}.



\paragraph{2. GLOBAL\_\allowbreak RETAINED\_\allowbreak DEFECT — C — Original retained-row defect bound LinkIRR}

Rows.DelpNonneg and delp$\le$exp(-x\textasciicircum{}(ez/7)) uniformly for every late actual row datum.

\emph{Shared objects and choice order.} one FamP; one class $\mathcal A$ and its target-compatible subclass; one coefficient field per family datum; one x and i at each application; original selected rows, words and levels; thresholds chosen before the common family; AllRows A; LabelCountGoodInjective; AllJoinA; CountAllRows

\emph{New completed inputs.} ARO1-ARO30

\emph{Current remaining subfrontier.} Actual common X and all original ties/RowsFinDim, plus CountIdempotentsClass or the exact three sufficient arithmetic counts and simultaneous family thresholds. The same-object weight and retained-defect consumers are available.

\emph{Source.} \href{https://github.com/CocoSgt/riemann-hypothesis-lean/blob/7cbf7a0c32fca1badf05931879a37c3f8a2b03b4/src/Formalism/SharpRestriction/GlobalTheoremInputs.lean#L55-L58}{LinkIRR}.

\emph{Source.} \href{https://github.com/CocoSgt/riemann-hypothesis-lean/blob/7cbf7a0c32fca1badf05931879a37c3f8a2b03b4/src/Development/AbelianVarietyTheory/Proofs/GlobalInterpolation/LinkIRROfJoinEx.lean#L17-L38}{linkIRR\_\allowbreak of\_\allowbreak join\_\allowbreak ex}.

\emph{Source.} \href{https://github.com/CocoSgt/riemann-hypothesis-lean/blob/7cbf7a0c32fca1badf05931879a37c3f8a2b03b4/src/Development/AbelianVarietyTheory/Proofs/GlobalInterpolation/LinkIRROfJoinEx.lean#L143-L166}{link\_\allowbreak irr}.



\paragraph{3. GLOBAL\_\allowbreak PERIOD\_\allowbreak HEIGHT — C — Full actual period and height comparison LinkDIVt}

Closed conditional TuSize -> DGamSize -> (LinkEta -> DIVinN -> LinkDIV), with every original LinkDIV\_\allowbreak at clause.

\emph{Shared objects and choice order.} one FamP; one class $\mathcal A$ and its target-compatible subclass; one coefficient field per family datum; one x and i at each application; original selected rows, words and levels; thresholds chosen before the common family; same RR target; same Bu/besS; same Nu and Bk

\emph{Source.} \href{https://github.com/CocoSgt/riemann-hypothesis-lean/blob/7cbf7a0c32fca1badf05931879a37c3f8a2b03b4/src/Formalism/SharpRestriction/GlobalTheoremInputs.lean#L77-L90}{LinkDIV\_\allowbreak at}.

\emph{Source.} \href{https://github.com/CocoSgt/riemann-hypothesis-lean/blob/7cbf7a0c32fca1badf05931879a37c3f8a2b03b4/src/Formalism/SharpRestriction/GlobalTheoremInputs.lean#L94-L96}{LinkDIVt}.

\emph{Source.} \href{https://github.com/CocoSgt/riemann-hypothesis-lean/blob/7cbf7a0c32fca1badf05931879a37c3f8a2b03b4/src/Development/AbelianVarietyTheory/Proofs/GlobalInterpolation/LinkBrest.lean#L7-L25}{GlobalTheory.Top.link\_\allowbreak brest}.

\emph{Source.} \href{https://github.com/CocoSgt/riemann-hypothesis-lean/blob/7cbf7a0c32fca1badf05931879a37c3f8a2b03b4/src/Development/AbelianVarietyTheory/Proofs/GlobalInterpolation/WireSel.lean#L11-L41}{GlobalTheory.Top.wire\_\allowbreak sel}.



\paragraph{4. GLOBAL\_\allowbreak LOCAL\_\allowbreak GAIN — C — Original global LinkG installation}

One uniform LinkG with actual local lattice lengths, served-orbit coverage and MeanKappaBound; delta(x)->0 before varying actual indices.

\emph{Shared objects and choice order.} one FamP; one class $\mathcal A$ and its target-compatible subclass; one coefficient field per family datum; one x and i at each application; original selected rows, words and levels; thresholds chosen before the common family; actual G.Inst; actual GK; actual lenMN; same orbit representatives

\emph{New completed inputs.} GLG1-GLG18

\emph{Current remaining subfrontier.} Realize the original common compatibility and full admitted inputs, including local first-growth and mean-deficit bounds on the same family. The full counting theorem and all written local-to-global consumers are available.

\emph{Source.} \href{https://github.com/CocoSgt/riemann-hypothesis-lean/blob/7cbf7a0c32fca1badf05931879a37c3f8a2b03b4/src/Formalism/SharpRestriction/GlobalTheoremInputs.lean#L122-L137}{LinkG\_\allowbreak at}.

\emph{Source.} \href{https://github.com/CocoSgt/riemann-hypothesis-lean/blob/7cbf7a0c32fca1badf05931879a37c3f8a2b03b4/src/Formalism/SharpRestriction/GlobalTheoremInputs.lean#L138-L141}{LinkG}.

\emph{Source.} \href{https://github.com/CocoSgt/riemann-hypothesis-lean/blob/7cbf7a0c32fca1badf05931879a37c3f8a2b03b4/src/Development/AbelianVarietyTheory/Proofs/GlobalInterpolation/LinkG.lean#L14-L27}{GlobalTheory.Top.link\_\allowbreak g}.



\paragraph{5. GLOBAL\_\allowbreak SPECTRAL — C — Original normalized kernel/Gram estimates LinkS}

Closed TuSize -> LinkS with exact Gamma normalization, off-diagonal Hilbert-Schmidt bound, diagonal error and Dall/Dstar estimate, errors exp(-x\textasciicircum{}(epsilon/8)).

\emph{Shared objects and choice order.} one FamP; one class $\mathcal A$ and its target-compatible subclass; one coefficient field per family datum; one x and i at each application; original selected rows, words and levels; thresholds chosen before the common family; actual KD; same S.Cu spaces; GamOf; Dstar; Dall

\emph{Source.} \href{https://github.com/CocoSgt/riemann-hypothesis-lean/blob/7cbf7a0c32fca1badf05931879a37c3f8a2b03b4/src/Formalism/SharpRestriction/GlobalTheoremInputs.lean#L145-L149}{LinkS\_\allowbreak at}.

\emph{Source.} \href{https://github.com/CocoSgt/riemann-hypothesis-lean/blob/7cbf7a0c32fca1badf05931879a37c3f8a2b03b4/src/Formalism/SharpRestriction/GlobalTheoremInputs.lean#L152-L157}{LinkS}.

\emph{Source.} \href{https://github.com/CocoSgt/riemann-hypothesis-lean/blob/7cbf7a0c32fca1badf05931879a37c3f8a2b03b4/src/Development/AbelianVarietyTheory/Proofs/GlobalInterpolation/LinkAlt.lean#L16-L50}{GlobalTheory.Top.link\_\allowbreak alt}.



\paragraph{6. TASK\_\allowbreak ONE\_\allowbreak CHAR\_\allowbreak PROOF — C — Reconstruct the original one-character argument}

From exact InpZ and AutomorphyUnramifiedCoefficients, prove the original OneChar trace/length estimate on the supplied coefficient datum and GD input.

\emph{Source.} \href{https://github.com/CocoSgt/riemann-hypothesis-lean/blob/7cbf7a0c32fca1badf05931879a37c3f8a2b03b4/src/Development/AutomorphicArithmetic/AutomorphicVarieties/UnitaryThreeGroups/KatoUse.lean#L54-L64}{OneChar}.

\emph{Source.} \href{https://github.com/CocoSgt/riemann-hypothesis-lean/blob/7cbf7a0c32fca1badf05931879a37c3f8a2b03b4/src/Development/AutomorphicArithmetic/AutomorphicVarieties/ModularityArgument/KatoQuadraticCharacters/OneChar.lean#L158-L178}{UpperTheory.V.Z.Kato.oneChar}.



\paragraph{7. TASK\_\allowbreak INP\_\allowbreak Z\_\allowbreak ASSEMBLY — B — Install the exact Kato input record}

Build the original CohInp/InpZ record using exactly lemM,lemMprime,lemP,RigII,AdVanishK and number on the same arithmetic objects. No rank-four replacement is permitted.

\emph{Source.} \href{https://github.com/CocoSgt/riemann-hypothesis-lean/blob/7cbf7a0c32fca1badf05931879a37c3f8a2b03b4/src/Development/AutomorphicArithmetic/AutomorphicVarieties/ModularityArgument/KatoQuadraticCharacters/KatoTwoOfUseN.lean#L55-L61}{InpZ}.



\paragraph{8. RIG\_\allowbreak II — C — Type-II polarized adjoint rigidity}

Marked original datum, actual type-II point and modularity of all source extensions imply vanishing for its exact adRep.

\emph{Source.} \href{https://github.com/CocoSgt/riemann-hypothesis-lean/blob/7cbf7a0c32fca1badf05931879a37c3f8a2b03b4/src/Development/AutomorphicArithmetic/AutomorphicVarieties/ModularityArgument/KatoQuadraticCharacters/KatoTwoOfUseN.lean#L49-L53}{RigII}.

\emph{Source.} \href{https://github.com/CocoSgt/riemann-hypothesis-lean/blob/7cbf7a0c32fca1badf05931879a37c3f8a2b03b4/src/Development/AutomorphicArithmetic/AutomorphicVarieties/ModularityArgument/KatoQuadraticCharacters/KatoTwoOfUseN.lean#L43-L47}{KatoTwo}.

\emph{Source.} \href{https://github.com/CocoSgt/riemann-hypothesis-lean/blob/7cbf7a0c32fca1badf05931879a37c3f8a2b03b4/src/Development/AutomorphicArithmetic/AutomorphicVarieties/ModularityArgument/KatoQuadraticCharacters/KatoTwoOfUseN.lean#L11-L41}{Used}.



\paragraph{9. AUT\_\allowbreak UNRAM\_\allowbreak COEF — C — Actual unramified-coefficient automorphy}

For each actual Dat with coefficient ramification one and fixed iota, choose one marked g1 supplying FinA, AutB and LemW together.

\emph{Source.} \href{https://github.com/CocoSgt/riemann-hypothesis-lean/blob/7cbf7a0c32fca1badf05931879a37c3f8a2b03b4/src/Development/AutomorphicArithmetic/AutomorphicVarieties/AutomorphicCases/AutomorphyUnramifiedCoefficients.lean#L24-L27}{AutomorphyUnramifiedCoefficients}.

\emph{Source.} \href{https://github.com/CocoSgt/riemann-hypothesis-lean/blob/7cbf7a0c32fca1badf05931879a37c3f8a2b03b4/src/Development/AutomorphicArithmetic/AutomorphicVarieties/AutomorphicCases/AutomorphyUnramifiedCoefficients.lean#L12-L13}{FinA}.

\emph{Source.} \href{https://github.com/CocoSgt/riemann-hypothesis-lean/blob/7cbf7a0c32fca1badf05931879a37c3f8a2b03b4/src/Development/AutomorphicArithmetic/AutomorphicVarieties/AutomorphicCases/AutomorphyUnramifiedCoefficients.lean#L15-L17}{AutB}.

\emph{Source.} \href{https://github.com/CocoSgt/riemann-hypothesis-lean/blob/7cbf7a0c32fca1badf05931879a37c3f8a2b03b4/src/Development/AutomorphicArithmetic/AutomorphicVarieties/AutomorphicCases/AutomorphyUnramifiedCoefficients.lean#L19-L22}{LemW}.



\paragraph{10. TASK\_\allowbreak CASE3\_\allowbreak AUT\_\allowbreak JOIN — B — Complete the original Case-III automorphy/image join}

Install the actual Case-III ProModAlt, SeedIVAlt and shared PMExAlt with the inherited CohomologyHyp/Ferrero-Washington/HFit/SeedII facts into the original whole Case-III automorphy/image conclusion.

\emph{Source.} \href{https://github.com/CocoSgt/riemann-hypothesis-lean/blob/7cbf7a0c32fca1badf05931879a37c3f8a2b03b4/src/Development/AutomorphicArithmetic/AutomorphicVarieties/ModularityArgument/ThresholdArguments/DenseThresholdCases/PromodularityCaseThree.lean#L43-L44}{AutomorphyRankFourDense}.

\emph{Source.} \href{https://github.com/CocoSgt/riemann-hypothesis-lean/blob/7cbf7a0c32fca1badf05931879a37c3f8a2b03b4/src/Development/AutomorphicArithmetic/AutomorphicVarieties/ModularityArgument/ThresholdArguments/DenseThresholdCases/CuFinal.lean#L44-L44}{main4\_\allowbreak alt}.



\paragraph{11. TASK\_\allowbreak CASE3\_\allowbreak PROMOD\_\allowbreak JOIN — C — Reconstruct the actual Case-III chain-to-promodularity proof}

Use the original tensor-factor tower, ChainRankFour, ChainDesc4, seeded rank-four chain reduction, actual PMExAlt and SeedIVAlt to construct ProModAlt on every original Ctx3 star point at the same fixed iota, preserving all finite-coefficient, inertia and support maps.

\emph{Source.} \href{https://github.com/CocoSgt/riemann-hypothesis-lean/blob/7cbf7a0c32fca1badf05931879a37c3f8a2b03b4/src/Development/AutomorphicArithmetic/AutomorphicVarieties/ModularityArgument/ThresholdArguments/DenseThresholdCases/PromodularityCaseThreeOfLemmas.lean#L131-L153}{UpperTheory.V.Cu.promodularity\_\allowbreak case\_\allowbreak three\_\allowbreak of\_\allowbreak chain\_\allowbreak reduction}.

\emph{Source.} \href{https://github.com/CocoSgt/riemann-hypothesis-lean/blob/7cbf7a0c32fca1badf05931879a37c3f8a2b03b4/src/Development/AutomorphicArithmetic/AutomorphicVarieties/ModularityArgument/ThresholdArguments/DenseThresholdCases/PromodularityCaseThree.lean#L40-L42}{PromodularityCaseThreeDense}.



\paragraph{12. SEED\_\allowbreak IV\_\allowbreak ALT — C — Actual Case-III seed with prechosen residual datum}

Base Ctx3/Four tower, iota and Amb -> w,s before every later Dies tower -> compatible finite-coefficient seed with exact residual support/root entries.

\emph{Source.} \href{https://github.com/CocoSgt/riemann-hypothesis-lean/blob/7cbf7a0c32fca1badf05931879a37c3f8a2b03b4/src/Development/AutomorphicArithmetic/AutomorphicVarieties/ModularityArgument/ThresholdArguments/DenseThresholdCases/PromodularityCaseThree.lean#L19-L25}{SeedIVAlt}.



\paragraph{13. TENSOR\_\allowbreak FACTOR\_\allowbreak TOWER — B — Actual tensor-factor tower and characters}

Original finite-field, tame, factor and character constraints; source-level construction is sizeable and has not been reconstructed here.

\emph{Source.} \href{https://github.com/CocoSgt/riemann-hypothesis-lean/blob/7cbf7a0c32fca1badf05931879a37c3f8a2b03b4/src/Development/AutomorphicArithmetic/AutomorphicVarieties/ModularityArgument/ThresholdArguments/DenseThresholdCases/PromodularityCaseThreeOfLemmas.lean#L131-L153}{UpperTheory.V.Cu.promodularity\_\allowbreak case\_\allowbreak three\_\allowbreak of\_\allowbreak chain\_\allowbreak reduction}.



\paragraph{14. CHAIN\_\allowbreak RANK\_\allowbreak FOUR — C — Actual SomeIn-to-AllIn propagation}

Same source Space and PM ideal, residual point Good4, original bound 30(C+1)\textasciicircum{}2 <= d; geometry must preserve all identities.

\emph{Source.} \href{https://github.com/CocoSgt/riemann-hypothesis-lean/blob/7cbf7a0c32fca1badf05931879a37c3f8a2b03b4/src/Development/AutomorphicArithmetic/AutomorphicVarieties/ModularityArgument/ThresholdArguments/DenseThresholdCases/PromodularityCaseThree.lean#L26-L32}{ChainRankFour}.



\paragraph{15. CHAIN\_\allowbreak DESC4 — B — Original residual chain descent}

Finite residual descent retaining source Chain4 compatibility.

\emph{Source.} \href{https://github.com/CocoSgt/riemann-hypothesis-lean/blob/7cbf7a0c32fca1badf05931879a37c3f8a2b03b4/src/Development/AutomorphicArithmetic/AutomorphicVarieties/ModularityArgument/ThresholdArguments/DenseThresholdCases/PromodularityCaseThreeOfLemmas.lean#L131-L153}{UpperTheory.V.Cu.promodularity\_\allowbreak case\_\allowbreak three\_\allowbreak of\_\allowbreak chain\_\allowbreak reduction}.



\paragraph{16. RANK4\_\allowbreak CHAIN\_\allowbreak REDUCTION — C — Seeded reduction to the original residual chain}

Select T1 before every later Deep tower; retain one compatible seed and Chain4 for every original residual lift over Q.

\emph{Source.} \href{https://github.com/CocoSgt/riemann-hypothesis-lean/blob/7cbf7a0c32fca1badf05931879a37c3f8a2b03b4/src/Development/AutomorphicArithmetic/AutomorphicVarieties/ModularityArgument/ThresholdArguments/DenseThresholdCases/PromodularityCaseThree.lean#L33-L39}{RankFourChainReductionDense}.

\emph{Source.} \href{https://github.com/CocoSgt/riemann-hypothesis-lean/blob/7cbf7a0c32fca1badf05931879a37c3f8a2b03b4/src/Development/AutomorphicArithmetic/AutomorphicVarieties/ModularityArgument/ThresholdArguments/DenseThresholdCases/CuPChainIV.lean#L16-L82}{UpperTheory.V.Cu.chainIV\_\allowbreak alt}.



\paragraph{17. TASK\_\allowbreak COUNT\_\allowbreak V\_\allowbreak INSTALL — B — Complete the original CountV finite assembly}

From exactly the listed source host/cohomology inputs, complete countV\_\allowbreak of\_\allowbreak finiteness\_\allowbreak duality\_\allowbreak euler\_\allowbreak host, including original host/frame formation, normalized length, the CaseV guard and R1<=2\textasciicircum{}c2.

\emph{Source.} \href{https://github.com/CocoSgt/riemann-hypothesis-lean/blob/7cbf7a0c32fca1badf05931879a37c3f8a2b03b4/src/Development/AutomorphicArithmetic/AutomorphicVarieties/HeightAndVanishing/Counts.lean#L69-L69}{CountV}.

\emph{Source.} \href{https://github.com/CocoSgt/riemann-hypothesis-lean/blob/7cbf7a0c32fca1badf05931879a37c3f8a2b03b4/src/Development/AutomorphicArithmetic/AutomorphicVarieties/ModularityArgument/ThresholdArguments/DenseThresholdCases/CountVOf.lean#L33-L60}{UpperTheory.V.UnitaryThreeInertia.countV\_\allowbreak of\_\allowbreak finiteness\_\allowbreak duality\_\allowbreak euler\_\allowbreak host}.



\paragraph{18. THM\_\allowbreak S6\_\allowbreak S — C — Original unitary-three ordinary support and automorphy}

Original host/frame and fixed iota -> one marked g1 -> every universal R,u -> weight-zero finiteness and every Type-I BPrime\_\allowbreak s point.

\emph{Source.} \href{https://github.com/CocoSgt/riemann-hypothesis-lean/blob/7cbf7a0c32fca1badf05931879a37c3f8a2b03b4/src/Development/AutomorphicArithmetic/AutomorphicVarieties/ModularityArgument/ThresholdArguments/DenseThresholdCases/ThmS6CoreS.lean#L17-L26}{ThmS6\_\allowbreak s}.

\emph{Source.} \href{https://github.com/CocoSgt/riemann-hypothesis-lean/blob/7cbf7a0c32fca1badf05931879a37c3f8a2b03b4/src/Development/AutomorphicArithmetic/AutomorphicVarieties/ModularityArgument/ThresholdArguments/DenseThresholdCases/ThmS6CoreS.lean#L12-L16}{BPrime\_\allowbreak s}.



\paragraph{19. HOST\_\allowbreak FRAME\_\allowbreak LENGTH — B — Original Case-V host/frame and normalized-length installation}

Actual HostData, frame, local-condition equality, abstract count and normalized original Selmer length must all refer to the same objects.

\emph{Source.} \href{https://github.com/CocoSgt/riemann-hypothesis-lean/blob/7cbf7a0c32fca1badf05931879a37c3f8a2b03b4/src/Development/AutomorphicArithmetic/AutomorphicVarieties/ModularityArgument/ThresholdArguments/DenseThresholdCases/CountVOf.lean#L33-L60}{UpperTheory.V.UnitaryThreeInertia.countV\_\allowbreak of\_\allowbreak finiteness\_\allowbreak duality\_\allowbreak euler\_\allowbreak host}.



\paragraph{20. HOST\_\allowbreak NUMBER\_\allowbreak S — C — Uniform original-host rank-three catalogue}

One exponent before all closures, coefficient rings, levels, characters, frames, embeddings and marked points; one separating trace function controls all later finite pairwise nonisomorphic collections.

\emph{Source.} \href{https://github.com/CocoSgt/riemann-hypothesis-lean/blob/7cbf7a0c32fca1badf05931879a37c3f8a2b03b4/src/Development/AutomorphicArithmetic/AutomorphicVarieties/ModularityArgument/ThresholdArguments/DenseThresholdCases/HostNumberS.lean#L18-L28}{HostNumber\_\allowbreak s}.



\paragraph{21. TASK\_\allowbreak HOST\_\allowbreak VAN\_\allowbreak E\_\allowbreak INSTALL — B — Transport K/duality to literal HostVanE}

Use the already supplied LemmaK/local-global duality and finiteness on the actual CountV host/frame and exact original coefficient modules.

\emph{Source.} \href{https://github.com/CocoSgt/riemann-hypothesis-lean/blob/7cbf7a0c32fca1badf05931879a37c3f8a2b03b4/src/Development/AutomorphicArithmetic/AutomorphicVarieties/UnitaryThreeDeformations/HostNumber.lean#L9-L12}{HostVanE}.

\emph{Source.} \href{https://github.com/CocoSgt/riemann-hypothesis-lean/blob/7cbf7a0c32fca1badf05931879a37c3f8a2b03b4/src/Development/AutomorphicArithmetic/AutomorphicVarieties/LevelCohomology/SoHostVanE.lean#L11-L18}{UpperTheory.V.X.So.hostVanE}.



\paragraph{22. TASK\_\allowbreak HOST\_\allowbreak XPP\_\allowbreak INSTALL — B — Transport CorX to literal HostXpp}

Preserve Psi\textasciicircum{}-2 eta, the prime-to-ell quotient, actual local conditions and the original CountV host while installing BR7.

\emph{Source.} \href{https://github.com/CocoSgt/riemann-hypothesis-lean/blob/7cbf7a0c32fca1badf05931879a37c3f8a2b03b4/src/Development/AutomorphicArithmetic/AutomorphicVarieties/UnitaryThreeDeformations/HostNumber.lean#L14-L17}{HostXpp}.

\emph{Source.} \href{https://github.com/CocoSgt/riemann-hypothesis-lean/blob/7cbf7a0c32fca1badf05931879a37c3f8a2b03b4/src/Development/AutomorphicArithmetic/AutomorphicVarieties/LevelCohomology/Xpp.lean#L11-L84}{UpperTheory.V.X.xpp}.



\paragraph{23. DESIGN\_\allowbreak CHOICES\_\allowbreak COMMON — C — Construct the original common design family and bundle}

For the one target HostExW, construct the original compatible DesignChoices on Cdes=WithTw(WithSm Cprime), where Cprime=WithChN(A intersect target class), A is the one Built construction class, and cT(eps,y)=cTmin(y). Retain every original field, tuple, word, metric and common bundle duty. DSZ10 and RMP11 give size/rank properties of class members, not member existence. The original RecExP/RecAtG and common object/row route remains coupled. Construct ONE original common design family R and its recursive/geometric half with original Inputs/Pins/Holds/ExAt/InGeo and required CoreP. PKT now reconstructs full exists\_\allowbreak pkT, rowPk\_\allowbreak nonempty and exists\_\allowbreak moved with Copy/Coh on those original given inputs and a supplied fitting later stage: one catalogue, both ChT, complete LfT, actual integral lattices and all metric/measure/period constants are retained; one numerical record precedes every eta and V/S/N/G. RDR supplies one Side0 before all readers, ReadPk/pins and full ReadL with exact o.eta=q.eta. The initial actual geometry, Fits/division/central conditions, primitive rational tests and eta agreement must still be produced. Inherited CSR supplies finite columns, one ReadC before later twists, period/evaluation and scalar/error transports, RowSideDetermined and same-R Bundle/BTw/BWit/BunU consumers. RPI supplies exact RowPairInt for actual pairs and finite translates with original mu\_\allowbreak b/fd0; first literal DivDef3 at the same updated R,F,Y,Z,C,sEE follows from its retained original readings, Pins/Holds, stored nonzero IsItW chain, five real-group predicates and group tie. Remaining common fields include host/frame, MultEq nonzero-part support, that original nonzero chain/group tie, eigensystem/carrier/parts and all-embedding logLam, HD/d3b, NEb, actual rational tests and base rows. Retain every simultaneous family hypothesis and eventual threshold. Stronger second DivRows3 remains separate; universal quotient-height CuspRapidDecay and actual common member existence are not inferred from the comparison constructors. Version 0.41 supplies the finite phase-correct rational theta route: exact coefficient extension, split representatives, square-law/PhiIs, Pins from CoversAt, class-sum/point-identity transfer to OneMapOverL, actual-good-projector RatG and both-level LiftTestB. Its original higher covers, all-tuple class-sum uniqueness including phase rescaling, one PointIdentCore before O, base ClassSumsSpan and geometric LatFull remain explicit producers. MLat is available from the original LatFull input. This sharpens the rational-test work without claiming actual common geometry, Fits or eta agreement. Version0.46 proves finite generation of the original whole integral cusp lattice by actual integral-frame evaluation. It also proves the exact identity-cover, analytic unfolding, original-field coefficient descent, global clearing-scalar and both-level lattice-transport implications. Full spanning retains the supplied original IsNab connection and the five explicit principal-chart, level-system, raw-Gaussian, short-rule/dictionary and ramified-jet producers. The higher tuple covers, actual common class sums, PointIdentCore and ClassSumsSpan remain separate.

\emph{New completed inputs.} ICL1-ICL21

\emph{Current remaining subfrontier.} Actual common geometry and higher covers/class sums/PointIdentCore/ClassSumsSpan, plus the supplied original IsNab and five exact producers retained for cusp-lattice spanning; existing host/chain/coefficient/Fits/eta obligations remain coupled.



\paragraph{24. LOCAL\_\allowbreak NBUILT\_\allowbreak COMMON — C — Actual local specialization on the common selected family}

Construct, with the required closed-class property on the same Des(C,3), one eventual family of actual NBuilt certificates for the original Fm.D(x,i), retaining the certificate before forgetting to TSel. HLP now reconstructs full HeckeLift0, all-index commutation, HeckeStd0 and HeckeIntTw from the original MtInput and complex Hecke facts, in the full ambient GenAt range k<=nV. Every maximal ideal is an actual supported residual character class; literal MtSub is the O-span of projections of the original MtLat generators, is finite free, and localizes onto the exact NCov.part. Literal Psi has its exact range and off-space projection error, both original span assertions hold, and ActEx/ActHK are constructed on the same lattice under original semisimplicity. Original MjLat relation quotients, universal coefficient-map relation fields, selected tame generators, coefficient embeddings and all identity transports give the complete nonempty\_\allowbreak spec constructor on its full printed input list. A single compositum combines the four original coefficient requirements with big fixed beforehand. Whole ndSpec\_\allowbreak exists still needs actual common geometric/choice MtInput, complex HeckeStdPlain and the no-Jordan producer, original F3 bridges with valuation/GoodCoef/iq, original F4 common-overfield reading of the entire coefficient field, and exact arithmetic character/reciprocity/admissibility suppliers. Original AutSpec, Galois specialization and the remaining NBuilt fields on that same installation also remain. Do not replace the actual projected lattice or maps by abstract objects with the same dimensions.



\paragraph{25. SECOND\_\allowbreak LAYER\_\allowbreak COMPATIBLE\_\allowbreak SELECTION — C — Original second-layer selector and class duties}

For the same A, supply Duty on A and a Closed common original TSel family whose selected records are the original true-column rec7 records (or have the entire dependent tuple transported), with their real witness and original coefficient-field root data. Existence/compatibility remains arithmetic; merely abstractly isomorphic fields or independently selected columns do not suffice. CSR21–22 now construct the same-R Bundle/BTw/BWit/BunU outputs from the exact selected realization/row-side packet and NEb, with the original EvP/ParEv quantifier order. Actual common selection, the five stated realization packets and NEb remain upstream; none is replaced by an independently chosen compatible record. Version 0.38 closes the selected analytic extension and reconstructs the first updated DivDef3 consumer; it leaves the original common selection, geometric readings, rational tests, base rows, host/parts/NEb and exact retained DI premises upstream, with unchanged EvP/ParEv thresholds. Version 0.39 supplies the original given-data reader constructors and exact RowPairInt/first DivDef3 on original premises. It leaves one actual common geometric/comparison/test member, eta agreement, compatible selection, rational tests, base rows, host/parts/NEb and original EvP/ParEv thresholds upstream. Version 0.40 also supplies original comparison/RowPk/fitting-stage Moved algorithms and exact common data transport. The actual initial geometric member, fitting tests, eta agreement and compatible second-layer selector with original eventual thresholds remain upstream. Version 0.41 supplies the finite phase-correct rational theta route: exact coefficient extension, split representatives, square-law/PhiIs, Pins from CoversAt, class-sum/point-identity transfer to OneMapOverL, actual-good-projector RatG and both-level LiftTestB. Its original higher covers, all-tuple class-sum uniqueness including phase rescaling, one PointIdentCore before O, base ClassSumsSpan and geometric LatFull remain explicit producers. MLat is available from the original LatFull input. This sharpens the rational-test work without claiming actual common geometry, Fits or eta agreement.



\paragraph{26. BASE\_\allowbreak THETA\_\allowbreak READINGS — C — Actual base reading on the common bundle}

For every family in the original class, eventually on the SAME selected Bundle and for each i<m, supply the actual R,F,Read,ReadL, weak models/Weil data and phase-one pins. Identify their LiftTie.theta, mtRead, mjRead and ThDef3 basis with that Bundle.P.A3, P.Zs, P.S0 and its stored maps. Exact field, embeddings, metrics, measures, row family and weak choices are retained. This remains one part of the coupled BTH12 producer. PKT now supplies the complete original comparison packet, RowPk and fitting-stage Moved with Copy/Coh on original given geometry, Inputs/Pins/Holds/ExAt/InGeo and Fits. Its ChT/LfT/Reps, actual lattice carriers, original Haar and norm constants, period maps, representatives and all-place RrT are constructed together; its one numerical record precedes every eta/V/S/N/G. RDR gives the original level-zero side, moved-Read and ReadPk/pin constructors on those objects, with one Side0 before every reader and all twelve Read ties. Full ReadL retains the exact o.eta(j,v)=q.eta(j,v) guard for every 1<=j<=m and place v; arbitrary Over does not imply it. Inherited CSR fixes one ReadC before all TwArgs, transports both metrics and all thirteen reading maps, and gives actual period/evaluation with kappa=(c*nP)\textasciicircum{}(-1/2). RPI supplies exact row-pair integrability, so first updated DivDef3 follows from the retained original Pins/Holds, nonzero chain, five real-group predicates and group tie. Initial theta extension remains complete. Remaining work is actual initial common geometry and its original predicates, Fits/primitive rational tests, eta agreement and outer family/eventual/same-Bundle duties, reFace compatibility, MultEq support, original chain/group data, actual base rows and any stronger DivRows3 theta square. No new compatible reading, family threshold or Haar normalization is substituted. Version 0.41 supplies the finite phase-correct rational theta route: exact coefficient extension, split representatives, square-law/PhiIs, Pins from CoversAt, class-sum/point-identity transfer to OneMapOverL, actual-good-projector RatG and both-level LiftTestB. Its original higher covers, all-tuple class-sum uniqueness including phase rescaling, one PointIdentCore before O, base ClassSumsSpan and geometric LatFull remain explicit producers. MLat is available from the original LatFull input. This sharpens the rational-test work without claiming actual common geometry, Fits or eta agreement.



\paragraph{27. BASE\_\allowbreak THETA\_\allowbreak RATIONAL\_\allowbreak TESTS — C — Simultaneous original base rational tests}

On the SAME original Bundle/reading/theta map, for every i<m supply one finite separating complex-linear test family whose values lie in the fixed embedded coefficient field on theta of EVERY original rational source vector and on EVERY stored target-coordinate basis vector. BTH8/BTH11 use this same LiftTestB for rational adjoint image and matrix coordinates. RTG now reconstructs prescribed-field coefficient extension, original split class representatives/RepsOK, the self-dual split-lattice square law and finite residue PhiIs; every Pins field follows from the actual CoversAt. With the base coefficient record literally fixed, exact class-sum existence/uniqueness including phase-rescaled O, PointIdentCore with common lambda,z before every O, and base ClassSumsSpan imply all-tuple finite lift span and OneMapOverL. The phase is exactly cj(twAt embLL A3.A.v0)(R.phs i). RatG uses the actual good projector eOp(y), and full LiftTestB follows at BOTH original levels. These finite implication proofs are complete. Remaining original producers are the prescribed higher covers, those actual class sums/uniqueness, PointIdentCore, base ClassSumsSpan and the base/common geometric datum; exact MLat is supplied only from original LatFull, whose geometric production remains separate. Finite dimension or an unrelated rational theta model cannot replace these same-object identities. Original outer eventual-family/Bundle quantifiers remain. Version0.46 proves finite generation of the original whole integral cusp lattice by actual integral-frame evaluation. It also proves the exact identity-cover, analytic unfolding, original-field coefficient descent, global clearing-scalar and both-level lattice-transport implications. Full spanning retains the supplied original IsNab connection and the five explicit principal-chart, level-system, raw-Gaussian, short-rule/dictionary and ramified-jet producers. The higher tuple covers, actual common class sums, PointIdentCore and ClassSumsSpan remain separate.

\emph{New completed inputs.} ICL1-ICL21

\emph{Current remaining subfrontier.} Actual higher covers, all-tuple class sums/uniqueness, common PointIdentCore, base ClassSumsSpan and common geometry; integral-lattice finite generation is proved, while spanning still needs the original IsNab and five named producer groups. Original both-level FWTie transport is available on those inputs.



\paragraph{28. SLR\_\allowbreak BASE\_\allowbreak ROWS\_\allowbreak PRODUCER — C — Actual base row identities on that same bundle}

For the SAME selected Bundle, common original fields and readings used by the other three BTH12 inputs, supply actual RowsHold0\_\allowbreak 3 uniformly eventually for every family. Its original primitive row realization, spectrum/type and basis identifications must be retained. BTH supplies a theta consumer only; it does not construct the base row identity. No separate witness may be chosen merely to satisfy this row requirement.



\clearpage
\section{The complete rank-three congruence integer}
\label{sec:rank-three-congruence-v46}

This chapter reconstructs the original $\mathrm{CNR3}$ from
\texttt{Congruence\allowbreak Number\allowbreak Rank\allowbreak Three.\allowbreak lean}, including its two uniform
trace catalogues.  Its elementary product argument will choose the
integer before the coefficient prime.  The catalogue proof uses only
ranks three and two; it does not assume the original all-rank
$\mathrm{CountK}$ or $\mathrm{CountQ}$ producers.

Fix an imaginary quadratic field $K$.  Put $d=|\operatorname{disc}K|$.
Write $\varepsilon$ for the cyclotomic character, with
$\varepsilon(\sigma)=q$ at an arithmetic Frobenius $\sigma$ at a
place of residue cardinality $q$ away from its coefficient prime.
Consequently geometric Frobenius is $\sigma^{-1}$.
For a spherical automorphic representation of rank $n$, our
algebraic Satake multiset is
\[
 A_v(\pi)=q_v^{(n-1)/2}\operatorname{Sat}_v(\pi).
                                                               \tag{RNC.1}
\]
It is the geometric Frobenius multiset.  The arithmetic multiset
$\beta_v(\pi)$ is its elementwise inverse.

The original $\mathrm{Level}(\ell,M,m,r)$ is retained literally.
It includes continuity; away from $\ell$, inertia is diagonal with
characters killing the subgroup which fixes every $M$-th root of
unity.  At each place above $\ell$ there is a complete ordered
Galois flag with diagonal characters $\kappa_i$ and
\[
 \kappa_i|_I=\varepsilon^{m_i}\chi_i|_I,
 \qquad \chi_i\text{ kills the subgroup fixing }\mu_M,
                                                               \tag{RNC.2}
\]
together with the disjunction
\[
 (\forall i,j,\ \kappa_i\ne\kappa_j\varepsilon)
       \quad\text{or}\quad r\text{ is potentially crystalline}.
                                                               \tag{RNC.3}
\]
The potentially crystalline clause has its original local bridge
and period-ring meaning.  No restriction such as $\ell\mid M$ is
added.  All fields, algebraic closures, scalar towers and algebraic
isomorphisms $\iota:\overline{\mathbb Q}_\ell\simeq\mathbb C$
below are the original ones.  In particular $\ell=3$ is allowed.

We use the named classical foundations already recorded with CX:
Clozel's cohomology and rationality theorem, Borel--Serre finite
arithmetic CW models, Borel--Labesse--Schwermer's injection of
cuspidal cohomology, strong multiplicity one and local genericity,
local Langlands conductor and spherical compatibility, the generic
newvector theorem, and finite-coefficient $p$-adic Hodge comparison
and weak admissibility.  Their numerical and normalization
applications are proved below, rather than supplied as a counting
theorem.  The extra rank-two adelization step uses the cuspidal
decomposition in Jacquet--Langlands, Proposition 10.9, and the
explicit classical-to-adelic construction in RNC2.  The exact
published statements and primary URLs are recorded separately in
the classical input registry.

% BEGIN RNC1
\paragraph{RNC1. The precise automorphic Galois input in the two required ranks.}
\label{claim:RNC1}
For an initial regular algebraic cuspidal polarized pair over
a CM or totally real field, with arbitrary initial integral
highest weight $\lambda$, the published theorem supplies a
continuous semisimple Galois representation at every prime,
the away-from-that-prime comparison (RNC.5), and Hodge--Tate
weights $\{\lambda_i+n-i\}$.  This general existence statement
is available before a flag has identified $\lambda$.
After that identification, its following bounded
specialization will be used.
For the following two types of cuspidal representations, the
automorphic Galois representation exists for every coefficient
prime and every coefficient isomorphism, is continuous and
semisimple, and has the stated Hodge--Tate multiset:
\[
\begin{array}{c|c|c|c}
 F,n& m&\lambda_i=-m_{n+1-i}-(n-i)&\mathrm{HT}\text{ multiset}\\ \hline
 K,3&(e-1,-1,-e-1)&(e-1,0,1-e)&\{e+1,1,1-e\}\\
 \mathbb Q,2&(e,-e)&(e-1,-e)&\{e,-e\}.
\end{array}                                                     \tag{RNC.4}
\]
Here $e\geq1$, the first representation is conjugate self-dual,
and the second is essentially self-dual.  Their infinitesimal
character is $\operatorname{infOf}(m)=\{(n-1)/2+m_i\}$ at every
embedding.  Local compatibility, including monodromy, is
\[
 \iota\,\mathrm{WD}(R|_{G_{F_v}})^{\mathrm{F\!ss}}
      \simeq\operatorname{rec}_{F_v}
                    (\pi_v|\det|^{(1-n)/2}) .                 \tag{RNC.5}
\]
At $v\nmid\ell$ the parameter is pure, of geometric weight $2$
in the first row and $0$ in the second.  A spherical place above
the coefficient prime gives a crystalline representation.
\begin{proof}
This is a checked application of the published theorems, not a
new automorphy assertion.  Barnet-Lamb--Gee--Geraghty--Taylor,
\emph{Potential automorphy and change of weight}, Theorem 2.1.1,
applies to regular algebraic cuspidal polarized pairs over CM
\emph{or totally real} fields.  Its definition in Section 2.1
asks for $\pi^c\simeq\pi^\vee\otimes(\chi\circ N_{F/F^+})$,
with the prescribed constant infinity sign.  For the first row
take the quadratic character of $K/\mathbb Q$ as multiplier:
its pullback by the norm is trivial and its real sign is $-1$,
as required in odd rank.  For rank two over $\mathbb Q$ take
the actual central character whenever the cusp satisfies
the stated essential self-duality.  RNC2 will verify this
input for the newform application; it is not an antecedent
for the general existence theorem here.  There is only one real place, so the
constant-sign condition imposes no extra choice there.

The weight convention of that theorem is that $\pi_\infty$
has the infinitesimal character of the dual of the highest-weight
module of weight $\lambda$.  The formula in (RNC.4) gives exactly
$\lambda_i+n-i=-m_{n+1-i}$, so it matches the original
$\operatorname{infOf}$ convention.  Both displayed $\lambda$
are dominant integral for $e\geq1$.  The theorem's Hodge formula
and conjugate symmetry give respectively geometric weights $2$
and $0$.  Its clauses (2)--(4) give purity and (RNC.5) away from
the coefficient prime, and crystallinity at spherical places
above it.  These statements require neither irreducibility of
every associated representation nor an automorphy-lifting lower
bound on the prime.

For (RNC.5) at the coefficient prime use the same authors'
\emph{Local-global compatibility for $l=p$, II}, Theorem A,
Theorem 1.1 and Corollary 1.2.  They give full Frobenius
semisimplification for Shin-regular weight.  Their definition is:
odd rank suffices; in even rank a strict inequality
$\lambda_i>\lambda_{i+1}$ for some odd $i$ at some embedding
suffices.  Rank three meets the first alternative even at $e=1$.
For rank two, $\lambda_1-\lambda_2=2e-1>0$, so the second
alternative holds.  No Iwahori hypothesis remains in this
theorem, and every prime is allowed.  Its reciprocity sends a
uniformizer to geometric Frobenius, giving precisely (RNC.1)
and (RNC.5).  We use this stronger conclusion only in these
two verified ranks and weights.
\end{proof}
% END RNC1

% BEGIN RNC2
\paragraph{RNC2. The actual newform supplies the required rank-two cusp.}
\label{claim:RNC2}
Let $e\geq1$, let $f$ be the original newform of weight
$k=2e+1$, and suppose $f$ is weakly attached to
$W\varepsilon^{-e}$ in the original sense.  Then there is a
regular algebraic essentially self-dual cusp $\pi$ on
$\mathrm{GL}_2/\mathbb Q$, of the second infinity type in
(RNC.4), weakly attached to the same $W$.  At every prime
$p\nmid f.M$, if $\alpha_p,\alpha'_p$ are the roots of
$f.\mathrm{heckePoly}(p)$, its algebraic Satake multiset is
\[
           A_p(\pi)=\{p^{-e}\alpha_p,p^{-e}\alpha'_p\}.
                                                               \tag{RNC.6}
\]
\begin{proof}
Here is the bounded adelization and its normalization.  The
original definition supplies the $\Gamma_0(f.M)$ nebentype law,
cuspidality, $a_1(f)=1$, and the coefficient identity
\[
 a_{pn}(f)+\chi(p)p^{k-1}{\bf1}_{p\mid n}a_{n/p}(f)
                         =a_p(f)a_n(f)\quad(p\nmid f.M).
                                                               \tag{RNC.7}
\]
On $g_\infty\in\mathrm{GL}_2^+(\mathbb R)$ put
\[
 F_\infty(g_\infty)=\det(g_\infty)^{k/2}
               j(g_\infty,i)^{-k}f(g_\infty i),
 \qquad j\!\left(\begin{pmatrix}a&b\\c&d\end{pmatrix},z\right)
                     =cz+d.
                                                               \tag{RNC.8}
\]
Strong approximation writes every adelic matrix as
$\gamma g_\infty u$, with $\gamma\in\mathrm{GL}_2(\mathbb Q)$
and $u\in K_0(f.M)$.  Define its value to be
$\chi(d_u)^{-1}F_\infty(g_\infty)$.  Two such decompositions
differ by an integral determinant-one matrix in
$\Gamma_0(f.M)$; its nebentype factor in (RNC.8) cancels the
change in $\chi(d_u)^{-1}$.  Thus this is a well-defined actual
left-rational-invariant adelic cusp form.  Its compact-open
finite transformation law is explicit.  Its unipotent constant
terms are the classical cusp constant terms and vanish.
The factor in (RNC.8) is invariant under positive real scalar
matrices and has rotation weight $k$; differentiating the
holomorphic identity gives the usual lowest discrete-series
vector, with infinitesimal character
$\{(k-1)/2,-(k-1)/2\}$.

Decompose this nonzero form in the cuspidal representation space
with its central character, by Jacquet--Langlands,
Proposition 10.9.  The fixed finite level and fixed infinity
type make the relevant space finite-dimensional, so at least
one irreducible cuspidal summand has a nonzero projection of
the form.  At $p\nmid f.M$, the $p+1$ standard double-coset
representatives act on (RNC.8) by the two terms of (RNC.7).
The normalized spherical parameters of every such nonzero
projection therefore have sum $p^{-(k-1)/2}a_p(f)$ and product
$\chi(p)$.  This also follows directly by comparing every
Fourier coefficient in (RNC.7); a holomorphic cusp form with
all coefficients zero vanishes.  These identities are preserved
by the projection, since the spherical Hecke algebra commutes
with the global representation decomposition.

Let $\pi^0$ be one of these summands and set
$\pi=\pi^0|\det|^{1/2}$.  At the finite uniformizer this twist
multiplies each spherical root by $p^{-1/2}$, while at infinity
it adds $1/2$ to both infinitesimal parameters.  As $k=2e+1$,
the infinity multiset becomes $\{e+1/2,1/2-e\}$ and (RNC.1)
gives (RNC.6).  This is the exact original
$\mathrm{NewformRep}$ specialization $j=e$, not a change to
the original form or Galois representation.

For any rank-two unramified parameter with product equal to
the central-character value, inversion is the same unordered
pair as multiplication by the inverse central character.
Strong multiplicity one consequently gives
$\pi^\vee\simeq\pi\otimes\omega_\pi^{-1}$.
The central character is algebraic at infinity for this
displayed integral highest weight, so $\pi$ is an admissible
polarized pair for RNC1.  Finally, at geometric Frobenius
$\varepsilon^{-e}$ has value $p^e$.  The original weak
attachment of $f$ therefore says that the eigenvalues of
$W(\sigma^{-1})$ are exactly (RNC.6), at the same cofinite
set of primes and with all original prolongations and lifts.
This proves the required weak attachment to $W$.
\end{proof}
% END RNC2

% BEGIN RNC3
\paragraph{RNC3. Weak attachment identifies the actual representation and its weights.}
\label{claim:RNC3}
Suppose an original continuous irreducible representation $r$
is weakly attached to an initial regular algebraic
conjugate-self-dual cusp in rank three, or to the rank-two
cusp of RNC2, and has the corresponding original level
and ordered weight $m$ from (RNC.4).  Then it
is isomorphic, in its original algebraic closure, to the
representation supplied by RNC1.  Its original complete flag
recovers exactly the highest weight in (RNC.4).  The statement
applies in rank three to $r=\rho\varepsilon^{-1}$ from
$\mathrm{Adm3}(\iota,e,N,\rho)$, and in rank two to the
actual $W$ from $\mathrm{Adm2}$.
\begin{proof}
First choose the representation attached by the published
theorem to the given regular algebraic cusp; at this point its
highest weight need not have been identified: this uses
the general-existence part of RNC1, not its later Shin-regular
specialization.  Its cofinite
Frobenius characteristic polynomials agree with those of $r$
by the definition of weak attachment and (RNC.1).
Both representations are continuous and semisimple.  ADK1
captures their compact images in a common finite coefficient
extension.  Chebotarev and continuity make their traces equal
on the whole original Galois group.  The semisimple character
argument of CAP6 then gives an actual intertwining matrix over
$\overline{\mathbb Q}_\ell$.  This argument neither changes
$\iota$ nor replaces the chosen algebraic closure.

The representation is therefore de Rham.  Its flag subobjects
and successive quotients are de Rham by exactness of the
de Rham category.  Each quotient in (RNC.2) is an unramified
character times a finite character times
$\varepsilon^{m_i}$.  Its Hodge--Tate weight at every relevant
embedding is $-m_i$.  Exactness gives the whole multiset
$\{-m_i\}$; it is not inferred merely from being an extension
of de Rham characters.  The automorphic Hodge formula now
identifies the decreasing integers $\lambda_i+n-i$ with
$-m_{n+1-i}$, giving (RNC.4) and the asserted infinity type.
Every local comparison is transported by this same matrix.

For rank three, twisting by $\varepsilon^{-1}$ preserves
irreducibility and continuity.  Away from $\ell$ it changes
no inertia character.  At $\ell$ it replaces $m=(e,0,-e)$
by $(e-1,-1,-e-1)$ and every $\kappa_i$ by
$\kappa_i\varepsilon^{-1}$.  The first alternative in
(RNC.3) is unchanged; the second is preserved because the
cyclotomic character is crystalline.  Thus the same original
integer $N$ is a level for $r$, and its automorphy is exactly
the $\mathrm{Adm3.aut}$ premise.  For rank two, RNC2 gives
the required weak cusp for the original $W$; its original
level and irreducibility are fields of $\mathrm{Adm2}$.
\end{proof}
% END RNC3

% BEGIN RNC4
\paragraph{RNC4. The flag detects every forbidden monodromy ratio.}
\label{claim:RNC4}
At a place $v\mid\ell$ of residue cardinality $q$, let $r$
be as in RNC3.  The semisimple Weil characters of its parameter
are the rank-one parameters $\mu_i$ of the ordered flag.
With arithmetic Weil degree $\deg(s)$, they satisfy
\[
 \mu_i(s)=q^{m_i\deg(s)}\iota(\psi_i(s)),\qquad
 \psi_i=\kappa_i\varepsilon^{-m_i}.                            \tag{RNC.9}
\]
Their inertia restrictions are $\iota\chi_i$.  If
$\mu_i(s)=q^{\deg(s)}\mu_j(s)$ for all Weil elements $s$,
then $\kappa_i=\kappa_j\varepsilon$ on the whole local Galois
group.  Consequently the first alternative of (RNC.3) forces
zero monodromy, both for the Galois and the automorphic
parameters.
\begin{proof}
Use a common finite coefficient extension as in ADK1.
RNC3 makes $r$ and all its flag terms de Rham, hence potentially
semistable.  Exactness of potentially semistable comparison
takes the flag to a flag of filtered modules and of raw Weil
actions; on its successive rank-one quotients the cyclotomic
character has arithmetic Weil value $q^{\deg(s)}$, a finite
character retains its inertia value, and an unramified
character retains its Frobenius value.  This gives (RNC.9),
including every original local bridge after scalar extension.
Frobenius semisimplification preserves the characteristic
polynomial on every Weil element, so the semisimple Weil
representation is the direct sum of these characters.

The continuous character $\psi_i$ has finite inertia image:
its restriction there is $\chi_i$.  More generally a continuous
character of the compact local Galois group has norm one,
because its norm image is a compact subgroup of
$\mathbb R_{>0}$ and all integer powers of any element are
bounded.  A proposed ratio in the assertion would give
\[
 \iota(\psi_i/\psi_j)(s)
                    =q^{(1+m_j-m_i)\deg(s)}.
                                                               \tag{RNC.10}
\]
Take an arithmetic Frobenius lift and apply $\iota^{-1}$,
which fixes the rational integer $q$.  The left side has
$\ell$-adic norm one, whereas $q$ has norm strictly less than
one.  Thus $m_i=m_j+1$.  Equation (RNC.10) then says that
$\psi_i/\psi_j$ is one on the Weil group.  Density and
continuity make it one on the whole decomposition group,
giving the claimed equality of $\kappa$'s.

In a Frobenius-semisimple Weil--Deligne module, a nonzero
monodromy entry maps some character space into its norm twist,
by $r(s)Nr(s)^{-1}=q^{\pm\deg(s)}N$; interchanging the two
indices covers the sign convention.  All such ratios were
excluded.  Therefore $N=0$.  Equality of the semisimple Weil
representations in (RNC.5) gives exactly the same exclusion
for the automorphic parameter.  The argument does not assert
that arbitrary extensions of characters are crystalline.
\end{proof}
% END RNC4

% BEGIN RNC5
\paragraph{RNC5. A polynomial automorphic level independent of the coefficient prime.}
\label{claim:RNC5}
The same cusps associated to the original representations have
nonzero principal-level fixed vectors at $M^{2n-1}$: level
$N^5$ for $\mathrm{Adm3}$ and level $M^3$ for the $W$ in
$\mathrm{Adm2}(\iota,e,M,\rho)$.  They are spherical at
every place not dividing $M$.  The original Galois level
remains $M$ throughout.
\begin{proof}
For a finite place $v$ write $a_v=v(M)$, normalized so that
$M\mathcal O_{F_v}=\mathfrak p_v^{a_v}$.  A character which
kills the subgroup fixing $\mu_M$ has local conductor at most
$a_v$.  Indeed local reciprocity identifies its restriction
with a character through the cyclotomic extension.  For
$x\in1+M\mathcal O_{F_v}$, the norm to the rational local
field is congruent to $1$ modulo $M$: multiplication by $x$
on an integral local basis has matrix $1+M B$, whose
determinant is $1$ modulo $M$.  The cyclotomic action is
therefore trivial.  If $a_v=0$, inertia fixes the prime-to-
residue-characteristic roots in question, so the conductor
is zero.  This proof also covers ramified places of $F$.

At $v\nmid\ell$ the original inertial action is a direct sum
of these finite characters.  Its monodromy is zero and its
Artin conductor is at most $n a_v$.  At $v\mid\ell$, RNC4
identifies the semisimple inertial characters of the parameter
with the same finite characters, so their conductor sum is
again at most $n a_v$.  For an arbitrary nilpotent monodromy
operator the full conductor is
\[
 a(r_v,N_v)=a(r_v)+\dim V^{I_v}-\dim(\ker N_v)^{I_v}
                              \leq n a_v+n-1.                 \tag{RNC.11}
\]
If $a_v\geq1$ this is at most $(2n-1)a_v$.

It remains to exclude an unwanted coefficient-prime factor
when $a_v=0$.  Under the first alternative of (RNC.3), RNC4
gives $N_v=0$.  Under the second, the Galois parameter has
zero monodromy by potential crystallinity, and the full
compatibility in RNC1 transfers this to the automorphic
parameter.  Inertia is already trivial when $a_v=0$.
The parameter is unramified, hence the local cusp component
is spherical.  This proves conductor at most $(2n-1)a_v$
at every place, with no condition $\ell\mid M$.

Every local component of a cuspidal $\mathrm{GL}_n$
representation is generic.  Local Langlands preserves its
conductor, and the generic newvector theorem supplies a
nonzero vector fixed by $K_1(\mathfrak p_v^{b_v})$ at its
conductor $b_v$.  The principal subgroup of depth
$(2n-1)a_v$ is contained in this subgroup.  Tensor these
vectors over the finitely many nonspherical places, using
the spherical vectors elsewhere.  The restricted tensor
product realization of the actual cusp gives a nonzero
global vector at the asserted principal level.  Thus this
is a bound on the auxiliary automorphic label set, not a
replacement of the original Galois level condition.
\end{proof}
% END RNC5

% BEGIN RNC6
\paragraph{RNC6. Fixed cells and explicit coefficient ranks give uniform finite catalogues.}
\label{claim:RNC6}
For fixed $K$ there are natural constants $c_3,c_2$ such that,
for every $e\geq1$ and $L\geq1$, the numbers of cusp
isomorphism classes with principal-level $L$ vectors and the
two infinity types of (RNC.4) are at most
\[
 \begin{split}
 c_3 L^{18}(e+1)^{18}&\quad(F=K,n=3),\\
 c_2 L^4(e+1)^4&\quad(F=\mathbb Q,n=2).
 \end{split}                                                   \tag{RNC.12}
\]
One may restrict the first set to conjugate-self-dual classes.
\begin{proof}
We give the bounded-rank version of the fixed-cell argument
underlying CX1.  In the first row, the highest-weight module
of weight $(e-1,0,1-e)$ is a determinant twist of the module
of weight $(2e-2,e-1,0)$.  HP1--HP2 realize the latter in
polynomials of degree $3e-3$ in the nine matrix coordinates.
Its dimension is at most $(3e-2)^9\leq[3(e+1)]^9$; this
inequality follows even by bounding each individual monomial
exponent between zero and $3e-3$.  Tensor over the two
embeddings of $K$ to obtain coefficient rank at most
$[3(e+1)]^{18}$.  The determinant twist changes no rank.
For the second row, shift $(e-1,-e)$ by $e$ to obtain
$(2e-1,0)$.  The same polynomial realization in four matrix
coordinates bounds the coefficient rank by
$(2e)^4\leq[2(e+1)]^4$.  These formulas include $e=1$.

Use principal level $3$ once and for all.  Its arithmetic
components may stabilize projective $\mathcal O_F$-lattices
rather than the standard free one.  Each is nevertheless
a free integral lattice of rank $dn$ after restriction of
scalars to $\mathbb Z$, and its principal level-$3$ action
is congruent to one on that same lattice.  Thus its arithmetic
groups are torsion-free: restriction of scalars embeds them
in the kernel of $\mathrm{GL}_{dn}(\mathbb Z)\to
\mathrm{GL}_{dn}(\mathbb Z/3)$, with $d=[F:\mathbb Q]$.
For completeness, if a finite-order matrix in this kernel
is not one, take a prime-order power of it and write it as
$1+3^a B$, $a\geq1$, with an entry of $B$ prime to $3$.
Expanding its prime-th power, for a prime other than $3$
the linear term has strictly smaller $3$-adic valuation
than all later terms; for prime $3$ the same holds because
$a\geq1$.  It cannot equal one.  This is a contradiction.

The level-$3$ adelic quotient has finitely many components,
using the finite ideal class group.  Remove the rational
split central direction.  Borel--Serre gives a finite CW
model for each of these fixed torsion-free components; let
$b_F$ be their total number of cells, in all degrees.
The cover at level $3L$ has total covering degree at most
\[
 [K_f(3):K_f(3L)]\leq
 |M_n(\mathcal O_F/(3L))|=(3L)^{dn^2}.                         \tag{RNC.13}
\]
The sum of the ranks of cellular cochains with coefficient
rank $D$ is therefore at most $b_F(3L)^{dn^2}D$.
Their total cohomology dimension is no larger.

Every cusp under consideration is cohomological for the
explicit coefficient just constructed, by Clozel's theorem
with the matching infinitesimal character.  A nonzero
level-$L$ vector is also a level-$3L$ vector.  The cuspidal
cohomology injection of Borel--Labesse--Schwermer places its
nonzero cohomological summand in ordinary cohomology at
that level.  Distinct global cusp classes give independent
summands; multiplicity one, or simply decomposition into
inequivalent irreducible cuspidal summands, ensures that
each class costs at least one dimension.  Applying (RNC.13)
with $dn^2=18$ and $4$, and the coefficient-rank bounds,
proves (RNC.12) after enlarging two fixed integer constants.
Neither constant depends on the coefficient prime, $e$, $L$,
or the place at which a Frobenius value will later be read.
\end{proof}
% END RNC6

% BEGIN RNC7
\paragraph{RNC7. Full cusp labels and algebraic Satake conjugation.}
\label{claim:RNC7}
Let $S_3(e,N)$ be the set of conjugate-self-dual cusp classes
of the first type in (RNC.4) with level-$N^5$ vectors.
Let $S_2(e,M)$ be all cusp classes of the second type with
level-$M^3$ vectors.  They are finite and are permuted by
every automorphism of $\mathbb C$.  At a place prime to the
respective level, this permutation carries $A_v$ and
$\beta_v$ to their coefficientwise conjugates.  Moreover
\[
 |S_3(e,N)|\leq c_3[N(e+1)]^{90},\qquad
 |S_2(e,M)|\leq c_2[M(e+1)]^{12}.                              \tag{RNC.14}
\]
\begin{proof}
Finiteness and the inequalities follow from RNC6, since
$N^{90}(e+1)^{18}\leq[N(e+1)]^{90}$ and
$M^{12}(e+1)^4\leq[M(e+1)]^{12}$.
The labels are full global cusp classes, not eigenvalues
at a single prime.  Clozel's rationality and conjugation
theorem applies to these regular $C$-algebraic cusps with
the displayed matching coefficient.  Both embeddings of
$K$ have the same weight, and $\mathbb Q$ has one embedding,
so the infinity condition survives every field automorphism.
The principal fixed-vector condition survives scalar
conjugation as well.  Duality, conjugation of $K$, and
scalar conjugation commute on the finite representation;
strong multiplicity one then preserves conjugate
self-duality in the first catalogue.  In rank two every
class is essentially self-dual by the last argument of
RNC2.  Applying the inverse field automorphism gives the
inverse permutation of each finite catalogue.

There is a possible square-root trap in raw normalized
Satake parameters, so we record the exact calculation.
If $T_{v,j}$ is the unnormalized characteristic-function
Hecke operator for
$\operatorname{diag}(\varpi^{[j]},1^{[n-j]})$ at spherical
level, then
\[
       e_j(A_v)=q_v^{j(j-1)/2}\,T_{v,j}(\pi).                 \tag{RNC.15}
\]
Indeed normalized Satake gives
$T_{v,j}=q_v^{j(n-j)/2}e_j(\operatorname{Sat}_v)$;
substitution in (RNC.1) proves (RNC.15).  The exponent in
(RNC.15) is an integer.  The Hecke correspondence has
rational coefficients, so scalar conjugation applies to
its eigenvalue and fixes that rational power of $q_v$.
The monic characteristic polynomial consequently has
coefficientwise conjugate roots.  Inverting them proves
the same assertion for $\beta_v$.  This also proves that
all these roots are algebraic, since the finite parts have
number-field models.  No automorphism is required to fix
an arbitrarily chosen square root of $q_v$.
\end{proof}
% END RNC7

% BEGIN RNC8
\paragraph{RNC8. Purity and both integral endpoint bounds.}
\label{claim:RNC8}
For the first catalogue at a place $w\mid p$ with
$p\nmid N$, each arithmetic root $b\in\beta_w$ has
$|b|=q^{-1}$ and both
\[
          q^{e+1}b,\qquad q^{e-1}b^{-1}
                       \quad\text{algebraic integers}.       \tag{RNC.16}
\]
For the second catalogue at $p\nmid M$, each arithmetic
root $b\in\beta_p$ has $|b|=1$ and both
\[
                    p^e b,\qquad p^e b^{-1}
                       \quad\text{algebraic integers}.       \tag{RNC.17}
\]
These assertions hold after every complex field automorphism.
\begin{proof}
The indicated places are spherical because principal level
is maximal compact there.  Apply RNC1 at a coefficient
prime different from $p$.  The parameter has zero monodromy,
so purity gives absolute value $q$ for geometric roots in
rank three, and $1$ in rank two.  Inversion gives the
displayed values for $b$.  RNC7 gives the same argument for
every conjugate label and hence every algebraic conjugate.

To check integrality, fix an embedding of the algebraic
roots into $\overline{\mathbb Q}_a$ for each rational prime
$a$.  For $a\ne p$ use the representation of RNC1 at
coefficient prime $a$.  After ADK1 captures its compact
image in a finite extension, the image preserves an
integral lattice.  One elementary construction is to take
the module spanned by the translates of a starting lattice:
compactness bounds them in one larger lattice, which is
Noetherian, so this span is a lattice.  Frobenius and its
inverse preserve it.  Their eigenvalues are thus integral
units, proving all required integrality away from $p$.

At $a=p$, use the same cusp's crystalline representation
from RNC1.  For a weakly admissible filtered Frobenius
module whose Hodge--Tate weights lie in $[h_-,h_+]$, all
Newton slopes lie in that interval.  Here is the needed
argument: on the minimal-slope subobject,
$t_N\geq t_H\geq h_-\dim$; the maximal-slope quotient has
$t_N\leq t_H\leq h_+\dim$, by the subobject inequality
and equality for the whole module.  Passing to a splitting
coefficient field and to a residue-degree power of Frobenius
gives the same inequalities for each root, with valuation
normalized by $v(q)=1$.

In rank three the geometric interval is $[1-e,e+1]$;
the arithmetic interval after inversion is
$[-e-1,e-1]$.  This proves (RNC.16) at every prime above
$p$.  In rank two both intervals are $[-e,e]$, giving
(RNC.17).  Together with the preceding unit calculation
these bounds prove integrality at every finite place of
a number field containing the roots, which is equivalent
to algebraic integrality.  The argument works for every
embedding and every conjugate catalogue label.  In
particular it proves all the original $\mathrm{EigData.int}$
and $\mathrm{intInv}$ clauses: their permitted exponents
are respectively at least $e+1,e-1$ in rank three and
at least $e,e$ in rank two.
\end{proof}
% END RNC8

% BEGIN RNC9
\paragraph{RNC9. The rank-three trace catalogue classifies every later coefficient realization.}
\label{claim:RNC9}
For fixed $e,N,p,w$ as above let $q=\operatorname{N}w$.
The already chosen labels $S_3(e,N)$ carry traces
\[
                    t_3(s)=\sum_{b\in\beta_w(s)}(qb)^{-1}.
                                                               \tag{RNC.18}
\]
They satisfy $|t_3(s)|\leq3$, $q^e t_3(s)$ is an algebraic
integer, and every automorphism of $\mathbb C$ permutes
the traces with their labels.  For every later odd
$\ell\ne p$, closure, $\iota$, prolongation $B$ of $w$
and arithmetic Frobenius $\sigma$, every
$\mathrm{Adm3}(\iota,e,N,\rho)$ has a label with
\[
        \operatorname{tr}\rho(\sigma^{-1})=\iota^{-1}t_3(s).
                                                               \tag{RNC.19}
\]
Equal labels for two such representations imply that they
are isomorphic.
\begin{proof}
RNC8 gives three summands of absolute value one in
(RNC.18).  It also shows that
$q^e(qb)^{-1}=q^{e-1}b^{-1}$ is integral.  RNC7 gives
the exact permutation assertion because $q$ is rational.
RNC3 identifies $r=\rho\varepsilon^{-1}$ with the
representation attached to its actual cusp, and RNC5
puts that cusp into the already fixed label set.
At $w\nmid\ell$ full local compatibility gives
unramifiedness at every prolongation and the geometric
multiset $\beta_w(s)^{-1}$.  Since
\[
             r(\sigma^{-1})=q\rho(\sigma^{-1}),
                                                               \tag{RNC.20}
\]
the trace is exactly (RNC.19), for every permitted lift.
If the two cusp labels agree, both twisted representations
are isomorphic to the same attached semisimple
representation in the same coefficient closure; equivalently
one may apply CAP6 directly to their cofinite polynomials.
Untwisting proves isomorphism of the original $\rho$'s.
Thus the map classifies full representations, even when
distinct labels happen to have the same trace at $w$.
\end{proof}
% END RNC9

% BEGIN RNC10
\paragraph{RNC10. The rank-two catalogue gives the exact restricted block trace.}
\label{claim:RNC10}
Let $M\geq1$, $p\nmid M$, and $q=p^h$ with $h\geq1$.
The labels $S_2(e,M)$ carry traces
\[
                  t_2(s)=1+\sum_{b\in\beta_p(s)}b^{-h}.
                                                               \tag{RNC.21}
\]
They have absolute value at most $3$, $q^e t_2(s)$ is
integral, and all complex automorphisms permute them with
multiplicity.  For every later original coefficient datum,
every $\mathrm{Adm2}(\iota,e,M,\rho)$ and every original
arithmetic Frobenius $\sigma$ at a $K$-place of norm $q$,
there is a classifying label with
\[
         \operatorname{tr}\rho(\sigma^{-1})=\iota^{-1}t_2(s).
                                                               \tag{RNC.22}
\]
Equal labels imply isomorphism of the rank-three block
representations, including their original ordering of blocks.
\begin{proof}
By RNC8 each of the two $b$'s has modulus one, so (RNC.21)
has modulus at most $3$.  Its scaled summands are
\[
 q^e t_2(s)=p^{he}+\sum_{b\in\beta_p(s)}(p^e b^{-1})^h,
                                                               \tag{RNC.23}
\]
which are integral.  RNC7 gives the permutation property.

Choose the actual $W$ and the actual conjugating matrix in
the original $\mathrm{Adm2}$ witness.  RNC2, RNC3 and RNC5
give its label in $S_2(e,M)$.  As $p\nmid M$ and
$p\ne\ell$, $W$ is unramified at $p$ and its geometric
Frobenius roots are $\beta_p(s)^{-1}$.  The restriction of
an arithmetic Frobenius of a residue-degree-$h$ extension
to $G_{\mathbb Q}$ differs from the $h$-th power of a
rational arithmetic Frobenius by inertia.  Applying $W$
therefore gives precisely the $h$-th powers of the same
two geometric roots.  In the ordered source matrix the
two-dimensional block occupies positions $0,2$, and the
middle diagonal entry is $1$.  Its trace is exactly
(RNC.21); conjugation by the given matrix preserves trace.

If labels agree, CAP6 or the common attached representation
identifies the two $W$'s.  Restrict its actual intertwiner
to $G_K$ and insert it in positions $0,2$, with $1$ in the
middle position.  Conjugating by the two original block
identifications gives an intertwiner of the two $\rho$'s.
No irreducibility of $W|_{G_K}$ is needed, and no converse
assertion about isomorphic restrictions is used.
\end{proof}
% END RNC10

% BEGIN RNC11
\paragraph{RNC11. One finite trace catalogue handles the original disjunction.}
\label{claim:RNC11}
For $e,N\geq1$ and a prime $p$ with $p\nmid Nd$, choose
$w\mid p$, $q=\operatorname{N}w$.  Before choosing any
coefficient prime, there is a finite set $S$ and numbers
$t_s\in\mathbb C$ such that
\[
 \begin{gathered}
 |S|\leq c_3[N(e+1)]^{90}
                   +c_2[Nd(e+1)]^{12},\\
 |t_s|\leq3,\qquad q^e t_s\text{ is an algebraic integer},
 \qquad \forall\tau\in\operatorname{Aut}(\mathbb C),\
          \exists u_\tau:S\simeq S,\quad
          \tau(t_s)=t_{u_\tau(s)}.
 \end{gathered}                                                 \tag{RNC.24}
\]
For every later original coefficient datum and Frobenius,
the predicate
$\mathrm{Adm3}(\iota,e,N,\rho)\lor
\mathrm{Adm2}(\iota,e,Nd,\rho)$ is classified by these
labels with trace $\iota^{-1}t_s$ and equal-label
isomorphism.
\begin{proof}
The residue degree of $w$ is $h\in\{1,2\}$ because $K$
is quadratic, so $q=p^h$.  Take the disjoint union of the
rank-three catalogue of RNC9 and the rank-two catalogue
of RNC10 at $M=Nd$.  Both are available since $p\nmid Nd$.
All three arithmetic assertions in (RNC.24) and its
cardinality bound follow summand by summand.  For a
representation satisfying the disjunction, use its
rank-three classifier when the first predicate holds;
otherwise use a chosen rank-two witness and its
classifier.  Equal labels cannot lie in different
summands, and the corresponding individual classifier
gives isomorphism.  This convention also covers a
representation satisfying both predicates, without
assuming the two classes of representations are disjoint.
The sets and their traces were constructed entirely from
$K,e,N,p,w$, before $\ell$, $\Omega$ or $\iota$.
\end{proof}
% END RNC11

% BEGIN RNC12
\paragraph{RNC12. Every Eisenstein comparison factor is nonzero and integral.}
\label{claim:RNC12}
Let $q\geq7$, $e\geq1$, let $z$ be a root of unity,
and suppose $|t|\leq3$ and $q^e t$ is an algebraic integer.
For the exact original expression
\[
 E(q,e,z)=q^{-e}z+1+q^e z^{-1},\qquad
              F=q^e\bigl(t-E(q,e,z)\bigr),                    \tag{RNC.25}
\]
$F$ is a nonzero algebraic integer and
\[
                         \log|F|\leq2e\log q+1.
                                                               \tag{RNC.26}
\]
\begin{proof}
A root of unity has modulus one, and its inverse is also
an algebraic integer.  The reverse triangle inequality
gives
$|E(q,e,z)|\geq q^e-1-q^{-e}\geq7-1-1/7>3$.
Therefore $t\ne E(q,e,z)$ and $F\ne0$.
Expand
\[
             F=q^e t-z-q^e-q^{2e}z^{-1}.                      \tag{RNC.27}
\]
Every term is integral.  If $T=q^e\geq7$, the triangle
inequality gives
$|F|\leq T^2+4T+1\leq2T^2$, since
$1+4/7+1/49<2$.  Taking logarithms, using
$\log2\leq1$ and $\log T=e\log q$, proves (RNC.26).
This proof needs no lower bound on individual trace gaps
at a coefficient prime.
\end{proof}
% END RNC12

% BEGIN RNC13
\paragraph{RNC13. The whole orbit product is a nonzero rational integer.}
\label{claim:RNC13}
Let $(S,t)$ be the catalogue in (RNC.24), and let
$z_o$ be roots of unity indexed by any finite set $\mathcal O$,
with the original condition that every automorphism of
$\mathbb C$ permutes this indexed family, with multiplicity.
Then
\[
       NN=\prod_{o\in\mathcal O}\prod_{s\in S}
                   q^e\bigl(t_s-E(q,e,z_o)\bigr)\in\mathbb Z
                                                               \tag{RNC.28}
\]
is nonzero and satisfies
\[
      \log|NN|\leq |\mathcal O|\,|S|\,(2e\log q+1).
                                                               \tag{RNC.29}
\]
\begin{proof}
Every factor is a nonzero algebraic integer by RNC12, so
the finite product is a nonzero algebraic integer.  For
an arbitrary field automorphism, use its permutation of
$\mathcal O$ and the independent permutation of $S$ in
(RNC.24).  Since $q$ is rational and the formula uses only
integer powers and inverses, these permutations carry
the factors of (RNC.28) bijectively to their conjugates.
The whole product is fixed.  An algebraic complex number
fixed by every complex field automorphism is rational:
any nontrivial algebraic conjugate extends to such an
automorphism by extending a transcendence basis.  A
rational algebraic integer is in $\mathbb Z$.
Finally the logarithm of the norm of a finite product
of nonzero complex numbers is the sum of the logarithms;
RNC12 gives (RNC.29).  Empty index sets cause no exception:
the empty product is $1$ and the logarithmic bound is $0$.
No coefficient prime has yet been chosen.
\end{proof}
% END RNC13

% BEGIN RNC14
\paragraph{RNC14. Selecting pairwise nonisomorphic families gives the required local product inequality.}
\label{claim:RNC14}
Choose any of the later original data with an odd
$\ell\ne p$, and for each $o$ let $X_o$ be a finite
family of the representations classified by RNC11,
pairwise nonisomorphic within that family.  Then every
original trace difference is nonzero and
\[
 |NN|_\ell\leq
 \prod_{o\in\mathcal O}\prod_{\rho\in X_o}
 \left|\operatorname{tr}\rho(\sigma^{-1})-
                         \iota^{-1}E(q,e,z_o)\right|_\ell.
                                                               \tag{RNC.30}
\]
\begin{proof}
For this fixed coefficient datum and Frobenius, choose
the classifier from RNC11.  Its restriction to each
$X_o$ is injective: equal labels imply isomorphism, and
the hypothesis on $X_o$ then identifies the two members.
The map $(o,\rho)\mapsto(o,s_\rho)$ is therefore an
injection into $\mathcal O\times S$.  Repeating the same
representation for different $o$ is harmless because
the first coordinate keeps the chosen factors distinct.
RNC12 and the exact trace identity in the classifier
give nonvanishing of every displayed difference.

Apply $\iota^{-1}$ to (RNC.28).  A monic integral
polynomial and the ultrametric inequality imply that
every algebraic integer in $\overline{\mathbb Q}_\ell$
has norm at most one.  Thus each factor has norm at most
one.  Removing the factors outside the indicated
injection can only increase the product of norms.
Finally $q=p^h$ and $\ell\ne p$, so $|q^e|_\ell=1$.
The norm of each selected factor is precisely the
corresponding unscaled trace-gap norm in (RNC.30).
This proves the inequality for every family at once,
with the same already chosen integer $NN$.
\end{proof}
% END RNC14

% BEGIN RNC15
\paragraph{RNC15. The constants are chosen before all varying arithmetic data.}
\label{claim:RNC15}
There are natural constants $c,A$, depending only on $K$,
such that (RNC.29) is at most
\[
                c\,|\mathcal O|\,[N(e+1)]^A\log p.
                                                               \tag{RNC.31}
\]
One permissible choice is
\[
             A=103,\qquad c=5(c_3+c_2d^{12}),                 \tag{RNC.32}
\]
with the fixed constants of RNC6 enlarged to natural
integers if necessary.
\begin{proof}
Put $Y=N(e+1)\geq1$.  Formula (RNC.24) gives
\[
 |S|\leq c_3Y^{90}+c_2d^{12}Y^{12}
                    \leq(c_3+c_2d^{12})Y^{102}.               \tag{RNC.33}
\]
Here the deliberately nonoptimal exponent is the sum
of the two catalogue exponents.  Since $q=p^h$ with
$1\leq h\leq2$, $\log q\leq2\log p$.  For $p\geq7$,
$\log p\geq1$, and $e+1\leq Y$.  Consequently
\[
 2e\log q+1\leq4e\log p+\log p
                         \leq5Y\log p.                       \tag{RNC.34}
\]
Insert (RNC.33)--(RNC.34) into (RNC.29).  The degree,
discriminant and cell constants were fixed with $K$;
none of $e,N,p,w,\mathcal O,z$, or any later coefficient
data occurs in (RNC.32).
\end{proof}
% END RNC15

% BEGIN RNC16
\paragraph{RNC16. The literal original $\mathrm{CNR3}$.}
\label{claim:RNC16}
For every original imaginary quadratic field $K$, choose
$c,A$ as in RNC15.  For every $e,N\geq1$, every rational
prime $p\geq7$ with $p\nmid N$ and
$p\nmid\operatorname{disc}K$, every $w\mid p$, and
every original finite root-of-unity orbit with its
full complex-automorphism permutation condition, there
exists $NN\in\mathbb Z\setminus\{0\}$ satisfying
(RNC.31).  For every subsequent odd $\ell\ne p$,
every original coefficient closure and scalar tower,
every $\iota$, every prolongation of $w$, every
arithmetic Frobenius $\sigma$, and every family
$X_o$ of pairwise nonisomorphic representations satisfying
\[
 \mathrm{Adm3}(\iota,e,N,\rho)
          \quad\text{or}\quad
 \mathrm{Adm2}(\iota,e,N|\operatorname{disc}K|,\rho),
                                                               \tag{RNC.35}
\]
all its original trace gaps are nonzero and (RNC.30) holds.
\begin{proof}
RNC11 constructs the finite catalogue with exactly
(RNC.35).  RNC13 constructs the integer solely from that
catalogue and the given indexed orbit.  RNC15 proves the
required size bound with constants chosen first.
RNC14 proves the original nonvanishing conjunction and
product inequality for every later coefficient datum
and family, at geometric $\sigma^{-1}$ as required.
The rank-two discriminant factor occurs only in its
original $M=N|\operatorname{disc}K|$; the rank-three
level is still $N$.  The auxiliary automorphic levels
used for counting have not altered these predicates.
These are all the clauses and the literal quantifier
order of the original $\mathrm{CNR3}$.
\end{proof}
% END RNC16

The proof leaves the original named classical and period
foundations in place, with the bounded applications above
checked explicitly.  It supplies the complete congruence
integer producer.  It does not supply the one-character
length theorem: the subsequent $\mathrm{PropGK}$ consumer
still needs that distinct arithmetic producer.


\section{The finite installation of the congruence integer into PropGK}
\label{sec:prop-gk-installation-v046}

\status{This chapter proves the complete original implication
\(\mathrm{CNR3}\Longrightarrow(\mathrm{OneChar}\Longrightarrow
\mathrm{PropGK})\). RNC16 supplies its first premise. The original
single-character Selmer estimate \(\mathrm{OneChar}\) remains an
explicit premise; its rigidity and cohomological producers are not
proved here. The result discharges the finite implication task,
not the separate one-character producer.}

The source revision is
\texttt{7cbf7a0c32fca1badf05931879a37c3f8a2b03b4}.
We retain the original coefficient datum: for a rational prime
\(\ell\), \(O\subset\overline{\mathbb Q}_\ell\) is a finite
\(\mathbb Z_\ell\)-algebra which is a discrete valuation ring.
A Dirichlet character modulo \(N\ge1\) means a character of
\(G_N=(\mathbb Z/N\mathbb Z)^\times\), extended by zero off its
units. For \(H\le G_N\), set
\[
 \mathcal O_O(H)=\{\psi:G_N\to O^\times:\ker\psi=H\},\qquad
 \mathcal O_{\mathbb C}(H)=
       \{\chi:G_N\to\mathbb C^\times:\ker\chi=H\}.
 \tag{PGK.1}
\]
The first is exactly the original finite set \(\operatorname{orb}(c,H)\).

We also retain the actual source length, rather than substituting
an arbitrary nonnegative number. Let
\(\Psi=\varepsilon_\ell^e\psi_{\rm Gal}\), where
\(\psi_{\rm Gal}\) is the rational Galois character obtained by
acting on the \(N\)-th roots of unity. At an element of \(G_K\),
this character is restricted along the given map to \(G_{\mathbb Q}\).
Then
\(\mathcal L_\psi=V\!\operatorname{len}(\Omega,c,e,\psi)\) is the
\(O\)-length, possibly infinite, of the source Selmer module
\(\operatorname{ShaU}(\operatorname{twist}_2(\Psi^{-1}),\ell N,\top)\).
Its normalized real length is
\[
 v_\psi=V_n(\Omega,c,e,\psi)
       =\frac{\mathcal L_\psi.\operatorname{toNat}}{e_O},
 \qquad e_O=v_O(\ell)>0 .                                      \tag{PGK.2}
\]
We prove finiteness before using this as an ordinary length.
For a nonzero \(a\in\overline{\mathbb Q}_\ell\), write
\(v_\ell(a)=-\log\lVert a\rVert_\ell/\log\ell\), so
\(v_\ell(\ell)=1\).

Here is the exact retained \(\mathrm{OneChar}\) premise. For
every original \(\Omega,K,\ell,O,N,e,\psi\) satisfying
\(\mathrm{GD}\), every field isomorphism
\(\iota:\overline{\mathbb Q}_\ell\simeq\mathbb C\), every
\(\gamma\in G_K\), and
\[
 t=\Psi(\gamma)+1+\Psi(\gamma)^{-1}
                  \quad\hbox{in }\overline{\mathbb Q}_\ell,
 \tag{PGK.3}
\]
suppose \(\operatorname{tr}\rho(\gamma)\ne t\) for every actual
rank-three representation in
\(\mathrm{Adm3}(\iota,e,N)\) or
\(\mathrm{Adm2}(\iota,e,N|\operatorname{disc}K|)\).
Then \(\mathcal L_\psi<\infty\) and there is a finite set
\(X_\psi\) of such representations, pairwise nonisomorphic, with
\[
 v_\psi\le2\sum_{\rho\in X_\psi}
             v_\ell(\operatorname{tr}\rho(\gamma)-t).
 \tag{PGK.4}
\]
The scalar maps in (PGK.3) are the given inclusions
\(O\hookrightarrow\overline{\mathbb Q}_\ell\). The rank-two
alternative keeps the original two-by-two block in positions
\(0,2\) and its middle entry \(1\), as in RNC10. We assume exactly
this original premise, including its coefficient and local conditions.

% BEGIN PGK1
\paragraph{PGK1. The complete character orbit and its coefficient realization.}
\label{claim:PGK1}
Put \(h=[G_N:H]\). The complex set in (PGK.1) is empty or has
cardinality \(\varphi(h)\). If \(\psi_0\in\mathcal O_O(H)\) and
\(f:O\hookrightarrow\mathbb C\) is any injective ring map, then
\[
 E_f:\mathcal O_O(H)\overset\sim\longrightarrow
       \mathcal O_{\mathbb C}(H),\qquad
                     E_f(\psi)=f\circ\psi .                    \tag{PGK.5}
\]
Every member of \(\mathcal O_O(H)\) is \(\psi_0^j\) for some
\(j\) coprime to \(h\), and \(\operatorname{ord}(\psi_0)=h\).

\emph{Proof.}
A finite subgroup of the multiplicative group of a field is cyclic.
Indeed its exponent is attained by multiplying elements of maximum
prime-power order, and the polynomial \(X^d-1\) has at most \(d\)
roots; an exponent smaller than the group's order is impossible.
A character with kernel \(H\) therefore identifies \(G_N/H\)
with a cyclic group of order \(h\). Fix a generator \(\bar g\)
and put \(\zeta=\psi_0(g)\). Its order is \(h\).
The \(h\) powers of \(f(\zeta)\) are all roots of \(X^h-1\)
in \(\mathbb C\), by degree and distinctness. Any faithful
character of \(G_N/H\) consequently sends \(\bar g\) to
\(f(\zeta)^j\) with \(\gcd(j,h)=1\). It is precisely
\(f\circ\psi_0^j\). Such powers already take their values in
the given \(O\); no enlargement of the coefficient ring is needed.
They have distinct values for distinct invertible residue classes
\(j\pmod h\), and all have kernel \(H\).
Injectivity of \(f\) also proves that composition preserves the
kernel and distinguishes characters. This proves (PGK.5), its
cardinality, and the assertion about powers. If the complex set
has no member its cardinality is zero. The argument includes
\(h=1\), where there is one trivial character and
\(\varphi(1)=1\). \hfill\(\square\)
% END PGK1

% BEGIN PGK2
\paragraph{PGK2. Every character with the same kernel has the original GD data.}
\label{claim:PGK2}
If \(\mathrm{GD}(\Omega,K,\ell,e,\psi_0)\) and
\(\ker\psi=\ker\psi_0\), then
\(\mathrm{GD}(\Omega,K,\ell,e,\psi)\).

\emph{Proof.}
We check all the fields of the original record. Its conditions are
\[
 \begin{gathered}
 \ell\ne2,\quad e\ge1,\quad 2\mid N,\quad
 \ell\nmid\operatorname{ord}\psi_0,\quad
       (\operatorname{comp}_\ell\psi_0)^2=1,\quad
       \psi_0(-1)=(-1)^e,\\
 [K:\mathbb Q]=2,\quad K\text{ totally complex},\quad
 \ell\nmid|\operatorname{disc}K|\,h_K,\quad
       (|\operatorname{disc}K|,N)=1,\quad \ell-1\nmid2e,\\
 \ker\bar\varepsilon_\ell\cap
     \ker(\operatorname{chiP}_\ell(\psi_0)_{\rm Gal})
                                      \not\subseteq G_K .
 \end{gathered}                                                  \tag{PGK.6}
\]
Here \(h_K\) is the class number, \(G_K\) denotes the subgroup
fixing the given copy of \(K\) in \(\Omega\), and
\(\operatorname{comp}_\ell\) is restriction to the
\(\ell\)-primary factor of the modulus, followed by inflation.
Equivalently it is composition with the original Chinese-remainder
idempotent on \(G_N\), so it commutes with products and powers.

By PGK1 write \(\psi=\psi_0^j\), \(\gcd(j,h)=1\), where
\(h=\operatorname{ord}\psi_0\). Its order remains \(h\).
Since \(\ell\nmid h\), the source's \(\ell\)-power part
\(\operatorname{kap}_\ell(\psi_0)\) is trivial and its
prime-to-\(\ell\) part \(\operatorname{th0}_\ell(\psi_0)\)
is \(\psi_0\); the same holds for \(\psi\). This also follows
directly from their defining exponents: the \(\ell\)-part of
\(h\) is \(1\), so the defining power is a multiple of \(h\).

Set \(\delta=\operatorname{comp}_\ell\psi_0\). Both
\(\delta^2=1\) and \(\delta^h=1\). Any element of order dividing
both \(2\) and \(h\) is fixed by raising to \(j\): if its order
is \(2\), then \(h\) is even and \(j\) odd. Thus
\(\operatorname{comp}_\ell\psi=\delta^j=\delta\), preserving
the square condition. The same argument applied to
\(\psi_0(-1)\) preserves the parity condition.
Furthermore
\[
 \operatorname{chiP}_\ell(\psi_0)=\psi_0\delta^{-1},\qquad
 \operatorname{chiP}_\ell(\psi)
                 =(\psi_0\delta^{-1})^j .                      \tag{PGK.7}
\]
The order of \(\psi_0\delta^{-1}\) divides \(h\), so taking
its \(j\)-th power preserves its kernel, also after composition
with the rational cyclotomic action. The last condition in
(PGK.6) is therefore unchanged. Every remaining displayed
condition depends only on \(K,\ell,e,N\) or on the order
already checked. This proves the whole record.
No additional prime bound has been imposed. In particular
\(\ell=3\) is included in the later universal quantifier;
its original \(\ell-1\nmid2e\) antecedent is impossible.
\hfill\(\square\)
% END PGK2

% BEGIN PGK3
\paragraph{PGK3. The precise Frobenius and Eisenstein scalar.}
\label{claim:PGK3}
Let \(p\nmid N\), \(\ell\ne p\), \(w\mid p\) in \(K\), and
\(q=\operatorname{Norm}(w)\). Choose a prolongation \(B\) to
\(\Omega\) and an arithmetic Frobenius \(\sigma\in G_K\) at
\(B\). For the image \(g\in G_{\mathbb Q}\),
\[
 \operatorname{cls}_N(g)=q\pmod N,\qquad
           \varepsilon_\ell(g)=q\quad\hbox{in }\mathbb Z_\ell^\times.
 \tag{PGK.8}
\]
For any \(\psi\), any \(\iota\), and
\(\chi=\iota\circ\psi\), put \(z_\chi=\chi(q)^{-1}\).
Then the original scalar satisfies
\[
 \begin{split}
 \iota^{-1}\bigl(q^{-e}z_\chi+1+q^ez_\chi^{-1}\bigr)
 &=\Psi(g^{-1})+1+\Psi(g^{-1})^{-1}\\
 &=q^{-e}\psi(q)^{-1}+1+q^e\psi(q).
 \end{split}                                                     \tag{PGK.9}
\]

\emph{Proof.}
The residue field at \(w\) has cardinality \(q=p^f\) for
some \(f\ge1\), hence \((q,N)=1\). For every positive integer
\(m\) prime to \(p\), roots of unity of order dividing \(m\)
are unramified at \(p\), and reduction is injective on them.
One way to see both statements is that \(X^m-1\) is separable
on the residue field and all its roots lift uniquely over the
maximal unramified extension. Arithmetic Frobenius raises
their reductions, and therefore the roots themselves, to
the \(q\)-th power. Apply this first to \(m=N\), then to
every \(m=\ell^a\), and take the inverse limit. This proves
(PGK.8) with the original arithmetic convention, including
its restriction from \(K\) to \(\mathbb Q\).

Consequently \(\Psi(g)=q^e\psi(q)\). Evaluating the
inverse character at \(g^{-1}\) gives (PGK.9), since a field
isomorphism fixes rational numbers. All quantities on the
right use the original inclusion of \(O\), without a change
of coefficient datum.

The choices used here exist for every later datum. Valuations
extend to an algebraic closure, and the local decomposition
group surjects onto the absolute Galois group of the finite
residue field, giving a lift of its arithmetic Frobenius.
Also \(\Omega\), already an algebraic closure of \(\mathbb Q\)
containing \(K\), is an algebraic closure of \(K\). A field
isomorphism \(\overline{\mathbb Q}_\ell\simeq\mathbb C\)
exists: both algebraically closed characteristic-zero fields
have cardinality and transcendence degree over \(\mathbb Q\)
equal to the continuum; match transcendence bases and extend
to their algebraic closures. No continuity of this isomorphism
is asserted or required. \hfill\(\square\)
% END PGK3

% BEGIN PGK4
\paragraph{PGK4. A single integer before all coefficient choices.}
\label{claim:PGK4}
Assume \(K\) is imaginary quadratic. There are natural numbers
\(c,A\), depending only on \(K\), such that for all
\(e,N\ge1\), \(H\le G_N\), and prime \(p\) with
\(p\nmid30N|\operatorname{disc}K|\), there is
\(NN\in\mathbb Z\setminus\{0\}\) with
\[
 \log|NN|\le
       c\varphi([G_N:H])\,[N(e+1)]^A\log p.                    \tag{PGK.10}
\]
It simultaneously has the trace-separation and norm-product
properties of RNC16 at every later original coefficient datum,
with the index set \(\mathcal O_{\mathbb C}(H)\) and values
\(z_\chi=\chi(q)^{-1}\).

\emph{Proof.}
Choose \(c,A\) from RNC16 for this \(K\). The prime restriction
implies \(p\ge7\), \(p\nmid N\), and
\(p\nmid\operatorname{disc}K\): the smaller primes are
\(2,3,5\), all dividing \(30\). Choose any place \(w\mid p\)
and write \(q=\operatorname{Norm}(w)\).
The finite set \(\mathcal O_{\mathbb C}(H)\) is determined
before \(\ell\) and \(O\). Each \(\chi(q)\) is a root of
unity, because \((q,N)=1\). An automorphism \(a\) of
\(\mathbb C\) sends \(\chi\) to \(a\circ\chi\), preserving
its kernel and permuting this finite set, with inverse given
by \(a^{-1}\). It sends \(z_\chi\) to \(z_{a\circ\chi}\).
Thus this is exactly an admissible indexed orbit for RNC16,
including repetitions of the value at \(q\).

Apply RNC16 to obtain \(NN\ne0\), its full product inequality,
and the logarithmic bound with
\(|\mathcal O_{\mathbb C}(H)|\) in place of \(\varphi([G_N:H])\).
PGK1 bounds that cardinality by the latter number. All remaining
factors in the bound are nonnegative, giving (PGK.10).
This invocation precedes every choice of \(\ell,\Omega,O,\iota,B\),
Frobenius, or finite collection of representations. An empty
complex orbit causes no difficulty; the original RNC16 product
is then the empty product \(1\). \hfill\(\square\)
% END PGK4

% BEGIN PGK5
\paragraph{PGK5. Installation of the original single-character estimates.}
\label{claim:PGK5}
Retain the integer from PGK4 and assume the exact
\(\mathrm{OneChar}\) premise (PGK.3)--(PGK.4). For every later
\(\Omega,\ell\ne p,O\) such that at least one
\(\psi_0\in\mathcal O_O(H)\) has \(\mathrm{GD}\), all the
original lengths \(\mathcal L_\psi\),
\(\psi\in\mathcal O_O(H)\), are finite. There are finite
pairwise-nonisomorphic sets \(X_\psi\) of the original admissible
representations for which, writing
\[
 t_\psi=\Psi(g^{-1})+1+\Psi(g^{-1})^{-1},\qquad
 a_{\psi,\rho}=\operatorname{tr}\rho(\sigma^{-1})-t_\psi,
 \tag{PGK.11}
\]
one has
\[
 a_{\psi,\rho}\ne0,\qquad
 v_\psi\le2\sum_{\rho\in X_\psi}v_\ell(a_{\psi,\rho}),\qquad
 \lVert NN\rVert_\ell\le
       \prod_{\psi\in\mathcal O_O(H)}\prod_{\rho\in X_\psi}
                         \lVert a_{\psi,\rho}\rVert_\ell .      \tag{PGK.12}
\]

\emph{Proof.}
The assumed \(\mathrm{GD}\) gives \(\ell\ne2\), so RNC16
applies to this coefficient prime, without an additional
restriction. Choose \(\iota,B,\sigma\) as in PGK3. Use PGK1
with the injective map
\(O\hookrightarrow\overline{\mathbb Q}_\ell
\xrightarrow{\iota}\mathbb C\) to obtain the bijection \(E\)
of the two complete character orbits. PGK2 supplies
\(\mathrm{GD}\) for every \(\psi\) in the coefficient orbit.

For each such \(\psi\), (PGK.9) identifies \(t_\psi\) with
\(\iota^{-1}E(q,e,z_{E\psi})\), using exactly the Eisenstein
expression in RNC16. Its universal nonvanishing conclusion
gives the antecedent in (PGK.3) for \(\gamma=\sigma^{-1}\).
To see that this precedes the choice of a finite set, apply
RNC16 with the singleton \(\{\rho\}\) at each orbit index
for an arbitrary admissible \(\rho\). Each such family meets
its pairwise-nonisomorphism condition. Thus every actual
admissible \(\rho\) has trace different from this scalar.

The stated \(\mathrm{OneChar}\) premise now yields
\(\mathcal L_\psi<\infty\) and a set \(X_\psi\) satisfying
its exact admissibility, nonisomorphism and estimate. These
choices are made for finitely many \(\psi\). Apply the product
clause of the same RNC16 integer to the family
\(X_{E^{-1}\chi}\) indexed by
\(\chi\in\mathcal O_{\mathbb C}(H)\). Its hypotheses are
exactly the properties just supplied by \(\mathrm{OneChar}\).
Reindexing through \(E\) and using (PGK.9) gives the last
inequality of (PGK.12), and its nonvanishing clause gives the
first. A representation may occur for several characters;
the original product retains its character index, so no
cross-character nonisomorphism is required. \hfill\(\square\)
% END PGK5

% BEGIN PGK6
\paragraph{PGK6. The whole original PropGK implication.}
\label{claim:PGK6}
Assuming the original \(\mathrm{OneChar}\), RNC16 proves the
literal \(\mathrm{PropGK}\): for every number field \(K\) there
are \(c_5,A\in\mathbb N\) such that, for every
\(e,N\ge1\), every \(H\le G_N\), and every prime
\(p\nmid30N|\operatorname{disc}K|\), a nonzero integer \(NN\)
satisfies (PGK.10) with \(c=c_5\), and, for every later
\(\Omega,\ell\ne p,O\),
\[
 \begin{split}
 &(\exists\psi_0\in\mathcal O_O(H),\ \mathrm{GD}(\psi_0))\\
 &\quad\Longrightarrow
   \left(\mathcal L_\psi<\infty\text{ for every }\psi\in
                        \mathcal O_O(H),\quad
     \sum_{\psi\in\mathcal O_O(H)}v_\psi
                           \le2v_\ell(NN)\right).
 \end{split}                                                     \tag{PGK.13}
\]
Here \(v_\ell(NN)=\operatorname{padicValInt}(\ell,NN)\).

\emph{Proof.}
First suppose \(K\) is imaginary quadratic. Choose the
constants and the integer by PGK4. For any later coefficient
datum with the displayed antecedent, PGK5 gives all the
finiteness assertions and (PGK.12). Every factor there is
strictly positive in norm, and \(NN\ne0\). We may therefore
take logarithms to base \(\ell>1\), obtaining
\[
 -v_\ell(NN)\le
       -\sum_{\psi\in\mathcal O_O(H)}\sum_{\rho\in X_\psi}
                            v_\ell(a_{\psi,\rho}).
 \tag{PGK.14}
\]
The identity on the left follows directly from the usual
integer factorization \(NN=\ell^r u\), \(\ell\nmid u\):
its norm is \(\ell^{-r}\). The logarithm of each finite
product is the sum of the logarithms, also for empty sets.
Multiplying (PGK.14) by \(-1\), summing (PGK.4), and
multiplying by \(2\) gives
\[
 \sum_\psi v_\psi
 \le2\sum_\psi\sum_{\rho\in X_\psi}v_\ell(a_{\psi,\rho})
 \le2v_\ell(NN).
\]
This proves exactly (PGK.13) for the same integer chosen
before the coefficient prime. No sign assumption on the
individual valuations was needed for the logarithmic step.

For any number field \(K\) which is not totally complex of
degree two, choose \(c_5=A=0\) and \(NN=1\) for all the
varying data. The logarithmic bound is \(0\le0\). The later
antecedent is impossible, since the \(\mathrm{quad}\) and
\(\mathrm{imag}\) fields in (PGK.6) assert precisely the
excluded conjunction. A totally complex quadratic field is
CM over \(\mathbb Q\), so these two cases exhaust all
number fields. This completes the literal original quantifier
order and the whole finite implication. Its only unsupplied
arithmetic premise is the original \(\mathrm{OneChar}\);
in particular this proof does not establish its type-II
rigidity producer. \hfill\(\square\)
% END PGK6

\subsection{Scope of the completed implication}

The construction has order
\[
 K\ \longmapsto\ (c_5,A)\ \longmapsto\ (e,N,H,p)
 \ \longmapsto\ NN\ne0\ \longmapsto\ (\Omega,\ell,O).
\]
The coefficient orbit and the sets in the one-character
estimate are chosen only after \(NN\). The two complete
orbits are identified by (PGK.5), and the arithmetic versus
geometric Frobenius sign is fixed by (PGK.9). Thus the
finite orbit/Frobenius/valuation task is proved with its
original hypotheses. The separate production of
\(\mathrm{OneChar}\), including the original
\(\mathrm{RigII}\) and other arithmetic inputs, remains outside
this component.


\section{The actual adelic projectors and one common arithmetic row family}
\label{sec:arithmetic-rows-v46}

This chapter constructs the row object on which the retained-weight
and cylinder-count arguments operate.  The input is one original
datum \(D\) and one original adelic \(\mathrm{SG.HidaSetting}(D)\),
denoted by \(X\).  This is the setting of lattices, adelic measures,
discrete-series representations and row operators.  It is a different
record from the ordinary Hecke-algebra setting used in the preceding
Hida chapters.  Nothing here identifies their objects merely because
both records have the word Hida in their name.

Write \(n=D.R.r\), \(f=D.R.f\), \(s=D.R.rS=D.R.m/2-1\),
\(k_0=D.R.k\), and \(q_j=D.\mathrm{host}.g(j)\), for \(j=0,1\).
Thus \(n=2s+k_0+1\), and \(q_0,q_1\) are distinct primes splitting
completely in the original CM field.  Transport through \(X.\iota_E\)
is understood throughout.  Put \(g=q_0q_1\), let \(K\) be the
original finite-adelic lattice stabilizer, and let
\[
 N=K(g)=\{x\in K:(x-1)\widehat L\subseteq g\widehat L\},
 \qquad
 \mathcal G=\prod_{j=0}^1\prod_{v\in D.\mathrm{nd}.\mathrm{PlG}(j)}
                  \mathrm{GL}_n(\mathbb F_{q_j}).
 \tag{ARO.1}
\]
Every measure, fundamental domain, infinity representation, frame,
and lattice in this chapter is the corresponding object of \(X\).
In particular \(\mu=X.B.\mu_f\) is never renormalized.  Let
\(V=\mu(N)\), with finite measures identified with their real values.
The finite labels are the actual monic degree-\(n\) polynomials with
nonzero constant coefficient; \(E_{j,P}=D.\mathrm{ntypes}.eSer(j,P)\)
are their original central orthogonal idempotents.  An invalid
polynomial has idempotent zero.  The normalized-type record, the
given setting \(X\), and, where expressly used below, the original
\(\mathrm{RowsFinDim}\), \(\mathrm{GapHyp}\), \(\mathrm{AutTie}\),
\(\mathrm{RowDictTie}\), and \(\mathrm{LabTie}\) are retained inputs.
Their simultaneous existence in a chosen global family is not assumed
to follow from this chapter.

\subsection{Reduction, the level character, and the exact Haar factor}

% BEGIN ARO1
\paragraph{ARO1. Reduction on the original lattice and its exact kernel.}
\label{claim:ARO1}
Choose once, for each selected place \(w=w(j,v)\), a basis of
\(L/wL\), and lift it to \(L\).  The action on these quotients defines
a homomorphism \(\rho:K\to\mathcal G\) with
\[
                         \ker\rho=N.                 \tag{ARO.2}
\]
The same assertion for one auxiliary prime says that the kernel of
the corresponding factor of \(\rho\) is \(K(q_j)\).
\begin{proof}
The localization of a full lattice at a finite prime is a free module
of rank \(n\) over its discrete valuation ring.  Reduction modulo its
maximal ideal has dimension \(n\), so the asserted bases exist.
The ordinary lattice maps onto this reduction: a localized
denominator is invertible modulo the maximal ideal and can be
replaced by its inverse in the residue field.  Since that field is
\(\mathbb F_{q_j}\), its coefficients can also be lifted by integers.
Every \(x\in K\) and its inverse preserve the completed lattice;
their reduced matrices are inverse.  Reducing products proves the
homomorphism assertion.  A different lifted basis conjugates the
reduced matrix and does not affect whether it is the identity.

Here are the two points needed to identify the kernel on the whole
adelic lattice.  At a selected split prime, the original frame and
dual frame have pairing matrix the identity modulo that prime.
This follows directly from \(\mathrm{gramIntegral}\),
\(\mathrm{formsAtZero}\), \(\mathrm{bw\_\allowbreak dual}\), and the
complementary residues of \(\mathrm{ew}\) and its conjugate.
Indeed an idempotent in a residue field is zero or one;
\(\mathrm{ew\_\allowbreak orth}\) and \(\mathrm{ew\_\allowbreak full}\) make the two
residues complementary.  Pairing a linear relation with the dual
vectors proves independence, and \(\mathrm{latticeSpan}\) proves
spanning.  Thus the pairing of the two conjugate local lattices is
perfect.  Unitarity identifies their action matrices as inverse
transposes in dual bases.  Identity reduction at \(w\) therefore
implies identity reduction at its conjugate \(\bar w\).

The selected places and their conjugates exhaust all primes above
\(q_0q_1\).  Complete splitting means these primes are distinct and
occur to exponent one.  In each local free lattice, zero reduction
of \((x-1)y\) means membership in its maximal-ideal multiple.
At the remaining primes \(g\) is a unit.  Consequently identity
reduction in every selected factor is equivalent to
\((x-1)\widehat L\subseteq g\widehat L\), component by component.
This proves (ARO.2); the proof with only the primes above \(q_j\)
gives the last assertion.  The argument uses the original completed
lattice, so it does not replace adelic congruence by congruence of a
chosen collection of rational vectors.
\end{proof}
% END ARO1

% BEGIN ARO2
\paragraph{ARO2. Surjectivity of the actual reduction map.}
\label{claim:ARO2}
The homomorphism of ARO1 is surjective.  Hence
\(K/N\simeq\mathcal G\), with exactly \(f\) factors above each
auxiliary prime.
\begin{proof}
At a completely split finite place the CM algebra is a product of
two copies of the base local field, interchanged by conjugation.
The unitary equation determines the action on the second factor
from its action on the first.  In the dual local bases supplied in
ARO1 it reads \(A'=(A^{-1})^{\mathsf t}\).  The lattice stabilizer
therefore identifies with \(\mathrm{GL}_n(\mathcal O_v)\).
An arbitrary matrix in \(\mathrm{GL}_n(\mathbb F_{q_j})\) lifts by
lifting its entries: its determinant remains a unit, so its inverse
also has integral entries.  This constructs an actual local unitary
lattice automorphism with the prescribed reduction at that place.

Do this at the finitely many selected places and take the identity
at every other finite place.  The resulting tuple is a finite adele,
satisfies the unitary equation at every place and preserves every
local lattice.  It is therefore an element of the original \(K\).
Changing from these convenient local bases to the bases defining
\(\rho\) only conjugates each target matrix; prescribe the conjugated
matrix when making the lift.  All factors can be prescribed
independently.  This proves surjectivity without a rational strong
approximation assertion.  ARO1 now gives the quotient isomorphism.
The number of factors is the original equality
\(|D.\mathrm{nd}.\mathrm{PlG}(j)|=f\).
\end{proof}
% END ARO2

% BEGIN ARO3
\paragraph{ARO3. Haar volume on each reduction fiber.}
\label{claim:ARO3}
One has \(0<V<\infty\).  Each fiber of \(\rho\) has measure \(V\).
If \(H\leq\mathcal G\), then
\[
 \mu(\rho^{-1}H)=|H|V,
 \qquad \mu(K)=|\mathcal G|V.                         \tag{ARO.3}
\]
\begin{proof}
The congruence subgroup \(N\) is open and contains the identity.
It is a closed subgroup of the compact lattice stabilizer \(K\).
A nonzero Haar measure is positive on every nonempty open set and
finite on compact sets, proving the inequalities.  By ARO1--ARO2
a reduction fiber is one coset \(xN\); left invariance gives measure
\(V\).  The fibers are disjoint measurable sets.  Finite additivity
over the fibers in \(H\) gives (ARO.3).  In particular no counting
measure of total mass one has been substituted for \(\mu\).
\end{proof}
% END ARO3

% BEGIN ARO4
\paragraph{ARO4. Inflation is independent of the reduction basis.}
\label{claim:ARO4}
For a central element \(a\in\mathbb C[\mathcal G]\), the original
function \(\mathrm{egA}(X,a)\) is exactly
\[
 \operatorname{inf}(a)(x)=
 \begin{cases}a(\rho(x))/V,&x\in K,\\0,&x\notin K.\end{cases}
 \tag{ARO.4}
\]
It is locally constant, compactly supported and bi-invariant under
\(N\).  For the original finite event idempotent \(E_S\), this
applies simultaneously to every set of tuples \(S\).
\begin{proof}
The original definition selects any tuple of reduction matrices in
any admissible lifted bases.  Such bases really are bases: their
spanning condition on a rank-\(n\) quotient of dimension \(n\)
implies independence.  Two choices are related by an invertible
change of basis in each factor, so their tuples are conjugate in
\(\mathcal G\).  Coefficients of a central group-algebra element are
constant under conjugacy.  ARO1 supplies one tuple for each \(x\in K\),
and the original denominator is precisely \(V\).  This proves
(ARO.4), including the zero value outside \(K\).
Finite fibers are open by ARO1; their union is \(K\), and \(K\)
is also closed.  This proves local constancy and compact support.
Multiplication on either side by \(N=\ker\rho\) changes no value.
RRT12 proves centrality of every actual \(E_S\), so the conclusion
holds for all these events with the same \(\rho\) and \(V\).
\end{proof}
% END ARO4

% BEGIN ARO5
\paragraph{ARO5. The exact finite convolution and involution comparison.}
\label{claim:ARO5}
For a function \(h\) supported in \(K\) and bi-invariant under
\(N\), define \(\widehat h(\gamma)=Vh(x_\gamma)\), where
\(\rho(x_\gamma)=\gamma\).  Then this definition is independent of
the representatives, is injective on such functions, and satisfies
\[
 \widehat{h_1*h_2}=\widehat h_1\widehat h_2,
 \qquad
 \widehat{h^*}=\widehat h^{\,*},
 \qquad h^*(x)=\overline{h(x^{-1})}.
 \tag{ARO.5}
\]
\begin{proof}
Representatives of a fiber differ by \(N\), by ARO1.
The values of \(h\) on \(K\) can be recovered by dividing
\(\widehat h\) by the positive number \(V\), proving injectivity.
For convolution partition its domain of integration into reduction
fibers.  At \(x_\gamma\) the integrand on \(x_\delta N\) is the
constant \(h_1(x_\delta)h_2(x_\delta^{-1}x_\gamma)\).  ARO3 gives
\[
 (h_1*h_2)(x_\gamma)
 =V\sum_{\delta\in\mathcal G}
           h_1(x_\delta)h_2(x_{\delta^{-1}\gamma}).
\]
Multiplication by the outer \(V\) gives the group-algebra product;
there are two factors \(V\), one from the integral and one from the
transform.  If \(x\notin K\), both nonzero factors in the integrand
would force \(x\in K\), so the convolution is zero.  The convolution
is again bi-invariant.  Inversion sends a fiber to the inverse
fiber, and \(V\) is real, proving the involution formula.  All
integrals here are finite sums of bounded functions on compact sets.
By ARO4, \(\widehat{\operatorname{inf}(a)}=a\).
\end{proof}
% END ARO5

% BEGIN ARO6
\paragraph{ARO6. The finite event algebra with its invalid-label convention.}
\label{claim:ARO6}
Let \(\mathcal T\) be the finite set of valid two-prime tuples.
For any set \(S\) of polynomial tuples, the literal \(E_S\) satisfies
\[
 E_SE_T=E_{S\cap T},\quad E_S^*=E_S,\quad E_{\mathrm{all}}=1,
 \quad E_S=\sum_{t\in\mathcal T\cap S}E_{\{t\}}.
 \tag{ARO.6}
\]
In particular arbitrary sets cause no infinite-sum issue.
\begin{proof}
A monic polynomial of degree \(n\) is uniquely determined by its
\(n\) coefficients below the leading coefficient.  Conversely each
coefficient tuple gives that polynomial.  The original definition of
\(E_S\) is consequently the sum of tensor products of the local
\(E_{j,P}\), indexed once by each monic tuple in \(S\).  Invalid
labels have a zero factor.  Removing them leaves exactly the finite
sum in (ARO.6).  RRT12 gives the product and total-sum identities.

For clarity self-adjointness follows even if it is not included as
a separate field of the normalized-type record.  In the left regular
representation a central element acts by a scalar on each irreducible
summand.  A central idempotent has scalar zero or one.  The regular
representation is unitary for its standard finite counting inner
product and faithful, so this idempotent equals its adjoint.
Tensor products and finite sums preserve that equality.  This proves
the involution assertion for all original event projectors.
\end{proof}
% END ARO6

% BEGIN ARO7
\paragraph{ARO7. The original frame gives dual bases at both orientations.}
\label{claim:ARO7}
At a selected place let \(b\in\{0,1\}\) be the residue of
\(X.\mathrm{ew}\).  A residue basis, in four consecutive blocks,
is
\[
 (e_0,\ldots,e_{s-1};\ c_0,\ldots,c_{k_0-1};\ z;
                         f_{s-1},\ldots,f_0),        \tag{ARO.7}
\]
where \(c_i=\beta_i\) if \(b=1\), and
\(c_i=\beta^d_{k_0-1-i}\) if \(b=0\).
Its dual frame is obtained by replacing \(e\) by \(f\),
\(z\) by \(-z\), and \(c\) by the complementary \(\beta\)-frame
with the same reversal.  These statements concern the original
frame vectors modulo the place.
\begin{proof}
The complementary-residue argument in ARO1 gives
\(\mathrm{ew}\equiv b\) and \(\overline{\mathrm{ew}}\equiv1-b\).
The identities \(\beta_i=\mathrm{ew}\,\beta_i\) and
\(\beta_i^d=\overline{\mathrm{ew}}\,\beta_i^d\) make one of the
two \(\beta\)-frames zero at this place and the other the active
frame.  The congruences in the setting give
\(h(e_a,f_c)=\delta_{ac}\), \(h(z,z)=-1\), and the corresponding
active dual \(\beta\)-pairings \(\delta_{ij}\); all off-block
pairings vanish.  Pairing (ARO.7) with its displayed dual therefore
gives the identity matrix.  Independence follows by pairing any
linear relation with the dual frame.  The original lattice-spanning
condition makes the active frame span, so it is a basis, not just an
independent subspace.  Reversal in the zero orientation is necessary
because conjugation interchanges the active and dual flags.
The changes from the basis of ARO1 to these bases are invertible
matrices fixed before the group element is varied.
\end{proof}
% END ARO7

% BEGIN ARO8
\paragraph{ARO8. Exact reduction of the concrete level subgroup.}
\label{claim:ARO8}
Let \(J_q\) be the subgroup of upper unitriangular
\(n\times n\) matrices over \(\mathbb F_q\) satisfying
\[
                  M_{a,s+k_0}=0\quad(s\leq a<s+k_0).
 \tag{ARO.8}
\]
There are matrices \(P_{j,v}\), fixed before \(x\), such that,
for every \(x\in K\),
\[
 x\in X.K_2\quad\Longleftrightarrow\quad
 P_{j,v}^{-1}\rho(x)_{j,v}P_{j,v}\in J_{q_j}
                    \quad\hbox{for every }j,v.       \tag{ARO.9}
\]
\begin{proof}
In the basis of ARO7 the defining level congruences say that
\((x-1)e_a\) is in the span of earlier \(e\)'s,
\((x-1)z\) is in their full span, and
\((x-1)\beta_i\) is in that span together with the earlier
\(\beta\)'s.  Thus, in the orientation \(b=1\), the \(e\),
active \(\beta\), and \(z\) columns have zero entries below the
diagonal, diagonal entry one, and the extra zero entries (ARO.8).
Applying the conjugate-place conditions and the identity
\(h(u,xv)=h(x^{-1}u,v)\) supplies the same conditions on the
reversed \(f\)-block.  In the orientation \(b=0\), the active
\(\beta^d\)-block is controlled by those conjugate inverse
conditions; the reversal in ARO7 turns them into upper-triangular
conditions.  The \(e\) and \(z\) columns are controlled directly.
For the \(z\) column the coefficients in the active middle block
are zero in both orientations.  These are exactly all entries below
the diagonal, all diagonal entries, and (ARO.8).

The converse can be checked without a volume-index argument.
Upper unitriangular matrices satisfying (ARO.8) form a subgroup:
the only possible intermediate terms in the forbidden column are
other forbidden entries.  Hence their inverses have the same shape.
The displayed shape gives the required flag inclusions for the
\(e\) and \(z\) vectors directly.  For \(b=1\) it gives them for
\(\beta\) directly; for \(b=0\) those vectors vanish at this
place and the conjugate-place inverse matrix supplies their
conditions.  Applying this at both places in every conjugate pair
proves the original inclusions modulo every prime above \(g\).
Chinese remaindering for the completed lattice then gives the
original inclusions modulo \(g\), with their original spans.
Outside these primes \(g\) is a unit and imposes no condition.

Finally multiply the basis by a diagonal matrix equal to one on
the \(e\) and \(z\) blocks, to \((-1)^i\) on the reversed
\(f\)-block, and also to \((-1)^i\) on the middle block when
\(b=0\).  This preserves the zero pattern and diagonal entries.
It is the sign adjustment used for the character below.  Compose
it with the fixed basis change of ARO7 to obtain \(P_{j,v}\).
This proves both directions of (ARO.9) on the original subgroup.
\end{proof}
% END ARO8

% BEGIN ARO9
\paragraph{ARO9. Factorization and nontriviality of the actual level character.}
\label{claim:ARO9}
For each selected place there is a nontrivial additive character
\(\psi_{j,v}:\mathbb F_{q_j}\to\mathbb C^\times\) such that, on
the subgroup of ARO8,
\[
 X.\chi_2(x)=\prod_{j,v}\psi_{j,v}
   \bigl(\phi_J(P_{j,v}^{-1}\rho(x)_{j,v}P_{j,v})\bigr).
 \tag{ARO.10}
\]
Here \(\phi_J\) is exactly the original sum of entries: the
successive entries in the \(e\)-chain, the two links through \(z\),
the successive entries in the reversed \(f\)-chain, and the
successive entries in the middle block.  Empty chains contribute zero.
\begin{proof}
Use Chinese remaindering in \(\mathcal O_E/g\mathcal O_E\) to
choose \(\varepsilon_{j,v}\) equal to one at \(w(j,v)\) and zero
at every other prime above \(g\), including its conjugate.
It is divisible by \(q_{1-j}\).  Define
\[
 \psi_{j,v}(a)=\psi_f(\varepsilon_{j,v}\widetilde a/g),
 \tag{ARO.11}
\]
where \(\widetilde a\) is any integer lift.  Changing the lift
adds an integral adele because \(q_{1-j}\mid\varepsilon_{j,v}\).
The given additive character is trivial on such adeles, proving
well-definedness and additivity.  Its value on the conjugate CRT
idempotent is the same, by the original rational/infinity character
identity and invariance of the real trace under CM conjugation.

Nontriviality can be checked by a trace calculation in the original
field.  Put \(a=\varepsilon_{j,v}/q_{1-j}\in\mathcal O_E\).
Modulo \(q_j\), its absolute trace is the sum of its residues at
all primes above \(q_j\), since the integral algebra there is the
product of copies of \(\mathbb Z_{q_j}\).  This sum is
\(q_{1-j}^{-1}\neq0\) in \(\mathbb F_{q_j}\).  The rational
character identity therefore gives
\(\psi_{j,v}(1)=\exp(-2\pi i\operatorname{Tr}_{E/\mathbb Q}(a)/q_j)
\neq1\).  This uses complete splitting, already part of the
original datum, and does not posit a new conductor normalization.

Let \(C(x)\) be the original argument of \(\chi_2\), namely
\[
 C(x)=\sum_a h(f_a,x\,\operatorname{next}(e,z)_a)
                 +\sum_{i<k_0-1}h(\beta_i^d,x\beta_{i+1}).
\]
It is integral because all frame vectors and their transforms are
in the integral lattice.  In the frame of ARO7 its residue at
\(w\) is the sum of the direct \(e\)-chain entries and, in
orientation one, the middle-chain entries.  At \(\bar w\), use
\(\overline{h(u,xv)}=h(v,x^{-1}u)\).  This gives the reversed
\(f\)-chain entries of the inverse matrix, a minus sign at the
\(z\)-link because the dual of \(z\) is \(-z\), and, in
orientation zero, the reversed middle-chain entries of the inverse.
For upper unitriangular matrices the first superdiagonal of the
inverse is the negative of that of the matrix.  The diagonal signs
at the end of ARO8 change these negatives into the positive
successive entries in \(\phi_J\).  The two links crossing the
middle block have no extra product term, precisely because the
forbidden column in (ARO.8) is zero.  Thus the sum of the two
conjugate residues of \(C(x)\) is the stated \(\phi_J\).

Subtract from \(C(x)\) the sum of its residues times the CRT
idempotents and their conjugates.  The difference lies in
\(g\widehat{\mathcal O}_E\).  Divide by \(g\) and apply \(\psi_f\).
The difference contributes one, and additivity, (ARO.11), and the
equal conjugate characters give (ARO.10).  Every sign and every
local factor is therefore tied to the original character.
\end{proof}
% END ARO9

% BEGIN ARO10
\paragraph{ARO10. The level volume and its finite character idempotent.}
\label{claim:ARO10}
Put \(c_j=|J_{q_j}|\), and let
\[
 a_{j,v}=P_{j,v}e_J(q_j,n,k_0,s,\psi_{j,v})P_{j,v}^{-1}
                 \quad\hbox{in }\mathbb C[\mathrm{GL}_n(\mathbb F_{q_j})].
\]
Then, for the original function \(e=X.B.e\),
\[
 \mu(X.K_2)=c_0^f c_1^f V,
 \qquad \widehat e=\bigotimes_{j,v}a_{j,v},
 \qquad e(1)=\frac1{c_0^f c_1^f V}>0.                 \tag{ARO.12}
\]
\begin{proof}
By ARO2 and ARO8 the reduction image of \(X.K_2\) is the full
product of the conjugated subgroups \(J_{q_j}\), and its kernel is
\(N\).  ARO3 gives the volume formula.  The identity belongs to
each \(J_q\), so \(c_j>0\).  The formula defining the original
\(e\) is \(\overline{\chi_2}/\mu(X.K_2)\) on \(X.K_2\),
and zero off it.  ARO9 expresses this character as the product of
the local \(\psi_{j,v}\circ\phi_J\).  Its conjugate is its
inverse, because finite additive characters have absolute value
one.  Multiplying by \(V\) and using the volume formula gives
exactly the normalized subgroup-character idempotents in (ARO.12).
Their values at the identity are \(1/c_j\), proving the last
formula.  The tensor equality is a coefficient equality at every
element, including those outside the level image.
\end{proof}
% END ARO10

% BEGIN ARO11
\paragraph{ARO11. Independence of the chosen nontrivial finite character.}
\label{claim:ARO11}
If \(Z\) is central in \(\mathbb C[\mathrm{GL}_n(\mathbb F_q)]\)
and \(\psi,\psi'\) are nontrivial additive characters, then
\[
 [1](e_J(\psi)Z)=[1](e_J(\psi')Z),
 \qquad [1]e_J(\psi)=[1]e_J(\psi')=1/|J_q|.
 \tag{ARO.13}
\]
\begin{proof}
The additive group of the prime field is cyclic of prime order.
Every nontrivial character is faithful and every character is
determined by its value at one.  Thus
\(\psi'(x)=\psi(cx)\) for some \(c\in\mathbb F_q^\times\).
Assign the following exponents to the four blocks of ARO7:
\[
 0,\ldots,s-1;\quad 0,\ldots,k_0-1;\quad s;\quad
                                  s+1,\ldots,2s.
\]
Let \(T\) be diagonal with entries \(c\) to these powers.
Every entry occurring in \(\phi_J\) has exponent difference one.
Therefore \(\phi_J(T^{-1}MT)=c\phi_J(M)\).  Diagonal conjugation
preserves both the upper unitriangular shape and its forbidden
zero column, so it permutes \(J_q\).  Reindexing the character
average shows that the two idempotents are conjugate by \(T\)
(with the direction reversed if the inverse convention is used).
The identity coefficient of a conjugate is unchanged, and \(Z\)
commutes with \(T\), proving the first equality.  At the identity
the character value is one, proving the second.  The displayed
exponents also cover \(s=0\) and an empty middle chain.
\end{proof}
% END ARO11

% BEGIN ARO12
\paragraph{ARO12. The original cylinder identity including its Haar normalization.}
\label{claim:ARO12}
Choose any nontrivial \(\psi_j^0\) on each \(\mathbb F_{q_j}\),
and put \(e_T(S)=e*\operatorname{inf}(E_S)\).  For every finite
set \(S\) of one-prime tuples,
\[
 \operatorname{Re}e_T(\operatorname{cyl}(j,S))(1)
   =e(1)\sum_{t\in S}\prod_v
      \frac{\operatorname{Re}[1](e_J(\psi_j^0)E_{j,t_v})}
           {\operatorname{Re}[1]e_J(\psi_j^0)}.
 \tag{ARO.14}
\]
This is the full original \(\mathrm{EtOne}\), on the given measure.
\begin{proof}
ARO4--ARO5 give \(\widehat{e_T(S)}=\widehat e E_S\).
ARO10 identifies \(\widehat e\) as the actual tensor product of
conjugated local character idempotents.  In the identity coefficient
of its product with a central label projector, the conjugating
matrix disappears.  ARO11 then replaces every local \(\psi_{j,v}\)
by the chosen \(\psi_j^0\).  Apply the finite tensor-cylinder
identity RRT12: all complementary-prime coordinates sum to their
full identity, leaving exactly the sum of products in (ARO.14)
times \([1]\widehat e\).  Invalid labels contribute zero by ARO6.

Each numerator is real and nonnegative.  Indeed the product of a
central self-adjoint idempotent and a subgroup-character idempotent
is a self-adjoint idempotent; its identity coefficient is
\([1](b^*b)=\sum_\gamma|b(\gamma)|^2\).  This is also the
positivity calculation of RRT2.  The denominators are \(1/c_j>0\).
Finally divide both finite identity coefficients by the same
\(V>0\), using ARO5 and ARO10.  This yields the literal adelic
formula with \(e(1)\), proving (ARO.14).  The choice of finite
characters changes neither the measure nor any projector \(e_T\).
\end{proof}
% END ARO12

\subsection{Operators, multiplicities, and the original good rows}

% BEGIN ARO13
\paragraph{ARO13. All event kernels are admissible orthogonal kernels.}
\label{claim:ARO13}
For all sets \(S,T\) of tuples,
\[
 e_T(S)*e_T(T)=e_T(S\cap T),\quad e_T(S)^*=e_T(S),
 \quad e_T(\varnothing)=0,\quad e_T(\mathrm{all})=e.
 \tag{ARO.15}
\]
Each kernel is locally constant, supported in \(K\), integrable,
and of positive type.  Thus \(\mathrm{EtUniv}\), \(\mathrm{EtInt}\)
and every original \(\mathrm{EtAlg}(S)\) are supplied.
\begin{proof}
By ARO10, \(\widehat e\) is a self-adjoint idempotent.  ARO6
makes every \(E_S\) central, so
\((\widehat eE_S)(\widehat eE_T)=\widehat eE_{S\cap T}\)
and \((\widehat eE_S)^*=\widehat eE_S\).  ARO5 transports these
identities to the kernels.  Their support and local constancy follow
from their finite reduction description; compact support makes
them integrable.  ARO6 gives the empty and universal events.

For positive type let \(h=e_T(S)\), let \(x_1,\ldots,x_m\) be
arbitrary finite-adelic elements, and let \(c_i\in\mathbb C\).
Using \(h*h=h\), \(h(x^{-1})=\overline{h(x)}\), and left Haar
invariance gives
\[
 \sum_{i,j}\overline{c_i}c_j h(x_i^{-1}x_j)
       =\int\left|\sum_j c_jh(z^{-1}x_j)\right|^2d\mu(z)
       \geq0.                                         \tag{ARO.16}
\]
Each product under the finite expansion is integrable, since its
support lies in a translate of a compact set.  Thus all exchanges
of the finite sum and the integral are justified.  This proves
positive type on the full group, not just on \(K\).
\end{proof}
% END ARO13

% BEGIN ARO14
\paragraph{ARO14. Bounded convolution on the given quotient Hilbert space.}
\label{claim:ARO14}
Let \(\mathcal H=X.B.I.L^2(\mu,X.B.\mathrm{fund})\) with its
original right regular action \(R\).  Every integrable finite kernel
\(h\) has a unique bounded operator \(R(h)\), characterized by
\[
 \langle u,R(h)v\rangle
       =\int h(x)\langle u,R(1,x)v\rangle\,d\mu(x),
 \qquad \|R(h)\|\leq\|h\|_1.                       \tag{ARO.17}
\]
Moreover \(R(h_1*h_2)=R(h_1)R(h_2)\) and
\(R(h^*)=R(h)^*\).
\begin{proof}
The local unitary groups are unimodular: the modular character is
the absolute determinant of the adjoint action on the Lie algebra,
and the invariant nondegenerate trace pairing makes this determinant
have absolute value one.  In the split case the same conclusion is
immediate on matrix units; in the nonsplit case the trace pairing
descends from the splitting extension.  Products and restricted
products preserve this assertion.  Thus the given Haar measure is
also right invariant and inversion invariant.

Right translation on the quotient preserves the measure specified
by the original measurable fundamental domain.  Consequently
\(R(ab)=R(a)R(b)\), \(R(1)=1\), and every \(R(a)\) is unitary.
The matrix coefficients are measurable: translate measurable
representatives and use a fundamental-domain representative,
then approximate \(L^2\) vectors by simple functions.  Their
absolute values are at most \(\|u\|\|v\|\).  Hence the right side
of (ARO.17) is a bounded sesquilinear form, conjugate linear in
\(u\), and Riesz representation supplies a bounded operator of
the indicated norm.  The same formula proves uniqueness.

For multiplication, the absolute value of the double integrand is
bounded by \(|h_1(x)h_2(y)|\|u\|\|v\|\), whose integral is
\(\|h_1\|_1\|h_2\|_1\|u\|\|v\|\).  Fubini and the change
of variable \(z=xy\) therefore give the convolution identity.
Taking the adjoint in (ARO.17), using \(R(1,x)^*=R(1,x^{-1})\),
and changing \(x\) to \(x^{-1}\) gives the involution identity.
This proves the original \(\mathrm{IsRConvF}\) specification,
including its existence and uniqueness assertions.
\end{proof}
% END ARO14

% BEGIN ARO15
\paragraph{ARO15. The literal event operators form one finite projection algebra.}
\label{claim:ARO15}
The original \(\mathrm{ReT}(X,S)\) equals \(P_S=R(e_T(S))\).
It satisfies
\[
 P_SP_T=P_{S\cap T},\quad P_S^*=P_S,\quad
 P_{\rm all}=X.B.R_e,
 \quad P_S=\sum_{t\in\mathcal T\cap S}P_{\{t\}}.
 \tag{ARO.18}
\]
Every \(P_S\) commutes with the entire infinity action and with
the original operator \(R_\phi=X.Rph\).
\begin{proof}
ARO13 supplies an integrable kernel, and ARO14 supplies its unique
operator.  Therefore the existential branch in the definition of
\(\mathrm{ReT}\) is realized; its zero fallback introduces no
alternative operator.  ARO13--ARO14 prove the product and adjoint
identities.  The universal kernel is \(e\), so uniqueness identifies
its operator with the given \(X.B.R_e\).  The finite split in ARO6,
linearity of inflation and convolution, and linearity of (ARO.17)
give the sum in (ARO.18).  Distinct singleton projectors therefore
have zero product.

The finite-adelic action \(R(1,x)\) commutes with \(R(y,1)\).
Inserting this equality in (ARO.17) proves that every \(P_S\)
commutes with the infinity action.  The defining weak integral for
\(R_\phi\) then gives \(P_SR_\phi=R_\phi P_S\): move the bounded
operator through the inner product and the scalar integral on
each side.  Thus all event and infinity projections concern the
same quotient Hilbert space and the same regular representation.
\end{proof}
% END ARO15

% BEGIN ARO16
\paragraph{ARO16. The exact infinity projector under the original gap.}
\label{claim:ARO16}
Assume the original \(\mathrm{GapHyp}(D)\).  Let \(\sigma\) be
the original Hilbert tensor product of the local representations,
let \(w=\bigotimes_vw_v\) be its specified unit vector, and put
\(m(y)=\langle w,\sigma(y)w\rangle\).  Then
\(\phi(y)=(\prod_vd_v)\overline{m(y)}\) is integrable for the
original product infinity Haar measure.  The supplied
\(T=R_\phi\) satisfies
\[
 T^*=T,\quad T^2=T,\quad
                 T R(y,1)T=m(y)T.                    \tag{ARO.19}
\]
The span of \(\sigma(y)w\), over all infinity elements \(y\),
is dense in the given tensor Hilbert space.
\begin{proof}
DC14 checks the literal ordered parameters and the original gap;
DC13 identifies every supplied local \(\mathrm{IsDS}\) with the
reconstructed model at that gap.  DC6 gives integrability of its
lowest coefficient, and DC12 gives full four-vector Schur
orthogonality with its actual Haar formal degree.  The quotient
defining the source's \(d_v=\mathrm{fdegOf}\) equals that degree:
set all four vectors equal to the given unit vector in the Schur
identity.  No additional formal-degree formula is needed here.

The product has finitely many places.  Fubini gives integrability
of \(\phi\) and the product Schur identity.  Applying that identity
to two rank-one coefficients gives, in the infinity convolution
algebra,
\(\phi^*=\phi\), \(\phi*\phi=\phi\), and the coefficient identity
whose integrated form is \(R(\phi)R(y)R(\phi)=m(y)R(\phi)\).
For the latter identity expand the two convolutions and use
\[
 \int\langle a,\sigma(x)b\rangle
             \overline{\langle a',\sigma(x)b'\rangle}\,dx
  = (\prod_vd_v)^{-1}\langle b',b\rangle\langle a,a'\rangle.
\]
Products of coefficients are integrable by Cauchy--Schwarz, and
\(\phi\in L^1\) permits the outer weak integrals.  The operator
identities follow as in ARO14.  Uniqueness of the given operator
integral identifies \(R(\phi)\) with the original \(T\).

In each local irreducible unitary representation, the closed span
of the orbit of the nonzero \(w_v\) is a nonzero invariant subspace,
hence the whole representation.  Tensor products of vectors in
these orbit spans are dense in the original tensor product, whose
defining tensor map has dense span.  They are linear combinations
of the global orbit of \(w\).  This proves the asserted cyclicity.
\end{proof}
% END ARO16

% BEGIN ARO17
\paragraph{ARO17. Evaluation identifies multiplicity maps with the actual row space.}
\label{claim:ARO17}
Under the hypotheses of ARO16, let \(Q\) be any of the event
operators of ARO15, or \(X.B.R_e\).  Let
\(\operatorname{Hom}_{G_\infty}(\sigma,\mathcal H)\) mean bounded
intertwiners, and set
\[
 \mathcal M_Q=\{Q\circ A:A\in
                    \operatorname{Hom}_{G_\infty}(\sigma,\mathcal H)\}.
\]
Then evaluation at the original \(w\) is a linear isomorphism
\[
              \mathcal M_Q\ \xrightarrow{\ A\mapsto A(w)\ }
                       \operatorname{im}(QT).        \tag{ARO.20}
\]
\begin{proof}
All members of \(\mathcal M_Q\) are intertwiners by ARO15.
Cyclicity in ARO16 proves injectivity: an intertwiner vanishing on
\(w\) vanishes on its orbit span and then, by continuity, everywhere.
Schur orthogonality in ARO16 says
\(\sigma(\phi)w=w\).  Move a bounded intertwiner through the
weak integral to get \(TA(w)=A(w)\).  Thus
\((Q A)(w)=QT A(w)\), proving inclusion in the stated image.

For the reverse inclusion first let \(z\in\operatorname{im}T\).
The rank-one compression identity (ARO.19) gives
\[
 \langle R(a,1)z,R(b,1)z\rangle
       =\|z\|^2\langle\sigma(a)w,\sigma(b)w\rangle.
 \tag{ARO.21}
\]
Indeed move \(R(a,1)\) to the other side and insert \(Tz=z\)
on both sides of the inner product, then apply (ARO.19) to
\(a^{-1}b\).  For a finite linear combination, (ARO.21) shows that
\[
 \sum_i c_i\sigma(a_i)w\longmapsto
                         \sum_i c_iR(a_i,1)z
\]
is well-defined and has norm \(\|z\|\) times the norm of its
argument.  It extends to a bounded map \(A_z\) on the full Hilbert
space, including the case \(z=0\).  Translation of the dense orbit
span proves equivariance; continuity extends it to all vectors.
One has \(A_z(w)=z\).  Given \(u=QT y\), take \(z=Ty\) and
use \(Q A_z\); its value at \(w\) is \(u\).  This proves
surjectivity and hence (ARO.20).  The proof establishes the exact
row-evaluation statement and does not identify multiplicities with
Hilbert dimensions merely by analogy.
\end{proof}
% END ARO17

% BEGIN ARO18
\paragraph{ARO18. Finite-dimensionality at the precise retained hypothesis.}
\label{claim:ARO18}
Assume \(\mathrm{RowsFinDim}\) and \(\mathrm{GapHyp}(D)\).
Then \(H_X=\operatorname{im}(X.B.R_e\,T)\) and
\(\mathcal M_e\) are finite dimensional.  Every
\(\mathcal M_{P_S}\) is finite dimensional as well.
\begin{proof}
The original \(\mathrm{RowsFinDim}\) quantifies over the given
Haar measure and fundamental domain, the regular representation,
a locally constant compactly supported finite kernel, the local
irreducible unitary representations, their formal degrees, unit
vectors, integrable lowest coefficients, and the specified infinity
operator integral.  The original setting supplies the measure,
fundamental-domain and representation fields; ARO13 supplies the
kernel properties and ARO16 supplies the degree and integrability
fields.  Thus this exact retained assertion applies to \(e\) and
\(T\), giving finite-dimensionality of \(H_X\).
ARO17 transfers it to \(\mathcal M_e\).

By ARO15, \(P_SX.B.R_e=P_S\).  Consequently every
\(P_S A\) is \(P_S(X.B.R_eA)\), and hence
\(\mathcal M_{P_S}\) is an image of the finite-dimensional
\(\mathcal M_e\).  This proves the last assertion for every set
\(S\).  The proof does not claim a new general automorphic
finite-multiplicity theorem: that is exactly the named retained
\(\mathrm{RowsFinDim}\) premise used in its original scope.
\end{proof}
% END ARO18

% BEGIN ARO19
\paragraph{ARO19. Additivity of the original row multiplicities.}
\label{claim:ARO19}
Under the hypotheses of ARO18, the original quantity
\(d(S)=\mathrm{dimX}(X,S)\) satisfies
\[
                  d(S)=\sum_{t\in\mathcal T\cap S}d(\{t\}).
 \tag{ARO.22}
\]
This proves the entire \(\mathrm{RowsAdd}(X)\).
\begin{proof}
By definition \(d(S)=\dim_{\mathbb C}\mathcal M_{P_S}\).
For an intertwiner \(A\), ARO15 decomposes \(P_SA\) as the
finite sum of \(P_{\{t\}}A\) over the indicated labels.  Conversely,
if \(A_t\in\mathcal M_{P_{\{t\}}}\), then
\(P_S A_t=A_t\); hence any finite sum belongs to
\(\mathcal M_{P_S}\).  The sum is direct: applying
\(P_{\{u\}}\) to a vanishing sum singles out its \(u\)-summand,
because the singleton projectors are pairwise orthogonal.
All spaces are finite dimensional by ARO18, so dimensions add.
The finite set \(\mathcal T\) is exactly the source's
\(\mathrm{labT}\), by the coefficient bijection and invalid-label
convention in ARO6.  This proves (ARO.22) for all sets, not only
for cylinders or finite subsets chosen in advance.
\end{proof}
% END ARO19

% BEGIN ARO20
\paragraph{ARO20. Good singleton ranges and their exact dimensions.}
\label{claim:ARO20}
Under the hypotheses of ARO18, let \(G\) be the original good operator,
namely convolution by
\(\operatorname{inf}(E_{\rm good})\), and put
\[
 V_t=H_X\cap\ker(P_{\{t\}}-1),\qquad
 H_{\rm good}=H_X\cap\ker(G-1).
\]
Then \(V_t\perp V_u\) for \(t\ne u\);
\(V_t\subseteq H_{\rm good}\) for every good tuple \(t\); and
\[
                         d(\{t\})=\dim_{\mathbb C}V_t.
 \tag{ARO.23}
\]
\begin{proof}
The good operator exists by ARO4 and ARO14.  Uniqueness identifies
it with any supplied \(\mathrm{GoodOp}\).  On the finite side,
\(E_{\rm good}(\widehat eE_{\{t\}})=\widehat eE_{\{t\}}\)
for a good \(t\).  ARO5 and ARO14 give
\(GP_{\{t\}}=P_{\{t\}}\), proving the good-range inclusion.
If \(v=P_{\{t\}}v\) and \(u=P_{\{u\}}u\) with distinct labels,
self-adjointness and zero product in ARO15 give \(\langle v,u\rangle=0\).

To prove the dimension assertion, identify the spaces exactly.
Since \(P_{\{t\}}\) commutes with \(T\), is idempotent, and is
absorbed by \(X.B.R_e\), its image after \(T\) lies in \(V_t\).
Conversely, if \(v=X.B.R_eTy\) and \(P_{\{t\}}v=v\), then
\(v=P_{\{t\}}Ty\).  Hence
\(\operatorname{im}(P_{\{t\}}T)=V_t\).
ARO17 identifies this image with \(\mathcal M_{P_{\{t\}}}\),
whose dimension is the original \(d(\{t\})\).  This proves (ARO.23)
for every valid tuple; no abstract replacement for the original
multiplicity space has been inserted.
\end{proof}
% END ARO20

% BEGIN ARO21
\paragraph{ARO21. The given good rows have exactly the singleton multiplicities.}
\label{claim:ARO21}
Retain the original isometry
\(\iota:D.s.HH\simeq H_X\), its \(\mathrm{RowDictTie}\) with
the good operator, and its \(\mathrm{LabTie}\).  Let
\(b_\sigma=\iota(D.sEE(D.\mathrm{what}_2.\mathrm{onb}_\sigma))\)
and let \(t_\sigma\) be the original tuple of labels of that row.
For every good valid tuple \(t\),
\[
                   \#\{\sigma:t_\sigma=t\}=d(\{t\}).
 \tag{ARO.24}
\]
Together with the label and goodness clauses this is
\(\mathrm{RowsGood}(X)\).
\begin{proof}
The literal row-dictionary tie is an isometric linear isomorphism
\(j\) onto the conjugate good space, followed by complex
conjugation.  Thus \(b_\sigma=\overline{j(\mathrm{onb}_\sigma)}\).
These vectors are linearly independent: conjugate a vanishing
linear relation, use independence of the original orthonormal
basis and injectivity of \(j\), then conjugate the coefficients
back.  They span \(H_{\rm good}\) for the same reason, using
surjectivity onto the conjugate good space.  This argument respects
the semilinear map; it does not treat conjugation as a complex-linear
identification.

The original \(\mathrm{LabTie}\) states that each \(t_\sigma\)
is valid and good and that \(b_\sigma\in V_{t_\sigma}\).
For a good \(t\), ARO20 puts \(V_t\) inside the span of all these
vectors.  If \(v\in V_t\), express it in this finite basis and
split the sum into labels equal to \(t\) and unequal to \(t\).
The first sum belongs to \(V_t\); subtracting it shows that the
second also belongs to \(V_t\).  By ARO20 every summand of the
second is orthogonal to \(V_t\), so its norm is zero.  It vanishes.
Thus the vectors with label \(t\) form a basis of \(V_t\).
ARO20 now gives (ARO.24).  Repeated labels are allowed: their
number is exactly the corresponding dimension.
\end{proof}
% END ARO21

% BEGIN ARO22
\paragraph{ARO22. Completing the label list while preserving every given good row.}
\label{claim:ARO22}
Let a finite label set \(\mathcal T\) have multiplicities
\(m_t\in\mathbb N\), a good-label predicate, and a given finite
row set \(I\) with good labels \(t_i\).  Suppose each good label
occurs in \(I\) exactly \(m_t\) times.  If
\(M=\sum_tm_t\), then there are labels on \(\{0,\ldots,M-1\}\)
and an injection from \(I\) into that set such that the original
labels are preserved, the good indices are exactly the image, and
every event \(S\) has exactly \(\sum_{t\in S}m_t\) indices.
\begin{proof}
Form the finite disjoint union
\[
                \mathcal R=I\ \sqcup\!
                  \coprod_{t\text{ not good}}\{t\}\times\{0,\ldots,m_t-1\}.
 \tag{ARO.25}
\]
Label an original row by its given label and a new pair by \(t\).
The multiplicity assumption on good labels shows
\(|\mathcal R|=\sum_tm_t=M\).  Choose any bijection between
\(\mathcal R\) and the standard \(M\)-element set.  Transport the
labels and the inclusion of the first summand through this one
bijection.  Its restriction to \(I\) is injective and preserves
labels; every added row has a nongood label and every original
row has a good label.  Counting an event in the two disjoint
summands gives the asserted formula.  This construction also works
when a multiplicity, \(I\), or the entire row space is zero.
It adds precisely the missing rows and makes no separate choice
of an enumeration for different events.
\end{proof}
% END ARO22

% BEGIN ARO23
\paragraph{ARO23. All fields of the original row-label record.}
\label{claim:ARO23}
Assume the hypotheses of ARO18 and ARO21 and the original
\(\mathrm{AutTie}(D,X.B,X.pl)\).  Then there is one label map
\(\mathrm{labAll}\) and one injection
\(\mathrm{inc}:\mathrm{Fin}(D_{\rm rows})\hookrightarrow
\mathrm{Fin}(D_{\rm all})\) such that
\[
 d(S)=\#\{\tau:\mathrm{labAll}(\tau)\in S\},
 \quad \mathrm{labAll}(\mathrm{inc}(\sigma))=t_\sigma,
 \quad \mathrm{labAll}(\tau)\text{ good}\ \Longleftrightarrow\
                           \tau\in\operatorname{im}(\mathrm{inc}).
 \tag{ARO.26}
\]
Every label is valid.  These are all fields of the original
\(\mathrm{RowsLab}(X)\).
\begin{proof}
By ARO15 the universal event operator is the original \(R_e\),
so \(d(\mathrm{all})=X.B.D\).  The dimension field of the given
\(\mathrm{AutTie}\) identifies this number with the original
\(D_{\rm all}\).  ARO19 therefore gives
\(D_{\rm all}=\sum_{t\in\mathcal T}d(\{t\})\).
ARO21 identifies every good-label fiber of the original
\(D_{\rm rows}\) rows with its required multiplicity.
Apply ARO22 with \(m_t=d(\{t\})\) and this exact row set.
Its event-count formula and ARO19 give the first identity in
(ARO.26); its inclusion and goodness statements give the other
two.  By construction labels belong to \(\mathcal T\), so they
satisfy every original \(\mathrm{IsLabel}\) clause.  The maps are
chosen once before \(S\), \(j\), or a cylinder is specified.
\end{proof}
% END ARO23

% BEGIN ARO24
\paragraph{ARO24. One actual AllRows object with all label and bundle ties.}
\label{claim:ARO24}
Under the hypotheses of ARO23 there is an original
\(A:\mathrm{AllRows}(D)\) with
\(\mathrm{LabelCountGoodInjective}(A)\) and
\(\mathrm{AllJoinA}(A)\).  It has \(A.B=X.B\), the event kernels
and operators of ARO12--ARO15, and
\[
                 A.eVec(j,v)=e_J(q_j,n,k_0,s,\psi_j^0)
                    \quad\hbox{for every }j,v,        \tag{ARO.27}
\]
where one nontrivial \(\psi_j^0\) is fixed for each prime.
\begin{proof}
Take, for example, \(\psi_j^0(a)=\exp(2\pi i\widetilde a/q_j)\).
This is nontrivial since \(q_j>1\).  The local idempotence in
(ARO.27) is the normalized subgroup-character identity of RRT1.
Use the label map and injection of ARO23, and use the given
\(X.B\).  The dimension identities \(A.B.I.n=n\) and
\(A.B.I.f=f\) are the original rank field and the degree equality
under \(X.\iota_E\).  ARO14--ARO15 give every operator specification
and the universal operator; ARO13 gives the universal kernel;
ARO12 gives the full cylinder identity.  These are every field
of the \(\mathrm{AllRows}\) structure.  ARO23 supplies all four
clauses of \(\mathrm{LabelCountGoodInjective}\).

For \(\mathrm{AllJoinA}\), its level clause is (ARO.27) with the
chosen nontrivial characters.  Its bundle clause is the original
\(\mathrm{RowsTie}\): use \(X.pl\) and the given \(\mathrm{AutTie}\),
the original Harish--Chandra parameter identities in that tie, and
complete splitting of \(q_0=X.B.g_0\) transported by
\(X.\iota_E\).  Complete splitting follows from the original
full list of residue-degree-one places; the sum of residue degrees
is the field degree.  Hence no extra ramification can occur.
All these fields refer to the same original bundle.  In particular
the row labels have not been supplied by an unrelated finite model.
\end{proof}
% END ARO24

\subsection{The count-idempotent interface and its common-family scope}

% BEGIN ARO25
\paragraph{ARO25. Partner-prime unipotence and the full admissibility fields.}
\label{claim:ARO25}
For every \(j\in\{0,1\}\) and every set \(S\) of its tuples,
\(e_T(\operatorname{cyl}(j,S))\) is an original admissible
idempotent at level \(g\), with distinguished unipotence prime
\(q_{1-j}\).
\begin{proof}
Decompose \(\mathcal G=\mathcal G_{1-j}\times\mathcal G_j\).
ARO6 shows that the cylinder projector is \(1\otimes Z_S\),
where \(Z_S\) is the central event projector in the \(j\)-factor.
ARO10 gives \(\widehat e=a_{1-j}\otimes a_j\), where each
\(a_i\) is the product of the corresponding conjugated character
idempotents.  Thus
\[
       \widehat{e_T(\operatorname{cyl}(j,S))}
                         =a_{1-j}\otimes(a_jZ_S).     \tag{ARO.28}
\]
If a coefficient is nonzero, its \((1-j)\)-part lies in the
product of the conjugated \(J\)'s.  Its reduction there is
unipotent.  The conjugate reductions are inverse transposes and
also unipotent.  Thus \((x-1)^n\widehat L\subseteq
q_{1-j}\widehat L\), which is the original unipotence-modulo-prime
condition.  The same conclusion follows by lifting that one
factor to \(X.K_2\) using ARO2 and ARO8 and comparing modulo
\(K(q_{1-j})\) by ARO1.

ARO13 supplies local constancy, convolution idempotence,
self-adjointness, positive type and support in \(K\).
ARO4--ARO5 supply invariance under multiplication on both sides
by \(K(g)\).  Together with the preceding unipotence these are
every field of the original \(\mathrm{IsAdmIdem}\).  They also
hold when the cylinder projector is zero.  The prime used for
unipotence is deliberately the complementary one, exactly as in
the original count application.
\end{proof}
% END ARO25

% BEGIN ARO26
\paragraph{ARO26. Every cylinder is an original counting idempotent.}
\label{claim:ARO26}
For every \(j,S\), the kernel of ARO25 satisfies the original
\(\mathrm{IsCountIdem}(\mu,q_{1-j},e_T(\operatorname{cyl}(j,S)))\).
Consequently the common object of ARO24 has the full
\(\mathrm{PCountTie}\), for every set \(S\).
\begin{proof}
Use exponents \(a=b=1\), second prime \(p=q_j\), and the
reduction datum
\[
 K\longrightarrow\mathcal G_{1-j}\times\mathcal G_j
\]
from ARO1--ARO2.  It is multiplicative and surjective; equality
of first, respectively second, coordinates is equivalent to
membership of the quotient \(x^{-1}y\) in \(K(q_{1-j})\),
respectively \(K(q_j)\).  Extending the reduction map by the
identity off \(K\) gives exactly the original total-map convention.

Set \(f_0=a_{1-j}\), \(\varepsilon_2=a_j\), and \(z=Z_S\)
in (ARO.28).  Each \(a_i\) is the character idempotent of the
product subgroup \(\prod_vP_{i,v}J_{q_i}P_{i,v}^{-1}\): its
character is the product of the local characters in ARO9, is
multiplicative there, and has modulus one.  Its coefficients are
the inverse character divided by the order of that subgroup.
Thus both \(f_0\) and \(\varepsilon_2\) satisfy all the original
\(\mathrm{IsCharIdemFin}\) clauses.  ARO6 makes \(z\) a central
idempotent.  ARO25 proves the required unipotence when the first
coefficient is nonzero.

Finally ARO1 and elementary coprime-lattice congruence give
\(K(q_0)\cap K(q_1)=K(g)\).  Equations (ARO.5) and (ARO.28)
therefore give exactly the required formula
\[
 e_T(\operatorname{cyl}(j,S))(x)=
 \begin{cases}
 \displaystyle\frac{f_0(\rho_{1-j}(x))
                 (z\varepsilon_2)(\rho_j(x))}
                 {\mu(K(q_{1-j})\cap K(q_j))},&x\in K,\\[6pt]
 0,&x\notin K.
 \end{cases}                                         \tag{ARO.29}
\]
The denominator is the original \(V\), not either one-prime
volume.  The two primes are distinct by the datum.  This supplies
every existential witness and every clause of
\(\mathrm{IsCountIdem}\).  The right disjunct of the literal
\(\mathrm{PCount}\) now gives \(\mathrm{PCountTie}\).
Arbitrary \(S\) is permitted because ARO6 reduces it to its
finite set of valid labels.
\end{proof}
% END ARO26

% BEGIN ARO27
\paragraph{ARO27. The exact common count parameters on the original bundle.}
\label{claim:ARO27}
Retain \(\mathrm{AutTie}\), an original size datum \(L\), the
original \(\mathrm{CountTheoremHypotheses}(D,L,n_0,A,C,c_b)\),
and \(4n<q_1\).  Then the full original
\(\mathrm{CountBase}(X,n_0,A,C,c_b)\) holds with
\[
 b=L.sc.b,\qquad h=Cn^{19}f^5\log(b\,\mathrm{rd}_E).
 \tag{ARO.30}
\]
The same two numbers work for both auxiliary primes and every
cylinder.
\begin{proof}
The retained count hypotheses assert the original discriminant
identity, \(n_0\leq n\), \((s+1)^2\leq n\),
\(\log\log\mathrm{rd}_E\leq n\),
\(|\lambda_{\tau,a}|\leq b^2\), \(g\leq b^A\),
\(4n<q_0\), complete splitting of \(q_0\), the gap at least
the value \(h\) in (ARO.30), and
\[
 \log f\leq\sqrt n,\qquad 2\leq b,\qquad
              \log\log b\leq c_bnf\log n.            \tag{ARO.31}
\]
Their discriminant clause is retained although this particular
transport does not need to re-use it.  The given \(\mathrm{AutTie}\)
identifies the bundle's rank, real degree, root discriminant,
ordered parameters and signs with these exact quantities.  In
particular, choose \(p\) and \(q\) to be the numbers of true and
false entries of the original \(\mathrm{epsS}\).  The
\(\mathrm{eps\_\allowbreak eq}\) field gives the same
\(p\mathrm{At}(v)=p\) and \(q\mathrm{At}(v)=q\) at every
infinity place, as required by \(\mathrm{SetBase}\).  The
number of negative signs is at most \(s+1\), so the signature
square bound follows from \((s+1)^2\leq n\).  The original
Harish--Chandra parameter conditions transport with the ordered
parameter; no reordering is made.

The group level is the given \(g=q_0q_1\geq1\).  Complete
splitting at both primes follows from the original datum as in
ARO24, and the two inequalities \(4n<q_0,q_1\) provide the
distinguished-prime guard for either choice of \(j\).
The gap, its lower bound by \(Cn^{19}f^5\log(b\mathrm{rd}_E)\),
and (ARO.31) then hold on the original bundle.  Finally the
\(D_*\)-field of \(\mathrm{AutTie}\) identifies
\(X.B.I.D_*(e)\) with the original positive \(D_*\).
These are all the clauses of \(\mathrm{CountBase}\), with the
single pair \((b,h)\) of (ARO.30).
\end{proof}
% END ARO27

% BEGIN ARO28
\paragraph{ARO28. Installation of the full cylinder count from its exact theorem.}
\label{claim:ARO28}
In addition to the hypotheses of ARO24 and ARO27, supply the
original assertion
\(\mathrm{CountIdempotentsClass}(\mathrm{PCount},D,n_0,A,C,c_3,c_b)\).
For any \(0<c\leq c_3\), the same \(A\) of ARO24 satisfies
\[
        \mathrm{CountAllRows}\bigl(A,\exp(-cnf\log n)\bigr).
 \tag{ARO.32}
\]
This is conditional on that exact arithmetic counting theorem;
the theorem itself is not declared proved here.
\begin{proof}
For each \(j,S\), ARO25 supplies its actual admissible kernel,
ARO26 supplies the \(\mathrm{PCount}\) predicate, and ARO27
supplies the same parameters \(b,h\).  The discrete-series,
unit-vector, Hilbert tensor, fundamental-domain, and regular-action
fields are those of \(X\), with the operator specification of
ARO14--ARO15.  Thus every argument of the supplied counting
theorem is available on this very bundle.  Its conclusion is
\[
 \left|d(\operatorname{cyl}(j,S))-
                D_*^A(e_T(\operatorname{cyl}(j,S)))\right|
 \leq e^{-c_3nf\log n}\,
                D_*^A(e_T(\operatorname{cyl}(j,S))).   \tag{ARO.33}
\]
This proves the original \(\mathrm{CountCylAt}\) on all cylinders.
Here \(n\geq2\) follows from the original gap hypothesis and
\(f\geq1\) from the degree.  Therefore
\(0<e^{-cnf\log n}<1\).  The left side of (ARO.33) is
nonnegative and its exponential factor is positive, so its right
normalizing \(D_*^A\) is nonnegative, including for zero kernels.
Since \(c\leq c_3\), replacing the exponential factor by
\(e^{-cnf\log n}\) preserves the inequality.  ARO27 supplies
\(D_*^A(e)>0\).  These are precisely all four clauses of
\(\mathrm{CountAllRows}\).  Neither a new row object nor a
different measure or formal degree is introduced in this step.
\end{proof}
% END ARO28

% BEGIN ARO29
\paragraph{ARO29. The completed local weights apply to this same row object.}
\label{claim:ARO29}
Retain the original normalized-type compatibility, admissibility
\(\mathrm{TieAdmG}\), and range \(\mathrm{RangeW}\), together
with the actual-rank threshold of FLC19.  Then the object of ARO24
satisfies the full \(\mathrm{WeightAllRows}\), with exactly the
original \(W_g\).  Its universal and two finite bad-cylinder
normalizations are
\[
 D_*^A(e)=M>0,\qquad
 D_*^A(e_T(\operatorname{cyl}(j,S_j)))=Mp_j,
 \quad p_j=\operatorname{locShare}(A,j,S_j)\geq0.      \tag{ARO.34}
\]
\begin{proof}
ARO24 provides the original \(\mathrm{AllJoinA}\), with its
precise \(e_J\) vectors and one chosen nontrivial character at
each prime.  Thus FLC19 applies with all of its original finite
type, range and admissibility inputs.  It supplies the unchanged
\(W_g\) estimate and nonnegativity for every local label and
every finite bad event.  In the selected eventual family its
additional absolute rank threshold is absorbed exactly as in FLC19,
using the original \(x^2\leq f\leq n\) growth.

For \(S_j\), take the finite image of the bad labels among the
all-row labels of ARO23.  Membership in the cylinder is exactly
badness of the row at prime \(j\); multiplicities are counted by
(ARO.26).  This is the finite-event construction in GQT2.
ARO12 gives \(e_T(\operatorname{cyl}(j,S_j))(1)=e(1)p_j\),
with real values and nonnegative \(p_j\).  By definition \(D_*^A(h)\)
is the same product of formal degree and quotient volume times
\(\operatorname{Re}h(1)\), for every \(h\).  Hence it gives
(ARO.34).  Its positivity for the full event is the original
\(\mathrm{AutTie}\) field used in ARO27.  In particular the
finite weight theorem and the arithmetic count theorem now refer
to exactly the same operators and labels.
\end{proof}
% END ARO29

% BEGIN ARO30
\paragraph{ARO30. Exact common-family consumer and the remaining existence inputs.}
\label{claim:ARO30}
Fix the original family \(\mathcal F\), its \(0<\epsilon\leq1\),
\(\epsilon_z=\epsilon/J\) with \(J\geq1\), and one prior
\(c>0\).  Suppose one threshold precedes every varying member
\(D=\mathcal F.D(x,i)\), and that each such member supplies one
\(X,G,\iota\) with the original \(\mathrm{AutTie}\),
\(\mathrm{RowDictTie}\) and \(\mathrm{LabTie}\), the gap and
count-parameter inputs of ARO27, the exact counting theorem of
ARO28 with \(c\leq c_3\), and the finite local hypotheses of
ARO29.  Retain the original \(\mathrm{RowsFinDim}\) and
\(\mathrm{SizesRec}(x,\epsilon_z,C_s)\).
Then the literal \(\mathrm{LinkIRR}(\mathcal F)\) follows.
\begin{proof}
At each member apply ARO24 once to its given \(X,G,\iota\).
ARO28 gives \(\mathrm{CountAllRows}\) for this object, and ARO29
gives \(\mathrm{WeightAllRows}\) for the same object.  ARO23
gives the exact original label injection and good-image condition.
Thus all inputs of GQT5 are supplied, with the original dimension
\(D_{\rm all}\), retained dimension \(D_{\rm rows}\), and stored
defect \(D.\mathrm{what}_2.\mathrm{delp}\).
Take the maximum of the common input threshold, the actual-rank
threshold in ARO29, and the numerical thresholds already proved
in GQT5.  None depends on the varying index \(i\).  GQT5 then
gives, for every later member, the original nonnegative defect and
its bound \(\exp(-x^{\epsilon_z/7})\), which are exactly
\(\mathrm{LinkIRR}\).

The same reasoning only needs the three arithmetic inequalities
for the full set and the two finite bad cylinders in GQT2,
provided their relative error is uniformly at most \(1/2\).
Their exact normalizing terms are now proved in ARO29.  This
weaker alternative does not assert that those three estimates
hold without an arithmetic count proof.
\end{proof}
% END ARO30

\paragraph{Scope of the completed producer.}
The original reduction, Haar normalization, level-character shape,
\(\mathrm{EtOne}\), event-operator algebra, \(\mathrm{RowsAdd}\),
\(\mathrm{RowsGood}\), \(\mathrm{RowsLab}\), common
\(\mathrm{AllRows}\) constructor and \(\mathrm{PCountTie}\) have
human proofs above.  The full pointwise count installation is proved
with the original counting theorem as an explicit antecedent.
The remaining global work includes producing the original common
\(X,G,\iota\) and all three ties in the selected family, and
supplying that arithmetic counting theorem or its three sufficient
special cases.  Original \(\mathrm{RowsFinDim}\) remains explicit.
Consequently this chapter does not by itself close the selected
\(\mathrm{GLOBAL\_\allowbreak RETAINED\_\allowbreak DEFECT}\) obligation.  Its contribution
is that these remaining inputs now feed a completely specified
mathematical row construction, with every normalization and
multiplicity retained.

\paragraph{Primary-source map.}
The attached metadata binds every claim to exact bytes of commit
\texttt{7cbf7a0c32fca1badf05931879a37c3f8a2b03b4}.
The principal source contracts are \path{ArithmeticCylinders.lean}
(\(\mathrm{AllRows}\)), \path{AllRowsOf.lean},
\path{AllRowsRhoX.lean}, \path{RhoSurjPf.lean},
\path{KppFwdPf.lean}, \path{AllRowsEtOne.lean},
\path{AllRowsAdd.lean}, \path{AllRowsDcount.lean},
\path{RowsGoodRed.lean}, \path{AllRowsLab.lean},
\path{AllRowsPCount.lean}, and \path{AllRowsAtRecord.lean}.
The direct reverse flag argument in ARO8 and the elementary trace
test in ARO9 simplify intermediate source bookkeeping; their full
original conclusions and original data are unchanged.


\section{Integral cusp lattices on the original geometric family}
\label{sec:icl-v46}

This chapter concerns the lattice defined by \emph{all} the original
finite-place integral-frame tests.  Its first part proves finite generation
by constructing separating evaluations on the actual moduli space.  Its
second part isolates the analytic and integral-theta inputs needed to prove
that this lattice spans, and proves their exact consequence for
\texttt{Lat\allowbreak Full}.  The distinction matters: the latter inputs are original
geometric producer statements, and are not consequences of a dimension
count or of the finite rational-theta implications in RTG.

Fix the source data \(D\), \(T:\mathrm{Core}(D)\), its original
\(Q,X_j,c,l\), a row \(j\in\mathrm{InR}(D)\), and a dominant weight
\(w\).  Put
\[
 K=c.F,\qquad E=D.\mathrm{host}.E,\qquad
 C_w=T.\mathrm{cuspA}(j,\mathrm{wtA}(w),w.n),\qquad
 \mathcal R=T.R(j).
\]
All scalar extensions below use the given inclusion of \(K\).  Auxiliary
number fields used to evaluate a section do not replace \(c\), the common
field \(D.\mathrm{what}_2.LL\), its fixed complex embedding, the chosen
Hermitian forms, the metric, the measure, or any tuple phase.  In particular
the last transport uses the original \(\mathrm{FWTie}\) and its map.

We use the following classical foundations with their stated scope.
The source interfaces and the exact places where they enter are recorded
in the accompanying foundation manifest.

\begin{enumerate}
\item Riemann's algebraic uniformization of polarized complex tori, the
algebraicity of complex abelian varieties with complex multiplication,
and their potential good reduction at every finite place.  For the last
assertion a primary reference is Serre--Tate, \emph{Good reduction of
abelian varieties}, Annals of Mathematics 88 (1968), Theorem 6(a).
\item The extension property for homomorphisms of abelian schemes over
a normal Dedekind base, the theorem of the cube for rigidified line
bundles, and cohomology and base change for an ample principal line
bundle on an abelian scheme.  These are used for the given polarization
and line bundle, including at primes ramified in \(E\).
\item The finite-dimensionality of holomorphic cuspidal automorphic
forms of a fixed highest weight on a fixed arithmetic quotient.  It is
used together with the original algebraic uniformization, not with a
newly selected space of automorphic functions.
\end{enumerate}

The theta-kernel and ramified jet inputs introduced later are kept
separate from these classical foundations.  In particular, this chapter
does not relabel the original kernel identity, chart bounds, or short
series gluing as published theorems.

% BEGIN ICL1
\paragraph{ICL1. The literal valuation lattice.}
\label{claim:ICL1}
For a finite prime \(\mathfrak p\) of \(K\), the source set
\(\mathrm{valSet}(X_j,c,\mathrm{algCls})(\mathfrak p)\) consists of the maps
\[
 f\longmapsto e^{-1}\operatorname{ord}_{V}
                 (\operatorname{ev}_{I}(y,f)),                 \tag{ICL.1}
\]
with value \(+\infty\) at zero.  Here \(V\) is a complete discrete
valuation ring, its residue field is finite, \(\iota:K\hookrightarrow
\operatorname{Frac}(V)\) has ramification index \(e>0\) over
\(\mathfrak p\), and \(y\) is an original \(\mathrm{IntFr}\).
An \(\mathrm{IntFr}\) includes an algebraic test presentation, an integral
semiabelian model, bases of each integral Lie eigenspace, and a generator
of its integral theta-section line with the source rigidified
identifications.  Thus the source \(\mathrm{supLat}(T,j,w)\) is exactly
\[
 L_w=\{f\in C_w:\ v(f)\geq0\text{ for every }\mathfrak p
           \text{ and every }v\in\mathrm{valSet}(\mathfrak p)\}. \tag{ICL.2}
\]
It is an \(\mathcal O_K\)-module.  Its finite generation and the equality
\(K L_w=C_w\) are precisely \(\mathrm{FullAt}(T,j,w)\).
\begin{proof}
The definition \(\mathrm{IntFr.IsAlg}\) is the finite-residue-field
condition; it imposes no fixed field of definition on the test point.
The definition of \(\mathrm{PtV}\) supplies \(e,\iota,y\), and
\(\mathrm{vAtV}\) is (ICL.1).  Conversely every such tuple is a
\(\mathrm{PtV}\) with its unique auxiliary \(\mathrm{Unit}\)-label,
so it occurs in \(\mathrm{valSet}\).  Evaluation is additive,
\(K\)-linear through \(\iota\), and multiplicative.  Consequently
\(v(f+g)\geq\min(v(f),v(g))\),
\(v(af)=\operatorname{ord}_{\mathfrak p}(a)+v(f)\) for \(a\ne0\),
and \(v(fg)=v(f)+v(g)\).  Since the prime valuations of an element
of \(\mathcal O_K\) are nonnegative, (ICL.2) is a module.  The
source \(\mathrm{latIn}\) is the image of the corresponding submodule
of \(C_w\) under its inclusion in \(\mathcal R\); that inclusion is
injective.  This proves both the set identification and the last
assertion, without assuming that the module already spans.
\end{proof}
% END ICL1

% BEGIN ICL2
\paragraph{ICL2. Finitely many integral evaluations suffice.}
\label{claim:ICL2}
Let \(C\) be a finite-dimensional \(K\)-vector space and
\(L\subseteq C\) an \(\mathcal O_K\)-submodule.  Suppose that for
every nonzero \(f\in C\) there are a finite number-field extension
\(K_f/K\) and a \(K\)-linear map \(\varphi_f:C\to K_f\) such that
\[
 \varphi_f(f)\ne0,\qquad \varphi_f(L)\subseteq\mathcal O_{K_f}.
                                                               \tag{ICL.3}
\]
Then \(L\) is finitely generated.  Neither saturation nor a spanning
assumption on \(L\) is needed.
\begin{proof}
Start with \(W_0=C\).  If \(W_i\ne0\), choose a nonzero
\(f_i\in W_i\), choose (ICL.3), and set
\(W_{i+1}=W_i\cap\ker\varphi_{f_i}\).  Its dimension is strictly
smaller.  In at most \(\dim_K C\) steps the common kernel is zero.
The resulting diagonal map embeds \(L\) into
\(\prod_i\mathcal O_{K_{f_i}}\).  Each factor is a finite
\(\mathcal O_K\)-module, and \(\mathcal O_K\) is Noetherian, so
its submodule \(L\) is finite.  If \(C=0\), use the empty product
and \(L=0\).  The argument concerns the whole module \(L\), not a
chosen finite-index submodule.
\end{proof}
% END ICL2

% BEGIN ICL3
\paragraph{ICL3. Rational Hermitian flags detect a nonzero section.}
\label{claim:ICL3}
Let \(X\) be one original Hermitian uniformization component, with
Hermitian space \(E^r\), CM type \(\Psi\), and signature
\((p_0,r-p_0)\) at every \(\tau\in\Psi\).  Rational positive
\(p_0\)-planes and rational orthogonal complementary negative planes
are dense in its period domain.  More precisely, if a continuous
section \(P\) on the original frame domain is Borel equivariant
with a nonzero character and is nonzero somewhere, there is a frame
\(x\) with \(P(x)\ne0\) whose positive and negative subspaces
come from orthogonal \(E\)-subspaces.  The spanning vectors can all
be taken in any prescribed full \(\mathcal O_E\)-lattice.
\begin{proof}
Choose a frame where \(P\ne0\).  Hermitian Gram--Schmidt, in the
order with the positive columns first, changes it by an invertible
Borel matrix.  The defining nonzero principal minors ensure that
each division is legal, and Borel equivariance preserves
nonvanishing.  It is therefore enough to approximate a frame with
orthogonal columns and the required signs.

The real-linear Minkowski map
\(E\otimes_{\mathbb Q}\mathbb R\to\prod_{\tau\in\Psi}\mathbb C\)
is an isomorphism.  Its image of \(E\) is dense: in a
\(\mathbb Q\)-basis this is the density of \(\mathbb Q^d\) in
\(\mathbb R^d\).  Approximate all columns at all embeddings by a
single matrix over \(E\).  Apply Gram--Schmidt over \(E\), using
the original Hermitian form.  The coefficients remain in \(E\),
since the scalar products and nonzero self-pairings lie there.
At all the finitely many embeddings this operation is continuous
near the chosen orthogonal frame.  A sufficiently close
approximation is still a basis, has the same signs, and lies in
the open nonvanishing set of \(P\).

Multiplying the finitely many resulting vectors by a common
positive integer puts them in the prescribed full lattice.  Keep
the approximating complex frame before this multiplication; its
subspaces are unchanged.  If the source
\(\mathrm{LamBasis}(X,\mathcal L,x,\lambda)\) notation is desired,
choose \(\mathcal L>0\) larger than the square roots of all the
finitely many absolute self-pairings.  All its qualitative
conditions hold.  No quantitative bound for \(\mathcal L\), no
maximum of \(P\), and no uniform lattice-basis radius is asserted.
The cases \(p_0=0\) and \(p_0=r\) use the same construction with
one of the two subspaces zero.
\end{proof}
% END ICL3

% BEGIN ICL4
\paragraph{ICL4. These flags give actual algebraic CM points.}
\label{claim:ICL4}
Under the original moduli uniformization, a frame supplied by
ICL3 lies over a CM point of the same \(X_j\).  Its underlying
polarized abelian variety, \(\mathcal O_E\)-action, level marking,
rigidified principal line bundle and any chosen algebraic frame
are defined over some finite extension of \(K\).  Every nonzero
original section has a nonzero restriction to the frame fibre
over at least one such point.
\begin{proof}
The rational orthogonal decomposition gives \(r\) one-dimensional
\(E\)-Hodge subspaces.  The lattice of the uniformized torus is
commensurable with the sum of its intersections with these
subspaces: clear denominators in the finitely many rational
projection matrices.  Thus the torus is isogenous to a product
of the corresponding CM tori.  Equivalently the product algebra
\(E^r\), of dimension twice the abelian dimension over
\(\mathbb Q\), acts on its rational Hodge structure.  Riemann
uniformization turns these rational Hodge endomorphisms into
algebraic endomorphisms.  The variety is therefore of CM type.
CM algebraicity gives a model over a number field.  Enlarge that
field to contain the fixed copy of \(K\) and to define the
finitely many action maps, polarization, line-bundle data and
level-torsion coordinates.  These are algebraic data of finite
presentation.  Their relations hold after complex base change
and hence over this field.  The moduli point is in the original
component with the original marking; an isogenous object is
used to prove its CM property, not substituted for it.

A natural algebraic rule restricted to a fixed frame fibre is a
regular function on
\(\prod_\tau\mathrm{GL}_{m_\tau}\times\mathbb G_m\).
Indeed, apply the rule to the universal matrices and the
universal invertible section scalar.  Their coordinate ring is
a finitely generated reduced \(K\)-algebra, and therefore is
one of the source test rings after a presentation.  Naturality
identifies every later evaluation with evaluation of this one
regular function.  This also proves compatibility with the
source \(\mathrm{FrChg}((Y_\tau)^t,s)\) convention.

Finally, a nonzero rule in the original section space has a
nonzero complex reading.  To check this assertion directly,
its values on a reduced finite-type test algebra are detected
by its complex points; uniformization and the nonzero Borel
weight express each such value as a scalar times a frame
reading.  Vanishing of all readings would make the rule zero
on every test algebra.  Apply ICL3 to a nonzero reading.  At its
CM point the regular function on the frame fibre is nonzero,
since its value at that complex frame is nonzero.
\end{proof}
% END ICL4

% BEGIN ICL5
\paragraph{ICL5. One integral model after a finite extension.}
\label{claim:ICL5}
Let \(x\) be a CM point from ICL4.  After a finite extension
\(K'/K\), the actual abelian variety over \(x\) extends to an
abelian scheme over \(\mathcal O_{K'}\), with its original action
and rigidified symmetric principal line bundle.  Its theta
section module and each integral Lie eigenspace are finite
projective modules of the expected ranks.  They commute with
further flat extensions of the base.
\begin{proof}
There are only finitely many places of bad reduction.  Potential
good reduction for a CM abelian variety gives a finite local
extension at each.  A finite global extension realizing these
local extensions can be chosen by weak approximation and
Krasner's lemma.  Take a normal closure to handle all places
above each chosen prime.  Good reduction is preserved by field
extension, so the resulting variety has good reduction
everywhere.  Its local abelian models glue by the extension
and uniqueness property of abelian schemes.  The finitely
many original action maps and their relations extend by the
same property.

For clarity, no change of the generic theta line is needed.
On the regular abelian scheme, close a divisor representing
that line on the generic fibre.  Regularity makes the closure
Cartier.  Tensor with the inverse of its restriction at the
identity to rigidify it.  The theorem of the cube extends its
generic cubical structure uniquely; a vertical ambiguity is
pulled back from the base and is removed by rigidification.
The generic polarization and its inverse extend, so the
resulting principal polarization is the same one on every
fibre.  The associated line bundle is relatively ample, by
the abelian-scheme polarization criterion.  Symmetry and all
rigidified generic identifications extend uniquely.  One can
equivalently use the original cubical principal-triple
extension theorem; its hypotheses are satisfied by this
everywhere-good model.

For an ample principal line bundle on an abelian variety,
the higher cohomology vanishes and the space of sections has
dimension one.  Cohomology and base change therefore makes
the relative section module an invertible
\(\mathcal O_{K'}\)-module, compatible with flat base change.
Let \(M\) be the Lie module.  After making all labels
\(\tau:E\hookrightarrow K'\) available, its \(\tau\)-eigenspace
is the kernel of the finite matrix
\[
 M\longrightarrow\bigoplus_a M,\qquad
 v\longmapsto (av-\tau(a)v)_a,
\]
where \(a\) ranges over integral generators of
\(\mathcal O_E\).  This kernel is saturated in the torsion-free
module \(M\), hence is finite projective over the Dedekind
ring.  Its generic rank is the original \(m_\tau=\mathrm{Dw.mult}(\tau)\).
Flatness commutes with the displayed kernel, so this is also
the whole integral eigenspace after localization, completion,
or further number-field extension.  We have not asserted
that the sum of these eigenspaces is the full integral Lie
module at ramified primes; that stronger statement is not
needed and is generally unavailable.
\end{proof}
% END ICL5

% BEGIN ICL6
\paragraph{ICL6. Simultaneous integral frames and their density.}
\label{claim:ICL6}
After one further finite extension, all the projective
modules in ICL5 admit bases.  Starting with these bases and a
theta generator, the union of their integral frame changes
over finite number-field extensions is Zariski dense in the
original generic frame fibre.
\begin{proof}
Over a Dedekind ring a projective module of positive rank is
the sum of free modules and one fractional ideal.  For each
of the finitely many ideals \(I\), choose \(h>0\) and \(a\ne0\)
with \(I^h=(a)\), and adjoin a root \(\alpha^h=a\).  At every
prime of the new number field the valuations of
\(I\mathcal O\) and \((\alpha)\) agree, so the ideals agree.
A compositum principalizes them all.  Rank-zero eigenspaces
need no basis.  The kernel and section-module base-change
statements in ICL5 show that these are bases of the actual
new modules, not of proper finite-index submodules.

For a matrix block of size \(d\geq2\), the integral
elementary matrices \(1+tE_{ab}\), \(t\in\mathbb Z\), are
Zariski dense in each root subgroup: a polynomial with
infinitely many integer zeros is zero.  These subgroups
generate \(\mathrm{SL}_d\).  Diagonal matrices whose entries
are roots of unity supply a Zariski dense subgroup of the
diagonal torus, because a nonzero Laurent polynomial in one
variable has only finitely many roots.  The same argument
handles \(d=1\), and \(d=0\) is trivial.  The Zariski closure
of all these matrices is therefore \(\mathrm{GL}_d\).
Their determinants are algebraic units, so both each matrix
and its inverse are integral.  Products across all blocks,
and roots of unity in the theta-generator scalar, are dense
in the whole product.  Every individual point uses only
finitely many roots of unity and hence is defined over a
finite extension.  Transposing the matrices gives exactly
the source frame-action convention and preserves all these
properties.
\end{proof}
% END ICL6

% BEGIN ICL7
\paragraph{ICL7. A nonzero integral-frame evaluation.}
\label{claim:ICL7}
For every nonzero \(f\in C_w\), there are a number field
\(K_f/K\) and an actual framed point \(y_f\) of \(X_j\) over
\(K_f\) such that
\[
 \operatorname{ev}_R(f,y_f)\ne0,\qquad
 \operatorname{ev}_R(g,y_f)\in\mathcal O_{K_f}\quad(g\in L_w).
                                                               \tag{ICL.4}
\]
The map \(g\mapsto\operatorname{ev}_R(g,y_f)\) is \(K\)-linear.
\begin{proof}
Take the CM point of ICL4 where the regular function defined
by \(f\) on its frame fibre is nonzero.  ICL5--ICL6 give an
integral frame outside its zero locus over some finite
extension; call it \(y_f\).  At every finite prime
\(\mathfrak P\) of \(K_f\), complete its integral model.
The chosen global eigenbases become bases of the entire
eigenmodules over this complete DVR, by the kernel argument
in ICL5.  The chosen theta generator becomes a generator of
the entire section module, by cohomology and flat base change.
Together with the rigidified identifications, these are
literally the witnesses in the source predicate
\(\mathrm{Integral}\).

The point has a \(\mathrm{FromTest}\) presentation: its
number field and finite list of coordinates have a reduced
finite-type presentation over \(K\), and its completion is
obtained from that presentation by base change.  The residue
field of the completion is finite, so
\(\mathrm{IntFr.IsAlg}\) holds.  If
\(\mathfrak p=\mathfrak P\cap\mathcal O_K\), the ordinary
ramification index \(e(\mathfrak P/\mathfrak p)>0\) gives the
source \(\mathrm{OverLa}\) condition.  Thus (ICL.1), applied
to this precise point, is among the tests in (ICL.2).
For \(g\in L_w\) it says that the \(\mathfrak P\)-valuation
of \(\operatorname{ev}_R(g,y_f)\) is nonnegative.  This is
true for every finite \(\mathfrak P\), so the value lies in
\(\mathcal O_{K_f}\).  Naturality identifies the completed
evaluation with the image of the same \(K_f\)-value.
Linearity follows from the source evaluation law.  No
denominator depending on \(g\) has been introduced.
\end{proof}
% END ICL7

% BEGIN ICL8
\paragraph{ICL8. Finite generation at every original weight.}
\label{claim:ICL8}
For every original row, dominant weight, pinned moduli datum
and supplied valid uniformization, the whole lattice
\(L_w\) is finitely generated over \(\mathcal O_K\).
\begin{proof}
The fixed-weight cusp space is finite-dimensional by the
classical analytic finite-dimensionality foundation stated
above.  Its application here uses the original finite set
of arithmetic uniformization components.  The reading map
preserves finite linear independence after extension
\(K\hookrightarrow\mathbb C\).  For each reduced finite-type
test algebra \(A\), a relation between readings makes the
corresponding element of \(A\otimes_{K,\sigma}\mathbb C\)
vanish at every complex point.  This scalar extension is
reduced because \(K\) has characteristic zero, and its
complex points detect its elements, so the element is
zero.  Applying \(K\)-linear coefficient functionals
\(\mathbb C\to K\) to this tensor equality gives a
relation between the original rules on every \(A\).
Independence of those rules makes every coefficient
functional vanish on each complex coefficient, hence
each coefficient is zero.  Thus a finite
independent family of rules maps to an independent family
of analytic cusp forms and is bounded in size by their
finite complex dimension.  This is the literal
\(\mathrm{cuspForms\_\allowbreak finiteDimensional\_\allowbreak of\_\allowbreak goodEmb}\)
application; it uses no rational-theta test family.

Apply ICL2 to the evaluations of ICL7.  This proves finite
generation of \(\mathrm{supLat}(T,j,w)\) itself.  The
argument is valid even if that lattice has not yet been
shown to span and even if the cusp space is zero.  It proves
the first conjunct of the original \(\mathrm{FullAt}\)
without any short-theta spanning or denominator assumption.
\end{proof}
% END ICL8

% BEGIN ICL9
\paragraph{ICL9. Clearing a positive power clears the section.}
\label{claim:ICL9}
Let \(\mathcal R\) be a commutative \(K\)-algebra with the
valuation families of ICL1, and let \(S\subseteq\mathcal R\).
Suppose every \(g\in S\) has a nonzero scalar multiple
nonnegative at every one of these valuations.  If
\(f^q\in\operatorname{span}_K S\) for an integer \(q\geq1\),
then a nonzero scalar multiple of \(f\) is nonnegative at
all of them.
\begin{proof}
Write the given power as a finite linear combination.  For
each term replace its generator by a scalar multiple whose
valuations are nonnegative.  Clear the finitely many new
coefficients with one nonzero \(t\in\mathcal O_K\).
Then \(t f^q\) has nonnegative valuation everywhere.  Fix
one valuation and write \(a=\operatorname{ord}_{\mathfrak p}(t)\geq0\).
If \(v(f)=+\infty\), the assertion for \(tf\) is immediate.
Otherwise \(v(f)=x\in\mathbb R\) and \(a+qx\geq0\).
For \(x\geq0\), also \(a+x\geq0\); for \(x<0\),
\(a+x=(a+qx)-(q-1)x\geq0\).  Thus \(tf\) works
simultaneously for every prime and every frame test.
This is why the clearing scalar is chosen integral.
No extraction of a \(q\)-th root and no enlargement of
the fixed field \(K\) is involved.
\end{proof}
% END ICL9

% BEGIN ICL10
\paragraph{ICL10. Powers reach the precise stable range.}
\label{claim:ICL10}
Let \(f\in C_w\) and \(w.n>0\).  Put
\[
 q=4([E:\mathbb Q]r)+4,\qquad r=D.R.r.
\]
There is a dominant exponent pair \((a',n')\), with
\(n'=q\,w.n\geq4([E:\mathbb Q]r)+4\), such that
\(f^q\) belongs to the original cusp space of that exponent.
\begin{proof}
The exponent pair is the sum of \(q\) copies of the original
one.  The defining weak monotonicity inequalities are
preserved by addition; the evenness condition on the theta
exponent is preserved as well.  Naturality and the Borel
weight law multiply, so the product is an algebraic section
of the summed weight.  On every original complex component,
products of the bounded weighted holomorphic readings are
bounded at the summed weight.  At a rational cusp their
Fourier expansions have the usual positive semidefinite
support, with no boundary constant term for a cusp factor.
In a boundary direction, a sum of these nonnegative
exponents can be zero only when every factor has zero
exponent there.  The cusp factor has no such term.  Thus the
product is again cuspidal.  This is the source
\(\mathrm{CuspMul}\) law, iterated on the original rules;
it does not require a chosen basis of the cusp space.
Finally \(w.n\geq1\) gives the displayed inequality.
\end{proof}
% END ICL10

% BEGIN ICL11
\paragraph{ICL11. The exact identity cover and its pins.}
\label{claim:ICL11}
For the original \((X_j,c,l)\) and an original canonical
connection \(\nabla\) satisfying \(\mathrm{IsNab}\), take
\[
 p_d=1,\quad I=\mathcal O_E,\quad l_2=l,\quad X_2=X_j,
 \quad c_2=c,\quad r_2=1,\quad m_2=1,
 \quad c_\delta=1,\quad\phi_\delta=1.                 \tag{ICL.5}
\]
This is the literal identity \(\mathrm{DeltaDat}\).  It
satisfies \(\mathrm{IsLev}\), \(\mathrm{UnitC}\),
\(\mathrm{PsiTie}(1)\), and \(\mathrm{PhiPin}(1,1,0)\).
Every field-valued evaluation of \(c_\delta\) is one.
\begin{proof}
The divisibility and ideal identities are identities at
the same level, \(\gcd(1,l.M)=1\), and
\(I+\overline I=\mathcal O_E\).  Both marked-section
conditions in \(\mathrm{EtaOK}\) reduce to first powers
and the identity maps.  A framed point lifts to itself;
the line-bundle identifications in \(\mathrm{FrOver}\) are
the original identity rigidifications.  The constant rule
one is a weight-zero section and evaluates to one on every
field point, which has a reduced test presentation.
This verifies the last two \(\mathrm{IsLev}\) clauses,
\(\mathrm{UnitC}\), and the trivial \(\mathrm{PsiTie}\).

Both half quotients \(\mathrm{LaI}\) and
\(\mathrm{LaJ}\) are zero.  Hence there is a single tuple
in the latter, its span is the full zero module, and its
constant function has value one.  In the finite Fourier
formula defining \(\mathrm{PhiPin}\), the exponent of the
primitive root includes \(l.M'/p_d=l.M'\).  Each term is
therefore one independently of the chosen representative
of the zero half quotient.  The sum has one term and is
\(\phi_\delta=1\).  This checks the literal pin, including
the source's representative choices.  The constant rule
has order one and \(3\nmid p_d\).  This construction is
an identity cover; it supplies no higher tuple cover.
\end{proof}
% END ICL11

\subsection{The analytic spanning argument and its exact inputs}

We now state the original inputs still required for this
part.  They are kept visible so that the final lattice
conclusion cannot be mistaken for an unconditional closure
of the retained RTG producer.

The principal-chart input supplies, on the same \(X_j,c\),
a valid uniformization \(U'\) whose components have the
same Hermitian matrix and imaginary scalar, whose lattice
stabilizers contain the original principal level, and
whose scalar multiplier \(\chi\) is one.  It also identifies
its metric with the original one.  This is the precise
\texttt{Sup\allowbreak Full.\allowbreak exists\_\allowbreak unf} construction and metric
comparison used in the source; it is not implied by an
arbitrary \(\mathrm{IsUnf}\).

The level-system input supplies the actual lattices
\(\mathcal L\), unitary group \(\Delta\), vacuum functions
\(\Phi(L,\Lambda)\), and unit multipliers \(\mu(\gamma,L)\)
with all fields of the original \(\mathrm{IsSystem}\).
In particular the finite class sum uses precisely its
positive weights \(\mathrm{wt}(M,\Lambda)\), and its norm
field has a summable off-diagonal majorant of total
strictly less than one.  The numerical range is
\(n\geq4([E:\mathbb Q]r)+4\), with \(n\) even.

The raw Gaussian input consists of the source
\texttt{kernel\_\allowbreak identity}, the integrability of that
kernel on a sufficiently thin orthogonal slab, and its
reproducing formula there.  With an invariant finite
harmonic family \(P_b\), write its kernel as \(\mathcal K\).
The formula and normalization used below are
\begin{align}
 \kappa\sum_{\gamma\in\Delta}
 C(L,\gamma M)\overline{\mu(\gamma,M)}
 \mathcal K(x,\gamma y)
 &=\sum_{[\Lambda],b}\mathrm{wt}(M_0,\Lambda)
       \Theta_{\Lambda b}(L,x)
       \overline{\Theta_{\Lambda b}(M,y)},             \tag{ICL.6}\\
 \int_S\mathcal K(x,y)H(y)\,dy&=\rho H(x),\qquad
                    \kappa\ne0,\quad\rho>0.          \tag{ICL.7}
\end{align}
Here \(M_0\) denotes the integer level, to distinguish it
from a lattice \(M\).  The series in (ICL.6) has its
asserted sum, and the kernel is absolutely integrable on
\(S\).  All measures, Gram matrices, and oscillator
normalizations in these formulas are the original ones.
The following claims prove the passage from these raw
inputs to spanning; they do not assume a spanning theorem.

% BEGIN ICL12
\paragraph{ICL12. The class-sum operator has a bounded inverse.}
\label{claim:ICL12}
For the actual supplied level system, set
\[
 C(L,M)=\sum_{[\Lambda]}\mathrm{wt}(M_0,\Lambda)
       \langle\Phi(L,\Lambda),\Phi(M,\Lambda)\rangle.
                                                               \tag{ICL.8}
\]
There are \(a>0\) and \(0\leq\theta<1\) such that
\[
 C(L,L)=a,\qquad
 \sum_{M\ne L}|C(L,M)|\leq\theta a.                  \tag{ICL.9}
\]
On the space of original bounded holomorphic forms with
the prescribed Borel weight and \(\mu\)-covariance,
\((\mathcal CH)(L,x)=\sum_M C(L,M)H(M,x)\)
is invertible and its inverse has norm at most
\(a^{-1}(1-\theta)^{-1}\) in the weighted supremum norm.
\begin{proof}
The class set is finite and nonempty, its weights are
positive, and their sum \(m_0\) is positive.  If the
vacuum norm in \(\mathrm{IsSystem.norm}\) is \(a_0>0\),
the diagonal is \(a=m_0a_0\).  Its majorant \(b(L,M)\)
gives \(|C(L,M)|\leq a b(L,M)\).  Summing gives
(ICL.9), including absolute convergence.  If
\(\mathcal L\) is empty the whole assertion about forms
is the assertion about the zero space; otherwise the
nonnegative off-diagonal sum also permits \(\theta\geq0\).
The unitary invariance of the vacuum inner product gives
\[
 C(\gamma L,\gamma M)\mu(\gamma,L)
                \overline{\mu(\gamma,M)}=C(L,M).
\]
Together with \(|\mu|=1\), this proves that
\(\mathcal C\) preserves the exact covariance convention
\(H(\gamma L,\gamma x)=\mu(\gamma,L)^{-1}H(L,x)\).

The space of bounded forms is complete in the weighted
supremum norm.  A Cauchy sequence converges uniformly after
weighting; on every compact subset of the frame domain
the weight and its inverse are bounded, so the limit is
locally uniformly holomorphic and satisfies the same
transformation laws.  Write \(\mathcal C=a(1+Q)\).
(ICL.9) implies \(\|Q\|\leq\theta<1\), and the uniformly
convergent series
\(a^{-1}\sum_{s\geq0}(-Q)^s\) is its two-sided inverse.
In particular the inverse preserves the holomorphic and
covariance conditions.  This gives the original Neumann
step with its actual class-sum coefficients.
\end{proof}
% END ICL12

% BEGIN ICL13
\paragraph{ICL13. Unfolding a finite-rank kernel.}
\label{claim:ICL13}
Let a countable group \(\Gamma\) act on a countable set
\(I\) and preserve a measure on \(Y\).  Suppose that the
diagonal action is free almost everywhere and has measurable
fundamental pieces \(D_i\) whose translates partition
\(I\times Y\) almost everywhere.  Let
\(f_z(i,y)\) satisfy
\(\sum_i\int_Y|f_z(i,y)|\,dy<\infty\) for every
parameter \(z\) in a set \(Z\).  If on the fundamental
pieces its orbit sum has finite rank in \(z\),
\[
 \sum_{\gamma\in\Gamma}f_z(\gamma i,\gamma y)
              =\sum_{b=1}^s e_b(z)t_b(i,y),            \tag{ICL.10}
\]
then \(z\mapsto\sum_i\int_Y f_z(i,y)\,dy\) lies in
the finite linear span of the functions \(e_b\).
For the original lattice/frame action the required
fundamental pieces exist and the measure is the original
Lebesgue measure on frames, restricted to the slab.
\begin{proof}
Absolute summability permits Tonelli and the changes of
variables \((i,y)\mapsto(\gamma i,\gamma y)\).  The
fundamental pieces partition the orbit space, so the
integral of the orbit sum on these pieces equals
\(\sum_i\int_Y f_z\).  There is a small integrability
point here: the individual \(t_b\) need not initially be
integrable.  Replace the \(e_b|_Z\) by a basis of their
span, and choose finitely many evaluation points
\(z_1,\ldots,z_t\) for which the evaluation matrix is
invertible.  Solving that matrix system expresses each
corresponding coefficient function on the pieces as a
fixed linear combination of the integrable orbit sums
at \(z_i\).  Each coefficient is consequently integrable.
Integrating (ICL.10) now proves the assertion.

For the lattice/frame application, give the countable
set of full lattices the discrete topology.  The action
on its product with the invertible-frame domain is free:
an element fixing a full frame is the identity.  It is
properly discontinuous.  On a compact set there are only
finitely many lattice labels; a matrix carrying one
fixed lattice to another has discrete coefficients in
fixed lattice bases, and the frame and inverse-frame
bounds bound these coefficients.  Only finitely many
such matrices can occur.  Choose a countable collection
of small open sets with disjoint translates and remove
earlier translates to obtain Borel fundamental pieces.
Unitary left multiplication has real Jacobian
\(\prod_\tau|\det\gamma_\tau|^{2r}=1\), and preserves
the slab defined by the original Gram matrices.  Thus
these are precisely the measure-preserving changes of
variables required above, with no rescaling of measure.
\end{proof}
% END ICL13

% BEGIN ICL14
\paragraph{ICL14. The raw original kernel implies theta spanning.}
\label{claim:ICL14}
Assume the level-system and raw Gaussian inputs just
stated.  Every original bounded holomorphic form \(F\)
is a finite linear combination of the actual functions
\(\Theta_{\Lambda b}\), simultaneously for every
\(L\in\mathcal L\) and every frame in the domain.
\begin{proof}
Set \(H=\mathcal C^{-1}F\) using ICL12.  Choose a slab
small enough for (ICL.6)--(ICL.7) and for the source
uniform Borel-factor bounds.  These bounds and boundedness
in the weighted norm give one ordinary bound for
\(|H(L,y)|\) on the slab, independent of \(L\).
For \(z=(L,x)\), apply ICL13 to
\[
 f_z(M,y)=C(L,M)\mathcal K(x,y)H(M,y).
\]
The diagonal term is integrable by the kernel bound;
the off-diagonal terms are dominated by its integral
times the uniform bound for \(H\) and the summable
series in (ICL.9).  This verifies the absolute
summability hypothesis, rather than exchanging
uncontrolled infinite sums.

In the orbit sum, the covariance of \(H\) contributes
\(\mu(\gamma,M)^{-1}=\overline{\mu(\gamma,M)}\).
Therefore (ICL.6) makes that sum a finite linear
combination, in \(z\), of
\(\Theta_{\Lambda b}(L,x)\).  Its coefficient at
\((M,y)\) is exactly
\(\kappa^{-1}\mathrm{wt}(M_0,\Lambda)
\overline{\Theta_{\Lambda b}(M,y)}H(M,y)\).
ICL13 now gives a finite expansion of the unfolded
integral.  On the other hand (ICL.7) and absolute
summability evaluate that integral as
\[
 \rho\sum_M C(L,M)H(M,x)=\rho F(L,x).
\]
Divide by the nonzero real number \(\rho\).  The result
is the asserted expansion, with coefficients independent
of both \(L\) and \(x\).  Its class representatives,
weights, phases and measure are those of (ICL.6).
This proves the source analytic spanning implication
from raw kernel and level-system data.  It does not
prove those raw data from the common geometric family.
\end{proof}
% END ICL14

% BEGIN ICL15
\paragraph{ICL15. Integral harmonic polynomials span the actual space.}
\label{claim:ICL15}
For any full original lattice \(\Lambda\subset E^n\),
the polynomials in \(\mathrm{harmL}(c,\Lambda,a)\)
span the original harmonic highest-weight space after
extension from \(K\) to its fixed complex embedding.
\begin{proof}
All polynomials under consideration have one fixed
finite total degree.  The row-weight, highest-weight
and harmonic conditions are kernels of finitely many
linear maps over \(K\): the original field \(c.F\)
contains the prescribed embeddings of \(E\) and their
conjugates.  Kernel formation commutes with a field
extension.  Choose a \(K\)-basis of that finite
solution space.

Choose a positive integer \(N\) with
\(N\mathcal O_E^n\subseteq\Lambda\), possible because
\(\Lambda\) is full.  Each variable multiplied by
\(N\) is one of the lattice-linear generators of
\(\mathrm{symO}(c,\Lambda)\).  For a basis polynomial
of degree \(d\), multiplication by \(N^d\), followed
by an integral scalar clearing its finitely many
coefficients, puts it in this subring.  The multiplier
is nonzero and preserves all the linear equations.
The resulting scaled polynomials belong to the
literal \(\mathrm{harmL}\) and retain their span.
Renaming their variables by the original letter
dictionary gives the corresponding row-highest
polynomials on the analytic side.  No independently
chosen harmonic representation is substituted.
\end{proof}
% END ICL15

% BEGIN ICL16
\paragraph{ICL16. The exact short-series span over the original field.}
\label{claim:ICL16}
Retain the original principal-chart and finite-level
dictionary input on \(U'\), the identity cover (ICL.5),
the chart-independent construction of its short rules,
and the raw inputs of ICL14.  If
\(n\geq4([E:\mathbb Q]r)+4\), then every section
\(f\in\mathrm{sec}(X_j,c,a,n)\) whose readings in
\(U'\) are bounded belongs to
\[
 \operatorname{span}_K\mathrm{shortGen}
                     (X_j,c,\nabla,1,0,a,n,d).        \tag{ICL.11}
\]
\begin{proof}
Here is the full use of the dictionary input.  It supplies
one reference Hermitian space, original component
lattices \(L_k\) in distinct \(\Delta\)-orbits, and
unitary coordinate changes \(\gamma_k\).  Its vacuum
multiplier is one on every stabilizer of an \(L_k\).
Pull the \(k\)-th reading back by \(\gamma_k\) and
divide by its fixed nonzero comparison scalar \(C_k\).
The principal-level inclusion, the trivial scalar
character, and the stabilizer multiplier show that it
is invariant under the stabilizer with exactly the
given Borel weight.  Extend it to the orbit of \(L_k\)
by the prescribed \(\mu^{-1}\)-covariance, and set it
zero on the other orbits.  The cocycle identity and
stabilizer invariance prove well-definedness.  There
are finitely many \(k\), so the resulting form has
a uniform weighted bound.  ICL14 expands it into
finitely many theta terms.

The lower-row and lattice-translation part of the
retained original dictionary rewrites those terms
using the actual allowed lattices \(\Lambda\) and
the original level \(M=l.M\).  ICL15 expands their
harmonic polynomials using the literal integral
\(\mathrm{harmL}\).  The retained short-rule
construction supplies, for each resulting \(\Lambda,N,P\),
a section \(g\) in \(\mathrm{shortGen}\).  Its exact
reading identity is
\[
 \operatorname{rd}_{U'}(g)_k(x)
   =C_k u_{\Lambda,N}
       \Theta_{\Lambda,P}(L_k,\gamma_k x),
            \qquad u_{\Lambda,N}\ne0.                 \tag{ICL.12}
\]
The scalar \(u_{\Lambda,N}\) is independent of the
component \(k\).  Dividing each theta coefficient by
that scalar gives one finite complex linear
combination of readings of original short rules,
valid on every component.  Off the frame domain all
these readings and the given reading are zero by
their definitions, so equality holds on the entire
source function space.

It remains to descend the coefficients; mere
injectivity of the reading would not suffice.  Its
preservation of finite \(K\)-linear independence
after complex extension was proved in ICL8 directly
from naturality and reduced test rings.  If (ICL.11)
failed, choose a basis of the span of the finitely
many rules in the complex expression and append
\(f\).  This is \(K\)-independent, while its readings
are complex-dependent, a contradiction.  Therefore
the coefficients can be chosen in the original
\(K\).  All assertions here concern the identity
cover; no higher \(\mathrm{CoversAt}\), arbitrary
tuple class sum, or rescaled phase has been supplied.
\end{proof}
% END ICL16

\subsection{The finite-place calculation and the full-lattice endpoint}

The remaining integral-theta input is the following original,
uniformly quantified statement.  For every allowed short
rule \(g=g_{\Lambda,N,P,x_0}\), and every test in (ICL.1),
one may pass to a finite extension of the complete DVR,
of ramification degree \(b>0\), and use the original
higher-level integral chart and torsion data.  With
\(e'=be\), its normalized valuation \(\nu\) satisfies
\(\nu(\iota' z)=e^{-1}\operatorname{ord}_{V}(z)\).
The precise short-value formula, after grouping the
original torsion sum, is
\[
 \iota'(\beta_H g(y))
       =\zeta\,c_0^{-1}\sum_{q}w_q X_q,
 \quad c_0=\#\bigl((\mathrm{Fin}\,n)\to
                              \mathrm{LevQ}(\mathrm{Dw},l_2)\bigr),
                                                               \tag{ICL.13}
\]
where
\[
 X_q=\iota_2\bigl(\mathrm{mIdx}(\Lambda,N,r)^{-1}
                                     N^{-\deg P}\bigr)
       \sum_i\mathrm{hv}\bigl(\mathrm{sqN}(F_{qi}),P\bigr).
\]
The scalar \(\beta_H=\mathrm{betaH}(c,\Lambda)\ne0\)
is the literal source scalar, transported by \(d.eF\)
on the cover.  The weights \(w_q\) are its literal finite
Fourier weights and are integral.  The scalar \(\zeta\)
has a positive finite order.  The torsion points satisfy
the source \(\mathrm{InZh}\), geometric injectivity,
\(\mathrm{IsUUOver}\), \(\mathrm{AllRat}\), and
half-ideal killing conditions, not merely an abstract
finite-set model.

For a squarefree \(M=l_2.M\), put \(\alpha=1\) if
the residue characteristic \(\ell\) does not divide
\(M\), and \(\alpha=\ell\) otherwise.  The local
jet input gives, for every \(q\),
\begin{align}
 \nu\bigl(\operatorname{nrm}(\alpha)X_q\bigr)&\geq-B,
       &\operatorname{nrm}(\alpha)&=\alpha^{d_w},       \tag{ICL.14}\\
 B&=\frac{d_w}{2}\sum_{p\mid\mathrm{condN}}
                        \nu(p)\,\mathrm{vdH}(p),
       &d_w&=\mathrm{DF}(\mathrm{Dw},w).               \tag{ICL.15}
\end{align}
This is the exact combination of the source
\texttt{class\_\allowbreak bound} in \texttt{Exists\allowbreak Smul\allowbreak Short\allowbreak Nonneg}
and \texttt{Host\allowbreak Chart.\allowbreak ch\allowbreak Bound\_\allowbreak host}.  It includes the
ramified case and every integral frame.  It is retained
as an original jet producer, not proved here by an
appeal to integrality at good primes.  The proof below
shows exactly why this bound, with these normalizations,
is sufficient and gives a single global clearing scalar.

% BEGIN ICL17
\paragraph{ICL17. A normalized finite average has a uniform bound.}
\label{claim:ICL17}
Let \(\nu\) be a normalized discrete valuation, let
\(u,c\ne0\), let \(w_q\) be integral, and let
\(\zeta^m=1\) with \(m>0\).  If
\(\nu(uX_q)\geq-B\) for all members of a finite set,
then
\[
 \nu\left(\zeta c^{-1}\sum_qw_qX_q\right)
                 \geq-B-\nu(u)-\nu(c).                \tag{ICL.16}
\]
The assertion remains valid when the average is zero.
\begin{proof}
The root-of-unity equation implies \(\zeta\ne0\) and
\(m\nu(\zeta)=0\), hence \(\nu(\zeta)=0\).
Each \(w_q\) has nonnegative valuation because it is
integral over \(\mathbb Z\) and the valuation ring is
integrally closed.  Thus
\(\nu(w_quX_q)\geq-B\), and repeated use of the
ultrametric inequality gives the same bound for their
sum.  Multiplication by \(c^{-1}u^{-1}\) subtracts
\(\nu(c)+\nu(u)\).  If any summand or the full sum
is zero its valuation is \(+\infty\), which is
compatible with each inequality.  No archimedean
triangle estimate or factor equal to the cardinality
of the summation set is inserted; the only cardinality
denominator is the literal \(c\) already in the formula.
\end{proof}
% END ICL17

% BEGIN ICL18
\paragraph{ICL18. One explicit scalar clears every original place.}
\label{claim:ICL18}
Assume (ICL.13)--(ICL.15) for a permitted original cover
with squarefree \(M=l_2.M\), and use the original host
bound \(0\leq\mathrm{vdH}(p)\leq\Delta_E\), where
\(\Delta_E=|\operatorname{disc}(E)|\).  Define
\[
 P_0=\prod_{p\mid\mathrm{condN}}p,\qquad
 A_w=P_0^{\Delta_E d_w}M^{d_w}c_0,\qquad
 s=A_w\beta_H\in K.                                 \tag{ICL.17}
\]
Then \(s\ne0\) and \(v(sg)\geq0\) for every
original \(v\in\mathrm{valSet}(\mathfrak p)\), at
every \(\mathfrak p\).  Thus the exact short-generator
denominator conclusion holds with the scalar (ICL.17).
\begin{proof}
The prime product is nonzero, \(M\ne0\), the finite
level quotient is nonempty and has positive cardinality,
and \(\beta_H\ne0\) by its source definition.  Hence
\(s\ne0\).  The stated host bound follows directly
from its definition
\(\mathrm{vdH}(p)=\operatorname{ord}_p(\Delta_E)/[E:\mathbb Q]\):
the denominator is at least one, and for each prime
\(p\) used here, \(p^{\operatorname{ord}_p(\Delta_E)}\)
divides the positive integer \(\Delta_E\).  Thus
\(0\leq\operatorname{ord}_p(\Delta_E)\leq\Delta_E\),
which proves both inequalities without a new local
estimate.  Fix any original test, make the finite
extension supplied by (ICL.13), and retain the normalized
valuation \(\nu=\operatorname{ord}_{V'}/e'\).
Because \(M\) is squarefree, \(\ell^2\nmid M\);
the chosen \(\alpha\) is a positive integer dividing
\(M\).  Write \(M=\alpha\alpha'\).
Apply ICL17 to (ICL.13), with
\(u=\operatorname{nrm}(\alpha)=\alpha^{d_w}\)
and \(c=c_0\).  It gives
\[
 \nu(\beta_Hg(y))\geq-B-d_w\nu(\alpha)-\nu(c_0).
\]
Add
\(\nu(A_w)=\Delta_Ed_w\sum_{p\mid\mathrm{condN}}\nu(p)
                  +d_w\nu(M)+\nu(c_0)\).
The resulting lower bound is
\[
 d_w\sum_{p\mid\mathrm{condN}}
        \nu(p)\left(\Delta_E-\frac{\mathrm{vdH}(p)}2\right)
                      +d_w\nu(\alpha')\ \geq\ 0.      \tag{ICL.18}
\]
Every factor is nonnegative.  Natural integers are
integral in the valuation ring, and the host bound
implies the displayed coefficient is nonnegative.
Here the literal degree is
\(d_w=\sum_\tau\sum_{i<\mathrm{Dw.mult}(\tau)}w(\tau,i)\in\mathbb N\);
its nonnegativity is part of its definition, including
when that sum is zero.
Normalized valuation is preserved by the finite
extension, so the original test also gives a
nonnegative value.  The scalar (ICL.17) was fixed
before the test was chosen.  It is therefore one
scalar for all places and all original frames.

For the identity cover of ICL11, all of its Fourier
weights are powers of the original \(\zeta_{l.M}\)
and are integral, \(c_\delta=1\), and one can take
the root order to be one.  Its level remains
\(l.M=g_0g_1\), squarefree by the two distinct
original pin primes.  Hence the assumptions involving
level, phase and squarefreeness used here are exactly
those needed in the source application.
\end{proof}
% END ICL18

In the next three claims the retained producer inputs are
quantified over every original row and every dominant
exponent used in the argument, including each positive
multiple introduced in ICL10.  The original \(D,T,c\),
field embeddings, metric, measure and phases are fixed
throughout.  A chart or short-rule witness is allowed to
depend on the row and exponent precisely as in the
displayed original producer statements; its existence
is not asserted to follow from the finite-generation
half of this chapter.

% BEGIN ICL19
\paragraph{ICL19. The original positive-weight full lattice, with its exact remaining inputs.}
\label{claim:ICL19}
Assume \(T.Q.\mathrm{Pins}\), a valid original
uniformization at the distinguished good embedding,
and the principal-chart, level-system, Gaussian,
short-rule/dictionary and integral-jet inputs specified
above.  For every original row \(j\), every dominant
weight \(w\) with \(w.n>0\),
\[
 \mathrm{FullAt}(T,j,w)
 \quad\text{and}\quad
 \mathrm{FullLat}(C_w,\mathrm{valSet},\mathrm{isVal})
                                                               \tag{ICL.19}
\]
hold on the original \(K\)-space and its whole
valuation-nonnegative lattice.
\begin{proof}
Finite generation is ICL8.  To prove spanning, fix
\(f\in C_w\), and take \(q,a',n'\) from ICL10.
The original cusp readings are bounded.  Their
boundedness passes to the supplied principal charts
without a change of metric.  More explicitly, at
orthonormal frames the theta metric is one and the
\(\tau\)-Lie metric is
\(2\pi|\tau(\xi_0)|\) times the identity.  A Borel
change between the two presentations has section
scalar of modulus one; its diagonal Lie scalars
have moduli equal to the square roots of the ratios
of these fixed positive constants.  The prescribed
weight multiplies the old bounds by the corresponding
finite product.  Taking a sum of these bounds over
the finitely many old components gives a bound on
each new component.  This is the exact metric
comparison used by \texttt{is\allowbreak Bdd\allowbreak Sec\_\allowbreak of\_\allowbreak is\allowbreak Unf}.

ICL16 now writes \(f^q\) in the \(K\)-span of
the original short generators at \((a',n')\).
ICL18 supplies a nonzero clearing scalar for each
one.  ICL9 gives \(t\ne0\) with all valuations
of \(tf\) nonnegative.  Since \(C_w\) is a
\(K\)-subspace, \(tf\in L_w\), and
\(f=t^{-1}(tf)\) belongs to \(K L_w\).  The
reverse inclusion follows from \(L_w\subseteq C_w\).
This proves \(\mathrm{FullAt}\).  The source
\(\mathrm{FullLat}\) places the same lattice inside
the subtype \(C_w\); the inclusion is an isomorphism
onto \(L_w\).  Finite generation and spanning
transport across that inclusion by ICL1.  This
proves both statements in (ICL.19), with no added
flatness, nonvanishing, or integral-basis premise.
\end{proof}
% END ICL19

% BEGIN ICL20
\paragraph{ICL20. The literal \texttt{Lat\allowbreak Full} and canonical-weight specializations.}
\label{claim:ICL20}
Under the inputs of ICL19, let \(T:\mathrm{CoreJ}(D)\)
and \(\kappa\) satisfy the original \(\mathrm{LawsJ}\)
and \(\mathrm{KapIs}\).  Then the literal predicate
\(\mathrm{LatFull}(T,\kappa)\) holds.  The same proof
gives the canonical-weight
\texttt{Integral\allowbreak Cusp\allowbreak Forms\allowbreak Full\allowbreak Lattice} endpoint.
\begin{proof}
Fix an arbitrary \(i\in\mathrm{Fin}(D.R.m)\) and
put \(j=i+1\).  The rank inequalities in \(D\)
give \(D.R.nV>0\).  Choose a member of its
nonempty CM embedding family.  The corresponding
\(\mathrm{KapIs}\) equation gives
\((\kappa_i).n=D.R.nV\); the weight is therefore
positive.  The fields of \(\mathrm{LawsJ}\)
give the original pins and valid uniformization
at this row and distinguished embedding.  ICL19
applies and says exactly
\[
 \mathrm{FullLat}\bigl(T.\mathrm{cuspA}(i+1,
                 \mathrm{wtA}(\kappa_i),(\kappa_i).n),
                 \mathrm{valSet}(X_{i+1},c,\mathrm{algCls})\bigr).
\]
The row was arbitrary, so this is the definition
of \(\mathrm{LatFull}\).  If there are no rows the
universal assertion is vacuous.  For the canonical
weight \(\mathrm{kaWt}\), use its literal equation
\(\mathrm{kaWt}.n=D.R.nV\) and its original
\(\mathrm{kappaExp}\).  The lattice is the same
\(\mathrm{supLat}\) by definition.  This supplies
that specialization without altering the common
coefficient record or selecting a new weight family.
The original geometric inputs named before ICL19
remain premises of both specializations.
\end{proof}
% END ICL20

% BEGIN ICL21
\paragraph{ICL21. Exact transport to the original \(M_K\) and both Boolean levels.}
\label{claim:ICL21}
Retain all inputs of ICL20, the original
\(t:\mathrm{FWTie}(T,R,\kappa,v_0)\), and an
arbitrary original face \(F\).  Then
\((R.FW_i).\mathrm{MKFull}\) holds for every row,
and \((R.\mathrm{toCx}(F)).\mathrm{MLat}(b)\)
holds for \emph{each} \(b\in\{\mathrm{false},
\mathrm{true}\}\).
\begin{proof}
Use the exact map \(\iota_i\) supplied by the tie.
Its rational range is the original \(W_{K,i}\).
The point and norm parts of that same tie identify
\(\iota_i(f)\in M_{K,i}\) with nonnegativity of
every original valuation of \(f\).  This is the
full test equivalence in RTG22, not just an
inclusion of one integral image.  Hence
\[
 M_{K,i}=\iota_i(L_{\kappa_i}).                         \tag{ICL.20}
\]
ICL20 and the source map now give finite generation
and the full rational span of this module.

For either Boolean level, the original good/type
subspace and coordinate map define \(M_j\) and
\(M_{\mathrm{coord}}\) by intersection and exact
inverse image.  RTG23 transports a full ambient
lattice to these original subspaces and through
the specified coordinate isomorphism.  Apply it
using (ICL.20), once with \(b=\mathrm{false}\)
and once with \(b=\mathrm{true}\).  The same
\(T,R,\kappa,v_0,F\), field embedding, group
operators and phase data are used in both
applications.  Thus this is the literal
\texttt{mk\allowbreak Full\_\allowbreak of\_\allowbreak tie}/\texttt{m\allowbreak Lat\_\allowbreak of\_\allowbreak tie}
consequence.  It changes neither the projector
nor either level's coefficientwise map.
\end{proof}
% END ICL21

\subsection{What this packet supplies and what it still requires}

ICL8 supplies the actual finite-generation half of
the original geometric lattice, from the stated
classical CM, abelian-scheme and fixed-weight
automorphic foundations.  It needs none of the
original short-theta producers.  Its simplification
does not claim the source's stronger quantitative
\(\mathrm{CMPt.Good}\) or bounded-radius theorem.
ICL11 constructs the actual identity cover.  ICL12--ICL18
prove the analytic unfolding, coefficient descent,
and global-denominator implications at the original
normalizations.  ICL19--ICL21 are complete implications
from the explicitly retained producer inputs; those
inputs have not all been supplied in this packet.

In particular, the actual principal-chart construction,
the common level system, the raw Gaussian identity and
reproducing estimates, the short-rule gluing and its
simultaneous reading dictionary, and the ramified
integral jet estimate remain a quantified frontier
for the spanning half.  Their exact source ranges
are in the contract and foundation manifests.
Consequently the whole original \texttt{Lat\allowbreak Full}
supplier is not marked unconditionally reconstructed.

The separate higher \(\mathrm{CoversAt}\), all-tuple
\(\mathrm{ClsEx}\) with phase rescaling, one
\(\mathrm{PointIdentCore}\) chosen before all tuple
records, and base \(\mathrm{ClassSumsSpan}\) also
remain.  This packet closes neither
\texttt{DESIGN\_\allowbreak CHOICES\_\allowbreak COMMON} nor
\texttt{BASE\_\allowbreak THETA\_\allowbreak RATIONAL\_\allowbreak TESTS}, and it
does not infer the outer common-family quantifiers,
eta agreement, or the stronger second division law.

\paragraph{Primary literature used for the classical inputs.}
J.-P. Serre and J. Tate, \emph{Good reduction of abelian
varieties}, Annals of Mathematics 88 (1968), 492--517,
especially Theorem 6(a),
\url{https://annals.math.princeton.edu/1968/88-3/p05}.
For the CM algebraicity and uniformization formulation,
J. S. Milne, \emph{Complex Multiplication}, Chapters I--II,
\url{https://www.jmilne.org/math/CourseNotes/CM.pdf}.
For the cubical extension and relative section-module
framework, L. Moret-Bailly, \emph{Pinceaux de varietes
abeliennes}, Asterisque 129 (1985), Chapters I and VI,
\url{https://www.numdam.org/item/AST_1985__129__1_0/}.
For fixed-type automorphic finite-dimensionality,
Harish-Chandra, \emph{Automorphic forms on semisimple
Lie groups}, Lecture Notes in Mathematics 62 (1968),
Theorem 1,
\url{https://doi.org/10.1007/BFb0098434}; the reductive
form is also stated and proved in A. Borel,
\emph{Automorphic forms on reductive groups}, Section 8.4,
\url{https://homepages.math.uic.edu/~rtakloo/papers/borel/borel3.pdf}.
The manuscript retains these foundational results;
the new proofs above spell out their application to
the exact integral-frame tests rather than claiming
to reprove their general theorems.


\section{The uniform logarithmic gain and its actual local-lattice interface}
\label{sec:global-local-gain-v046}

\status{The logarithmic counting theorem in this chapter is a theorem
about every original admitted family, with the original constants and
the error \(b^{x^{-\epsilon/2}}\). The subsequent local-lattice
installation retains the original comparison and growth hypotheses.
Neither a replacement family nor an independent collection of local
lattices is introduced. The whole global local-gain obligation is
closed only if those actual-family hypotheses are supplied.}

The source revision throughout is
\texttt{7cbf7a0c32fca1badf05931879a37c3f8a2b03b4}.
Source declarations specify the objects being reconstructed; the
arguments below supply the mathematical proofs. Previously proved
local projection theorems TC6, EA9, GI8, GII8, BC8 and GJ4 are reused.
In particular the counting argument below does not posit a new local
projection theorem.

\subsection{The original quantities}

Fix \(0<\epsilon<1/52\), an original member \(D\), its size datum \(L\),
and its instance \(X\). Write
\[
 \begin{gathered}
 b=\exp(x^\epsilon),\quad a=xb,\quad
 X'=\exp(x^{\epsilon/3}),\quad b'=a/X',\\
 \delta=x^{-\epsilon/8},\quad
 w=b'/(xn),\quad X''=a/w=xnX',\quad
 p_0=k^2x^{26},\\
 N_m=k+m(t+1),\qquad S_c=4k\Delta,\qquad c=4S_c,
 \qquad \eta=L.\mathrm{etaSz}.
 \end{gathered}                                                    \tag{GLG.1}
\]
Here \(m\ge2\), \(k\ge m+3\), \(n=|HH|\ge1\), \(f=2n\),
and \(t\) is the original positive odd integer. A prime always means
a rational prime. No bound on the functions \(b_s\) from above will
be imposed. The original window is
\[
 a\le a_s\le a+\frac{a}{kx^{13}},\qquad b_s\ge b'
                         \quad(s\in TT).                         \tag{GLG.2}
\]
The \(a_s,b_s\) are integers, although all estimates below are real.

The constants \(K\) are the original nonnegative record
\(\mathrm{Consts}\). To make the formulas readable set
\[
 \begin{gathered}
 K_D=K.\mathrm{DII},\quad K_I=K.\mathrm{DI},\quad
 K_T=K.\mathrm{tII},\quad K_P=K.\mathrm{Pr1},\\
 K_1=K.\mathrm{kappaThresholdOne},\quad
 K_2=K.\mathrm{kappaThresholdTwo}.
 \end{gathered}
\]
All other fields, notably \(K.\mathrm{dl1}\), \(K.\mathrm{dl2}\),
\(K.\mathrm{La}\) and \(K.\mathrm{kII}\), retain their original
meaning. Put
\[
 \begin{split}
 r&=\left\lceil8(4+8\epsilon)/\epsilon\right\rceil+1,\\
 L_1(p)&=\left\lceil\log_p((kx^{26})^4)\right\rceil,\qquad
 \Lambda=K.\mathrm{La}\,N_1^2\log a,\qquad Q=Q_p .
 \end{split}                                                       \tag{GLG.3}
\]
The \(N_1,G_p,L_2\) below are exactly the original orbit-window count,
good-window length and twice the supremum of \(L_1(p)\) over inert
primes. All ceilings of nonnegative real numbers are natural ceilings.

For \(1\le j\le m\), put
\(\kappa_j=[k+(2j-1)(t+1)]/2\).
The original filling function is \(F_j(z)=X.\mathrm{fill}(b,j,z)\):
at a top \((0,s)\) it is \(b_s+o_j(0,s)\), where
\[
 b'/x\le o_j(0,s)\le b'/x+\eta.
\]
At a thin place its offset lies in \([w,w+\eta]\), and at a thick
place in \([b,b+\eta+1]\). Set
\(F_j^G=F_j-\kappa_j\), \(\alpha_s=a_s-2c\), and
\(U_j(s,e)=F_j(\operatorname{partner}(p,s,e))-\kappa_j+c\).
Thus \(F_j,F_j^G,U_j\) are the source's \(\mathrm{fill}\),
\(\mathrm{fillG}\), and \(\mathrm{capU}\).

The additive orbit step is \(h_p\) at inert primes and \(2h_p\)
at signal primes. Let \(R_p\) be the chosen representatives, let
\(T_v\) be the tops with representative \(v\), and let \(N_v=|T_v|\).
The source's \(n_p\) is \(\operatorname{ord}(h_p)\) at inert primes,
and \(2\operatorname{ord}(2h_p)\) at signal primes. The subset \(V\)
of a cyclic orbit consists of its even positions, measured from the
same chosen representative. A \emph{candidate} prime is inert,
\(p<X'\), of order at least \(x^{\epsilon/4}\), and satisfies the
original exclusion-band condition for \(p\ge p_0\). In particular
its order \(n'\) is even. These are the original \(\mathrm{Cand}\)
and \(\mathrm{IsV}\), with the orbit dictionary proved in GJ3.

Write \(u_+=\max(0,u)\) and \(\gamma_p=p/(p^2-1)\). The relevant
local quantities are
\[
 \begin{split}
 A_j(s)&=\sum_{\substack{e\ge1\\ e\ {\rm odd}}}
       (a_s/p^e-F_j(\operatorname{partner}(p,s,e)))_+,\\
 A_j^G(s)&=\sum_{\substack{e\ge1\\ e\ {\rm odd}}}
       (a_s/p^e-F_j^G(\operatorname{partner}(p,s,e)))_+,\\
 M_j(s)&=\sum_{\substack{e\ge1\\ e\ {\rm odd}}}
       \min(a_s/p^e,F_j(\operatorname{partner}(p,s,e))),\\
 d_v&=K.\mathrm{dl1}\,n_pN_m^2\log_p(2a)
          +K.\mathrm{dl2}\,N_v\log_p(2a)^2,\\
 g_v^0&=\sum_{s\in T_v}
       \sum_{\substack{e\ge1\ {\rm odd}\\\neg\mathrm{Excl}(p,s,e)}}
          (\alpha_s/p^e-U_j(s,e))_+-d_v,\\
 g_v^U&=\gamma_p\sum_{s\in T_v}\alpha_s
       -\sum_{s\in T_v\setminus V}\sum_{e\ {\rm odd}}
          \min(\alpha_s/p^e,U_j(s,e))-d_v .
 \end{split}                                                       \tag{GLG.4}
\]
They are \(\mathrm{askP},\mathrm{askG},\mathrm{Mn},\mathrm{dl},
\mathrm{g0},\mathrm{gU}\), respectively. Positive-part sums have
finite support; the minimum sums converge by comparison with
\(\sum a_s/p^e\).

\subsection{The finite counting argument}

% BEGIN GLG1
\paragraph{GLG1. numerical bounds on the original filling data.}
\label{claim:GLG1}
Under the original \(\mathrm{ParameterChoice}'\),
\[
 \begin{gathered}
 k\ge5,\quad N_m\le S_c,\quad S_c\ge100,\quad
 500\le w,\quad kS_c\le w\le b'/x\le b',\quad 2b'\le a,\\
 0\le\kappa_j\le N_m\le S_c,\qquad
 F_j\ge1,\quad U_j\ge0,\quad
 0\le\alpha_s\le6a/5. 
 \end{gathered}                                                     \tag{GLG.5}
\]
Also \(F_j^G\ge1\) whenever \(b_s\ge0\) on \(TT\) and
the original size condition \(N_m+1\le w\) holds.

\emph{Proof.}
The rank conditions give \(k\ge5\) and
\(\Delta\ge N_m m^{10}\ge N_m\ge k\).
Since \(S_c=4k\Delta\), this gives \(N_m\le S_c\) and
\(S_c\ge100\). The parameter inequality \(kS_c\le w\)
then gives \(w\ge500\). The formula for \(w\), \(n\ge1\)
and \(x>1\) give its two upper bounds. Since \(X'>2\),
\(2b'\le a\). The formula for \(\kappa_j\), \(1\le j\le m\)
and \(N_m=k+m(t+1)\) give \(0\le\kappa_j\le N_m\).
Every filling is at least \(w\): a top is at least \(b'/x\),
a thin place is at least \(w\), and a thick place is at least
\(b>w\), because \(w=b/(nX')\). This proves the filling assertions.
Furthermore \(c=4S_c\le w\le b'\), so \(2c\le a\).
Finally \(kx^{13}\ge5\), and (GLG.2) gives
\(a_s\le a+a/5=6a/5\). Subtracting \(2c\) proves the last
claim. All these comparisons hold for the actual parameters;
no asymptotic replacement of \(a_s\) or \(b_s\) has occurred.
\hfill\(\square\)
% END GLG1

% BEGIN GLG2
\paragraph{GLG2. finite shifts of the odd geometric sum.}
\label{claim:GLG2}
Let \(p\ge2\), \(c\ge0\), \(2c\le y\le A\), \(A\ge1\).
Suppose \(W_e\ge1\), \(0\le U_e\le W_e+c\). For any subset
\(E\) of the positive odd integers,
\[
 \sum_{e\in E}(y/p^e-W_e)_+
 \le\sum_{e\in E}((y-2c)/p^e-U_e)_+
                                  +3c\log_p A.                    \tag{GLG.6}
\]
If \(0\le\kappa\) and \(W_e-\kappa\ge1\), then
\[
 0\le\sum_{e\ {\rm odd}}(y/p^e-W_e+\kappa)_+
             -\sum_{e\ {\rm odd}}(y/p^e-W_e)_+
                      \le\kappa\log_p A                           \tag{GLG.7}
\]
for \(0\le y\le A\). Moreover
\[
              A_j(s)+M_j(s)=\gamma_p a_s .                         \tag{GLG.8}
\]

\emph{Proof.}
In (GLG.6), a positive term on the left requires \(p^e<y\le A\).
For each such exponent the difference of the positive parts is at
most \(2c/p^e+c\le3c\). There are at most \(\log_p A\)
positive integer exponents with \(p^e<A\). All other left terms
vanish, while their right terms are nonnegative. This proves
(GLG.6), including \(A<p\), when both contributing sets on the left
are empty. For (GLG.7), the pointwise difference is between zero
and \(\kappa\); it can be nonzero only if
\(y/p^e>W_e-\kappa\ge1\). The same exponent count applies.
For (GLG.8) use
\((u-v)_++\min(u,v)=u\) for \(u,v\ge0\), and sum
\(\sum_{e\ge1,\ e\ {\rm odd}}p^{-e}=p/(p^2-1)\).
The geometric bound also justifies the interchange with every
finite sum over tops.
\hfill\(\square\)
% END GLG2

% BEGIN GLG3
\paragraph{GLG3. the actual cyclic orbit and the nearest even predecessor.}
\label{claim:GLG3}
For a candidate prime and a representative \(v\), the map
\[
 \phi:\mathbb Z/n'\mathbb Z\longrightarrow HH,\qquad
       i\longmapsto v+i h_p,\qquad n'=\operatorname{ord}(h_p),
                                                                    \tag{GLG.9}
\]
is injective, identifies the original tops and even positions,
and sends addition of \(e\) to the original partner
\((0,\phi(i)+e h_p)\). For an odd top \(y\), let \(d(y)\) be
the least positive backward distance to an even top, or zero
if no such top exists. Put \(x(y)=y-d(y)\) when \(d(y)>0\).
For nonnegative \(a_x,W_y\),
\[
 \sum_{\substack{y\ {\rm odd\ top}\\d(y)>0}}
          \min(W_y,a_{x(y)}/p^{d(y)})
 \le\sum_{x\ {\rm even\ top}}\sum_{e\ {\rm odd}}
          \min(a_x/p^e,W_{x+e}).                                  \tag{GLG.10}
\]
Let
\[
 S_{\rm kink}=\sum_{\substack{y\ {\rm odd\ top}\\d(y)>0}}
                      \min(b_y,a_{x(y)}/p^{d(y)}).
 \]
Then, on the actual orbit,
\[
 \sum_{s\in T_v\setminus V}(F_j^G(0,s)-b_s)+S_{\rm kink}
       \le\sum_{s\in T_v\cap V}M_j(s)+N_v(b'/x+\eta).              \tag{GLG.11}
\]

\emph{Proof.}
If \(\phi(i)=\phi(i')\), the order of \(h_p\) divides \(i-i'\),
so the residues coincide. The representative function is constant
on an orbit and fixes its representative. Its fiber is therefore
exactly the image of (GLG.9). This proves the dictionary, including
the cardinal \(N_v\), without replacing the selected representative.
The evenness of \(n'\) makes parity well defined modulo \(n'\);
thus \(d(y)\) is odd. The pair \((x(y),d(y))\) determines
\(y=x(y)+d(y)\), so the assignment in (GLG.10) is injective.
Its terms occur among the nonnegative terms on the right.

For (GLG.11), use \(W_y=F_j(0,y)\) in (GLG.10).
At a top \(F_j(0,y)=b_y+o_j(0,y)\ge b_y\), while
\[
 F_j^G(0,y)-b_y=o_j(0,y)-\kappa_j\le b'/x+\eta.
 \]
Consequently each left summand of (GLG.11) is bounded by
the corresponding summand of (GLG.10) plus \(b'/x+\eta\).
For a top with no predecessor only this last bound is used.
There are at most \(N_v\) such tops. The partner dictionary
identifies the right side of (GLG.10) with the even-top
minimum sums \(M_j\), proving the assertion.
\hfill\(\square\)
% END GLG3

% BEGIN GLG4
\paragraph{GLG4. an explicit numerical budget for one good orbit.}
\label{claim:GLG4}
Set \(l=\log x\), \(l_p=\log_p a\), \(N=N_v\), and
\[
 \begin{gathered}
 G_0=C_{31}l/\delta,\qquad U=w+\eta+c+N_m,\qquad c_0=2b'+2c,\\
 E_1=\Lambda a/(pk)+ap^{-L_1(p)},\qquad
 E_2=a/p^2+\Lambda a/(pk),\\
 W_D=K_D(N_1^3+r^5+k^2)l_p,\qquad
 E_I=K_Ix^{26}(N_1^5+k^3)l_p,\\
 D_I(n')=\lceil40n'\delta E_I\rceil,\qquad
 c_6(n')=8(D_I(n')+1)+40K_2n'\delta k l_p
                       +160n'\delta c+2n'c .
 \end{gathered}                                                     \tag{GLG.12}
\]
Assume the original \(\mathrm{adm1}(C_{31})\), \(C_{31}>0\), and
the following numerical conditions:
\[
 \begin{gathered}
 L_2<\lceil1/\delta\rceil,\quad 3G_p\le x^{\epsilon/4},
 \quad l\ge1,\quad G_0U\le b',\quad c_0\le3b',\\
 kQ_p\le b'\ (p<p_0),\qquad W_D\le b'\ (\text{every prime }p),\\
 C\ge3,\quad C\ge K_Tr,\quad
 C\ge13+4K_T+160C_{31},\quad C\ge K_PC_{31},\\
 K_1N_1\le Cl ,
 \end{gathered}                                                     \tag{GLG.13}
\]
and, for \(p_0\le p<X'\) and every \(n'\le n\),
\[
 \begin{split}
 B(n',p):={}&K_PG_0(\eta+c+N_m)
       +K_2(Q_pc_6(n')+kQ_p+kl_p+c)\\
       &+2G_0c+40C_{31}lE_I+3\le b_o .
 \end{split}                                                       \tag{GLG.14}
\]
These are precisely the original \(\mathrm{NumStar}\) inequalities.
For every good candidate orbit,
\[
              N\ge1,\qquad n'\delta\le NC_{31}l,\qquad n'\le NG_0 .
                                                                    \tag{GLG.15}
\]
Moreover
\[
 \sum_{s\in T_v}A_j(s)+\sum_{s\in T_v\cap V}M_j(s)
                  -3Nc\log_p(2a)\le g_v^U+d_v .                   \tag{GLG.16}
\]

\emph{Proof.}
The order bound and the first two conditions in (GLG.13) activate
the original good-orbit count in \(\mathrm{adm1}\):
\(n'\delta/(C_{31}l)\le N\).
Its positive numerator gives \(N>0\), hence \(N\ge1\);
the other two assertions follow by multiplication.

For an odd-position top, apply (GLG.6) with
\(y=a_s\), \(A=2a\), and \(U_e=U_j(s,e)\). The hypotheses
hold by GLG1, and the geometric identity for \(a_s-2c\)
turns the right side into
\[
 \gamma_p\alpha_s-\sum_{e\ {\rm odd}}
                         \min(\alpha_s/p^e,U_j(s,e)).
 \]
For an even-position top, (GLG.8) gives
\(A_j(s)+M_j(s)=\gamma_p(\alpha_s+2c)\).
As \(p<X'<a\), \(\log_p(2a)\ge1\); also \(\gamma_p\le1\).
Thus \(2c\gamma_p\le3c\log_p(2a)\).
Sum the two kinds of tops and use the definition of \(g_v^U\).
This proves (GLG.16) with its original loss, independently of
the size of any \(b_s\) above its lower bound.
\hfill\(\square\)
% END GLG4

% BEGIN GLG5
\paragraph{GLG5. the good-orbit inequality below \(p_0\).}
\label{claim:GLG5}
Under GLG4, let a good candidate prime satisfy \(p<p_0\).
The original quantities on this orbit are
\[
 \begin{split}
 \Pr_v&=r\{N(2b'+c)+n'U\},\\
 t_v&=K_TrN(E_1+c_0+kQ)+2rn'c,\\
 D_p&=\left\lceil40n'\delta W_D+2rn'c
                                      +40rn'\delta c_0\right\rceil .
 \end{split}                                                       \tag{GLG.17}
\]
For every original lift \(j\),
\[
 \sum_{s\in T_v}\{A_j(s)-\operatorname{lo}(C,b_o,p)\}
       \le g_v^U+d_v-\Pr_v-\lceil t_v\rceil-D_p-1.                 \tag{GLG.18}
\]
This is the full original \(\mathrm{StarAt}\) in this prime range.

\emph{Proof.}
All arguments of the ceilings are nonnegative. From (GLG.15)
and \(G_0U\le b'\), both \(n'U\) and \(n'c\) are at most
\(Nb'\). The deliberately coarse bound
\(\Pr_v\le6rNb'\) follows. The part of \(t_v\) other than
\(K_TrNE_1\) is at most \((4K_T+2)rNb'\).
The argument of the final ceiling is at most
\[
       40C_{31}Nl b'+2rNb'+120rC_{31}Nl b'.
 \]
The two ceilings and the final \(1\) cost at most \(3\).
Since \(N,r,l,b'\ge1\), the sum of these non-\(E_1\) terms
is at most
\[
              (13+4K_T+160C_{31})Nr b'l .
 \]
The remaining \(E_1\) term is at most \(NClE_1\), using
\(K_Tr\le C\) and \(l\ge1\). Consequently
\[
       \Pr_v+\lceil t_v\rceil+D_p+1
                        \le NCr b'l+NClE_1.                     \tag{GLG.19}
\]
In this prime range the original definition is
\[
 \operatorname{lo}(C,b_o,p)
             =C(c+\eta)\log_p(2a)+Crb'l+ClE_1.
 \]
Its first term absorbs \(3c\log_p(2a)\) by \(C\ge3\).
Subtract (GLG.19) from (GLG.16) and discard the nonnegative
even-top sum \(M_j\). This is (GLG.18).
\hfill\(\square\)
% END GLG5

% BEGIN GLG6
\paragraph{GLG6. the good-orbit inequality from \(p_0\) to \(X'\).}
\label{claim:GLG6}
Under GLG4, let \(p_0\le p<X'\) be a good candidate prime.
The original quantities are
\[
 \begin{split}
 \Pr_v&=\sum_{s\in T_v\setminus V}(F_j^G(0,s)-b_s)+K_Pn'U,\\
 t_v&=S_{\rm kink}+K_1NN_1E_2
            +K_2N(Qc_6(n')+kQ+kl_p+c)+2n'c,\\
 D_p&=D_I(n').
 \end{split}                                                       \tag{GLG.20}
\]
Then the same \(\mathrm{StarAt}\) inequality (GLG.18) holds.

\emph{Proof.}
Use GLG3 to combine the first term of \(\Pr_v\) with
\(S_{\rm kink}\). Their sum is at most
\(\sum_{T_v\cap V}M_j+N(b'/x+\eta)\).
The term \(K_Pn'w\) is at most
\[
 NK_PG_0w\le NCw\,l/\delta.
 \]
The \(K_1NN_1E_2\) term is at most \(NClE_2\).
All remaining nonnegative terms, including the two ceiling
errors and the final \(1\), are at most
\[
 \begin{split}
 N\{&K_PG_0(\eta+c+N_m)
       +K_2(Qc_6(n')+kQ+kl_p+c)\\
       &+2G_0c+40C_{31}lE_I+3\}\le Nb_o.
 \end{split}
 \]
Here \(D_I(n')\le40n'\delta E_I+1\),
\(\lceil t_v\rceil\le t_v+1\), and (GLG.15) were used.
Thus
\[
 \Pr_v+\lceil t_v\rceil+D_p+1
 \le \sum_{T_v\cap V}M_j
       +N\{b'/x+\eta+Cwl/\delta+b_o+ClE_2\}.                      \tag{GLG.21}
\]
The original definition in this range is
\[
 \operatorname{lo}(C,b_o,p)
   =C(c+\eta)\log_p(2a)+b'/x+\eta+Cwl/\delta+b_o+ClE_2 .
 \]
Combine (GLG.16) and (GLG.21) and again use \(C\ge3\).
Every term, including the coefficient \(1\) of \(b_o\), is
unchanged. In particular no upper restriction on the \(b_s\)
was needed in the cancellation.
\hfill\(\square\)
% END GLG6

% BEGIN GLG7
\paragraph{GLG7. a single threshold proves the original \(\mathrm{RestStar}\).}
\label{claim:GLG7}
Define
\[
 \begin{split}
 c_G(\epsilon)&=\{(5+8\epsilon)10^{6r+1}+6\}(6r+1)2,\\
 c_N(\epsilon)&=18c_G(\epsilon)+2(5+8\epsilon)+2,\\
 C_*&=3+K_Tr+(13+4K_T+160C_{31})
                    +K_PC_{31}+K_1|c_N(\epsilon)|.
 \end{split}                                                       \tag{GLG.22}
\]
For every original family \(\mathrm{Fm}\), nonnegative \(K\),
\(C_{31}>0\) and real \(C_e,C_m\), choose \(C_0=C_*\).
For every \(C\ge C_0\), and every pair of original families
of instances \(X,Y\) satisfying the full
\(\mathrm{AdmFam}(\mathrm{Fm},K,X,Y,C_{31},C_e,C_m)\),
there is \(x_0\) such that, for every \(x\ge x_0\) and every
member \(i\), (GLG.18) holds simultaneously for every
original \(a_s,b_s\) in (GLG.2), every candidate prime,
every good orbit and every lift, with
\[
          b_o=b^{\,x^{-\epsilon/2}}=\exp(x^{\epsilon/2}).           \tag{GLG.23}
\]
This proves the literal \(\mathrm{RestStar}\), including the
order \(C_0\), then \(C\), then \(X,Y\), then \(x_0\).

\emph{Proof.}
Only \(\mathrm{adm1}\), \(\mathrm{adm5}(\epsilon/48)\),
the original parameter conditions and \(\mathrm{SizesAt}\)
are needed from the admitted record. The other fields are
retained in the quantified antecedent.
The definitions with their natural ceilings give, for \(x\ge e\),
\[
 G_p\le c_G(\epsilon)\log x\,x^{\epsilon/8},
             \qquad N_1\le c_N(\epsilon)\log x.                  \tag{GLG.24}
\]
Indeed the small-multiplicity ceiling is at most
\((5+8\epsilon)10^{6r+1}\log x+1\), and
\(\lceil x^{\epsilon/8}\rceil\le2x^{\epsilon/8}\).
Their product gives the first estimate. For the second,
the two arguments of the orbit-window maximum are
\(6eG_p\delta\) and \((5+8\epsilon)\log_2x\);
use \(e<3\), \(\log2>1/2\), and allow \(2\) for its
ceiling and added \(1\). This gives exactly \(c_N\).
Also \(k\le x^6\) gives
\(L_1(p)\le\lceil186\log x\rceil\) for every prime
(\(4(6+26)/\log2<186\)); hence \(L_2<514\log x\).

For completeness the only limiting facts used are, for
\(\alpha>0\) and fixed \(A\),
\[
 A\log x=o(x^\alpha),\qquad
 Ax^\beta=o(x^\alpha)\quad(\beta<\alpha).
 \]
Writing \(x=e^z\), the first follows, for example, from
\(e^{\alpha z}\ge(\alpha z)^2/2\); the second is
\(Ax^{\beta-\alpha}\to0\). Therefore one threshold gives
\[
 \begin{gathered}
 514\log x\le x^{\epsilon/8},\qquad
 3c_G\log x\le x^{\epsilon/8},\\
 4C_{31}\log x\le x^{1-\epsilon/8},\qquad x\ge e .
 \end{gathered}
 \]
The original sizes give \(\eta\le w\) and \(N_m,c\le w\).
Thus \(U\le4b'/x\), and these inequalities prove
\(G_0U\le b'\), \(c_0\le3b'\) for \(x\ge2\), and the
first five conditions in (GLG.13). Equations (GLG.22)--(GLG.24)
give all the stated lower requirements on \(C\).

Here is the full uniform estimate for the remaining terms.
Put \(\beta=23\epsilon/48\) and \(Z=\exp(x^\beta)\).
After increasing one threshold, \(Z\ge130\); all fixed
nonnegative constants used below are at most \(Z\); and
\[
 \begin{gathered}
 \eta,c,N_m,n,k,x^{13},\log x,\log_p a,G_0,N_1\le Z,\\
 40n'\delta\le Z\quad(n'\le n),\qquad Q_p\le Z\quad(p<X'),\\
 Q_p\le Z\quad(p<p_0).
 \end{gathered}                                                     \tag{GLG.25}
\]
Every assertion is uniform in \(i,p,n'\). Here are explicit
checks of the possibly non-polynomial ones. The original
sizes say \(n\le x^3\), \(k\le x^6\),
\(\log\eta\le x^{\epsilon/4}\), and
\(\log\Delta\le x^{\epsilon/4}/2\).
The inequality \(14\log x\le x^\beta\) bounds all powers
up to \(x^{14}\) by \(Z\). Since \(16k\le x^7\) for
\(x\ge40\), it gives \(c=16k\Delta\le Z\).
Next \(\log_pa\le2\log a\le x^2\),
\(G_0=C_{31}\log x\,x^{\epsilon/8}\le x^2\),
\(N_1\le x\), and \(40n'\delta\le x^4\) eventually.
Finally \(\mathrm{adm5}(\epsilon/48)\) gives
\[
 \log Q_p\le x^{11\epsilon/24+\epsilon/48}=x^\beta
       \quad(p<X'),
 \]
and its small-prime clause gives the still smaller exponent
\(\epsilon/8+\epsilon/48\) when \(p<p_0\).

It follows that \(kQ_p\le Z^2\) and
\[
         W_D\le Z^5+Z^7+Z^4\le Z^8 .
 \]
Choose the same threshold so that
\[
 14x^\beta\le x^\epsilon,\quad
 4x^{\epsilon/3}\le x^\epsilon,\quad
 13x^\beta\le x^{\epsilon/2}.
 \]
These are possible because \(\beta<\epsilon/2<\epsilon\).
They imply
\(Z^8\le x\exp(x^\epsilon-x^{\epsilon/3})=b'\) and
\(Z^{13}\le b_o\).
To check the entire expression (GLG.14), first use
\[
 D_I(n')\le
 Z\cdot Z\cdot Z^2(Z^5+Z^3)Z+1 .
 \]
Substitute (GLG.25) and this ceiling bound into \(B(n',p)\).
Its resulting polynomial is
\[
 8Z^{12}+40Z^{11}+8Z^{10}+40Z^9+Z^6
                     +6Z^4+5Z^3+19Z^2+3
          \le130Z^{12}\le Z^{13}\le b_o.                           \tag{GLG.26}
\]
This proves every condition of GLG4. Intersect the finitely many
scalar thresholds with the admitted-family threshold at
\(\epsilon/48\), and with the family's parameter threshold.
GLG5 and GLG6 then apply to all members and all window functions
at once. In particular the threshold has not been selected
after the orbit, lift, prime, or top weights.
\hfill\(\square\)
% END GLG7

% BEGIN GLG8
\paragraph{GLG8. from good-orbit inequalities to the original gain.}
\label{claim:GLG8}
Assume \(\mathrm{StarAll}(X,K,C,b_o)\), the original parameter
conditions, \(b_o\ge0\), \(C\ge3\), and
\[
 C\ge2K.\mathrm{dl1}+C_e',\qquad C\ge4K.\mathrm{dl2},\qquad
 e_p\le C_e'f k^2\log_pa\quad(p\le X'').                           \tag{GLG.27}
\]
The \(e_p=X.\mathrm{ep}(p)\) are the original nonnegative
integer exponents. Then for all \(a_s,b_s\) in (GLG.2), all
inert or signal primes \(p\le X''\), and all lifts,
\[
 0\le\mathrm{gain}_j(p)\le4|TT|a/p,                               \tag{GLG.28}
\]
and
\[
 \sum_{j=1}^m\sum_{s\in TTP(p)}
                 \{A_j(s)-\operatorname{lo}(C,b_o,p)\}
       -m\,\operatorname{fl}(C,p)
        \le\sum_{j=1}^m\mathrm{gain}_j(p)-\tau_p .                 \tag{GLG.29}
\]
Here every term is the original term:
\[
 \begin{split}
 \operatorname{fl}(C,p)&=CfN_m^2\log_pa+Cf(\log_pa)^2,\\
 \mathcal V_p&=\begin{cases}
      \{\text{good representatives}\},&
      p\text{ candidate and }\mathrm{Sig}(p)\le x^{-\epsilon},\\
      \varnothing,&\text{otherwise},
      \end{cases}\\
 \tau_p&=\begin{cases}
      0,&\mathcal V_p=\varnothing,\\
      e_p+m\sum_{v\in\mathcal V_p}(\lceil t_v\rceil+D_p+1),
                                  &\mathcal V_p\ne\varnothing,
      \end{cases}\\
 \mathrm{gain}_j(p)&=
      \sum_{v\in\mathcal V_p}\max(g_v^U-\Pr_v,(g_v^0)_+)
      +\sum_{v\in R_p\setminus\mathcal V_p}(g_v^0)_+ .
 \end{split}                                                       \tag{GLG.30}
\]
The last expression is replaced by zero if \(p\) is inert
and \(h_p=0\), or if \(p>X''\). The original \(TTP(p)\)
is all \(TT\) at signal primes or \(p\ge X'\), and otherwise
the tops whose representative lies in \(\mathcal V_p\).

\emph{Proof.}
The representatives partition \(HH\) into equal-length orbits.
Hence \(|R_p|n_p\le f\) (the signal convention multiplies
the orbit length by \(2\)), while \(\sum_{v\in R_p}N_v=|TT|\le f\).
Since \(\log_p(2a)\le2\log_pa\), summing (GLG.4) gives
\[
 \sum_{v\in R_p}d_v
 \le2K.\mathrm{dl1}\,fN_m^2\log_pa
                      +4K.\mathrm{dl2}\,f(\log_pa)^2 .
 \]
As \(k\le N_m\), (GLG.27) implies
\[
                     \sum_{v\in R_p}d_v+e_p
                                      \le\operatorname{fl}(C,p). \tag{GLG.31}
\]
Every \(\Pr_v\) is nonnegative. In the large-prime branch
its top terms equal \(o_j-\kappa_j\ge b'/x-S_c\ge0\);
in the small branch this is immediate. The other losses
are also nonnegative. Thus both \(g_v^0\) and \(g_v^U\)
are at most \(\gamma_p\sum_{T_v}\alpha_s\), and so are
the positive maxima appearing in the gain. Since
\(\gamma_p\le4/(3p)\) and \(\alpha_s\le6a/5\), summing
over the partition even gives the stronger upper bound
\(8|TT|a/(5p)\). This proves (GLG.28), including the
branches where the gain is defined to be zero.

If \(\mathcal V_p\ne\varnothing\), \(p\) is a candidate
and \(\mathcal V_p\) is the whole good-representative set.
Sum \(\mathrm{StarAt}\) over these representatives, and
then over the \(m\) lifts. The left side is the sum over
the actual \(TTP(p)\). The \(g_v^U-\Pr_v\) terms are
bounded by the selected gain maxima. The accumulated
rounding terms are exactly \(\tau_p-e_p\). Finally
\[
 e_p+m\sum_{\mathcal V_p}d_v
   \le m\bigl(e_p+\sum_{R_p}d_v\bigr)
   \le m\,\operatorname{fl}(C,p),
 \]
because \(m\ge1\) and \(e_p,d_v\ge0\). This proves (GLG.29).

Suppose next \(\mathcal V_p=\varnothing\). If \(TTP(p)\)
is empty, (GLG.29) follows from nonnegative gain and
\(\operatorname{fl}\). Otherwise \(p\) is signal or
\(p\ge X'\). At a signal prime the excluded terms in
\(\mathrm{g0}\) do not occur: an odd partner has first
coordinate \(1\), whereas an excluded partner is a top
with first coordinate \(0\). For \(p\ge X'\), every
excluded partner is a top and, for \(e\ge1\),
\[
 \frac{a_s}{p^e}
 \le\frac{a+a/(kx^{13})}{X'}
       =b'+\frac{b'}{kx^{13}}
       \le b'+b'/x\le F_j(\operatorname{partner}(p,s,e)).
 \]
Thus all excluded positive parts vanish. Apply GLG2 to
the remaining terms to obtain
\[
 \sum_{s\in T_v}\{A_j(s)-3c\log_p(2a)\}
                                      \le g_v^0+d_v .
 \]
The original \(\operatorname{lo}\) is at least
\(3c\log_p(2a)\): its base term has \(C\ge3\), and all
its other branches are nonnegative. Sum over all
representatives and lifts and use (GLG.31).
If in addition \(p\) is inert and \(h_p=0\), every partner
is \((0,s)\), the preceding top bound makes all \(A_j(s)\)
zero, and the inequality follows directly despite the
gain's special zero definition. These cases exhaust the
original prime quantifier.
\hfill\(\square\)
% END GLG8

% BEGIN GLG9
\paragraph{GLG9. the full original remaining counting theorem.}
\label{claim:GLG9}
For every original \(\mathrm{Fm}\), every nonnegative \(K\),
and every \(C_{31}>0\), \(C_e,C_m\ge0\), define
\[
 C_1=\max\{C_*,3,\,2K.\mathrm{dl1}+C_e+1,\,
                                      4K.\mathrm{dl2}\},
                  \qquad C_{\rm count}=C_1+1.                    \tag{GLG.32}
\]
This constant is chosen before the instance families \(X,Y\).
For every such pair satisfying the full original
\(\mathrm{AdmFam}\), there is a common \(x_0\) such that
for every \(x\ge x_0\), every member \(i\), and every pair
of integer top functions in (GLG.2), the exact original
\(\mathrm{GainBoundAll}(X_{x,i},K,C_{\rm count},
                 b^{x^{-\epsilon/2}})\) holds.
Thus the complete source endpoint
\(\mathrm{RemainingCountingStatement}(\mathrm{Fm})\)
has been proved.

\emph{Proof.}
GLG7 supplies \(\mathrm{StarAll}\) at \(C_1\).
The admitted record at the fixed auxiliary exponent \(1\)
supplies \(\mathrm{adm9}\) and \(N_m+1\le w\).
Its \(\mathrm{adm9}\) clause says
\(e_p\le(C_e+1)fk^2\log_pa\) for \(p\le a\).
Since \(w\ge500\), \(X''=a/w\le a\), so GLG8 applies.

It remains to replace \(\mathrm{askP}\) by the actual
\(\mathrm{askG}\), not to assume they are equal. By (GLG.7),
GLG1, and \(a_s\le2a\),
\[
           A_j(s)\le A_j^G(s)
                \le A_j(s)+N_m\log_p(2a).                         \tag{GLG.33}
\]
This also proves the full \(\mathrm{RequiresAll}\) assertion
for arbitrary integer \(a_s\le2a\), \(b_s\ge0\): when
\(a_s<0\) both positive-part sums are zero.
Increasing \(C_1\) by \(1\) increases
\(\operatorname{lo}\) by at least
\((c+\eta)\log_p(2a)\ge N_m\log_p(2a)\), and increases
\(\operatorname{fl}\). Each additional branch coefficient
is nonnegative; the \(b_o\), \(b'/x\) and \(\eta\) terms
whose coefficient is \(1\) remain unchanged.
Since \(TTP(p)\subseteq TT\), (GLG.33) converts (GLG.29)
term by term into exactly the \(\mathrm{askG}\) inequality.
The gain bounds are unchanged. Take the maximum of the
two family thresholds. This preserves every quantifier
in the stated endpoint and the original error (GLG.23).
\hfill\(\square\)
% END GLG9

\subsection{Actual finite places and the comparison lattice}

% BEGIN GLG10
\paragraph{GLG10. served indices, coverage, and the original \(S_v\) bound.}
\label{claim:GLG10}
On the actual degree functions \(a_L,b_L\), for every prime
and every served representative,
\[
                         S_v\le 2mN_va/p .                        \tag{GLG.34}
\]
The original served-set coverage condition holds identically.
For the original family, \(\mathrm{ServedGIdx}\) holds for every
member once
\[
              x\ge\max\{\mathrm{Fm}.x_1+1,\,
                                      4^{\,5/\mathrm{Fm}.ez}\}.   \tag{GLG.35}
\]

\emph{Proof.}
A served representative occurs only at a candidate prime, so
\(p<X'\). By definition
\[
 S_v=\left\lceil\sum_{j=1}^m
       \max\{0,g_v^U-\Pr_v-\max(g_v^0,0)\}\right\rceil .
 \]
GLG1 and the nonnegativity argument of GLG8 give
\[
 \max\{0,g_v^U-\Pr_v-\max(g_v^0,0)\}
       \le\gamma_p\sum_{T_v}\alpha_s\le\frac85N_va/p.
 \]
If \(N_v=0\), \(g_v^U=-d_v\le0\), and \(S_v=0\).
Otherwise \(N_v\ge1\), and \(a/p>b'\ge w\ge500\).
The ceiling costs at most \(1\), which is bounded by
\((2-8/5)mN_va/p\). This proves (GLG.34) with the
original coefficient \(2\), including an empty-top orbit.
If a prime is a candidate and its total deficit is at most
\(x^{-\epsilon}\), the definition in (GLG.30) makes its
served set the entire good-representative set. This is the
exact coverage assertion.

To check the local index, a nonempty served set implies a
candidate inert prime and hence a positive even Frobenius
order. The original arithmetic Frobenius identification
therefore gives the nonzero host component required by
\(\mathrm{GIdx}\). The original level-prime condition says
that level primes split, so the inert prime is not a level
prime. Under (GLG.35) the supplied split bound is at least
\(\lfloor x^{ez/5}\rfloor\ge4\). An inert prime below \(5\)
would be \(2\), which is in the specified small-prime set,
or \(3\), which splits by this bound; both are impossible.
Thus \(p\ge5\). Finally, if \(p\ne T_u\), primality of
\(p,T_u\) implies \(p\nmid T_u\). These are exactly the
fields of \(\mathrm{GIdx}\), for every word \(u\); the
threshold precedes all words and finite places.
\hfill\(\square\)
% END GLG10

% BEGIN GLG11
\paragraph{GLG11. the finite class sum and the original global loss.}
\label{claim:GLG11}
Assume the original \(\mathrm{ParametersFill}\) identifications,
including \(\Lambda=L.\mathrm{LaG}\). Then the actual finite
von Mangoldt class sum is exactly \(A_j(s)\), not merely
asymptotic to it. For every real \(C,b_o\), every prime \(p\),
and \(C'=4|C|+1\),
\[
 \operatorname{lo}(C,b_o,p)\le\operatorname{loG}(C',b_o,p),
 \qquad
 \operatorname{fl}(C,p)\le\operatorname{flG}(C',p).                \tag{GLG.36}
\]
Here the global loss is literally
\[
 \begin{split}
 \operatorname{loG}(C',b_o,p)
 ={}&C'(S_c+k+\eta)\log_p(2a)+P_G(p)+O_G(p),\\
 P_G(p)={}&
 \begin{cases}
 C'rb'\log x,&p\text{ inert},\ p<p_0,\\
 C'\{b'/x+\eta+S_c+(w+t)\log x/\delta\}+b_o,
                          &p\text{ inert},\ p_0\le p<X',\\
 0,&\text{otherwise},
 \end{cases}\\
 O_G(p)={}&
 \begin{cases}
 C'\log x\{\Lambda a/(pk)+ap^{-L_*(p)}\},
                          &p\text{ inert},\ p<p_0,\ p<X',\\
 C'\log x\{\Lambda a/(pk)+a/p^2\},
                          &p\text{ inert},\ p_0\le p<X',\\
 0,&\text{otherwise},
 \end{cases}\\
 L_*(p)={}&\left\lceil\log_p(k^2x^{78})\right\rceil+1,\qquad
 \operatorname{flG}(C',p)=C'fN_m^2\log_pa+C'f(\log_pa)^2 .
 \end{split}                                                       \tag{GLG.37}
\]

\emph{Proof.}
The supplied filling identity identifies the summand of
the finite class sum with
\((a_s/p^e-F_j(\operatorname{partner}(p,s,e)))_+\) for
odd \(e\), and with zero for even \(e\). Its range ends
at \(\operatorname{Nat.log}_p(a_s)+2\). Beyond this range
\(p^e>a_s\); since the filling is at least \(1\), every
omitted positive part is zero. The top degree is a positive
integer by (GLG.2), so the natural-log convention causes
no exceptional case.

For (GLG.36), \(C'\ge1\), \(C'\ge C\), and
\[
              C(4S_c+\eta)\le C'(S_c+k+\eta).
 \]
All multiplying logarithms are nonnegative.
The small-prime \(P_G\) term follows from \(C'\ge C\).
In the middle range, \(C'\ge1\) absorbs \(b'/x+\eta\),
\(C'\ge C\) absorbs \(Cw\log x/\delta\), and the
additional \(S_c\) and \(t\) terms are nonnegative.
The coefficient of \(b_o\) stays \(1\), even if \(b_o<0\).
When \(p<p_0\),
\[
 \log_p(k^4x^{104})-\log_p(k^2x^{78})
                       =\log_p(k^2x^{26})>1,
 \]
so \(L_*(p)\le L_1(p)\) and
\(ap^{-L_1(p)}\le ap^{-L_*(p)}\).
This proves the tail comparison; if \(C<0\), first use
\(C\le C'\) on its nonnegative tail, and then the
nonnegative coefficient \(C'\) on this exponent comparison.
The remaining prime branches are zero on both sides.
The formula for \(\operatorname{flG}\) proves its inequality
directly. This verifies the exact envelope for arbitrary
real \(C\), rather than imposing a new positivity hypothesis.
\hfill\(\square\)
% END GLG11

% BEGIN GLG12
\paragraph{GLG12. the actual local depth inclusion.}
\label{claim:GLG12}
Let \(w\mid p\) be served and \(p\ne T_u\), and use the
original \(\mathrm{LatCmp}\) and \(\mathrm{NuServedDef}\).
Then
\[
 \mathrm{LocLe}\!\left(w,\lfloor e_w\tau_p\rfloor,
                              \mathrm{NuLat}_u,M_{{\rm ram},u}\right)
                                                                    \tag{GLG.38}
\]
holds, where \(e_w=v_w(p)\). Consequently the exact
\(\mathrm{NuServed}\) predicate is supplied by these original
comparison data, with no rational-operator assumption.

\emph{Proof.}
Write \(K_0\) for the global coefficient field,
\(A=\mathcal O_{K_0,w}\), \(O\subset F\) for the comparison
DVR and its fraction field, \(\iota:K_0\hookrightarrow F\)
for the supplied embedding, and \(\phi\) for the supplied
semilinear comparison. By \(\mathrm{LatCmp}\),
\[
 \iota(z)\in O\ \Longleftrightarrow\ z\in A,\qquad
 M=O\phi(L_u),
 \]
and the images of a \(K_0\)-basis remain independent over \(F\).
These conditions imply the following descent fact:
if \(v\in V_u\) and \(\phi(v)\in O\phi(L_u)\), then
\(v\in A L_u\). To see this, choose an \(A\)-basis of
the finite free localization \(A L_u\). Its images are
independent over \(F\). The unique coordinates of
\(\phi(v)\) in this basis belong to \(O\), hence the
coordinates of \(v\) belong to \(A\). Clearing their
finitely many denominators produces \(s\in\mathcal O_{K_0}\),
\(s\notin w\), with \(sv\in L_u\).

Now take \(z\in\mathrm{NuLat}_u\).
The served defining identity gives
\(\phi(\nu_{\rm in}z)\in P\), where
\[
                  P=J\cap p^T M\cap U_F,\qquad
                  T=\max(0,\lceil\tau_p\rceil).
 \]
Thus \(\phi(\nu_{\rm in}(p^{-T}z))\in M\).
By the descent fact there is \(s\notin w\) such that
\(s\nu_{\rm in}(p^{-T}z)\in L_u\). It is already in the
rational ramified subspace, since \(\nu_{\rm in}\) has
that image. Hence \(sp^{-T}z\in M_{{\rm ram},u}\).
As \(p\in w^{e_w}\) and
\(\lfloor e_w\tau_p\rfloor\le e_wT\),
\[
          sz=p^T(sp^{-T}z)
                     \in w^{\lfloor e_w\tau_p\rfloor}
                                          M_{{\rm ram},u}.
 \]
This is the exact denominator-cleared definition of
\(\mathrm{LocLe}\). In particular the factor is \(e_w\),
not the valuation of \(p\) in the larger comparison DVR.
\hfill\(\square\)
% END GLG12

% BEGIN GLG13
\paragraph{GLG13. the original local operators preserve both lattices.}
\label{claim:GLG13}
For every original \(G=GAtP(p)\) and each local vertex \(v\),
\[
                        y_vM_{\rm na}\subseteq M_{\rm na},
                  \qquad y_vM\subseteq M .                        \tag{GLG.39}
\]
These are the two exact fields of
\(\mathrm{FormalGroupSystem.adm12}\), for the original
dual-Weyl and Weyl integral coefficient models.

\emph{Proof.}
We first verify the elementary integral step, including the
Weyl model for which preservation cannot be inferred merely
from containment in the dual-Weyl model. Let \(\lambda\) be
decreasing of length \(q\), \(e=\lambda_q\), and
\(\nu_i=\lambda_i-e\ge0\). Let \(z\ne0\) in \(O\), and
\[
                 D_r(z)=\operatorname{diag}
                         (\underbrace{z^{-1},\ldots,z^{-1}}_r,
                                           1,\ldots,1).
 \]
The tableau basis in CB2 is homogeneous for row degrees.
In every one of its monomials the total degree in the
first \(r\) rows is at most \(N=\sum_{i\le r}\nu_i\):
each tableau column has distinct row entries, so it uses
at most \(r\) of those entries, and summing over columns
gives \(\sum_i\min(r,\text{column height}_i)=N\).
Thus the corresponding diagonal substitution, multiplied
by \(z^N\), preserves the integral dual-Weyl module.

For the Weyl module use its defining integral generators:
words in row divided-power operators applied to the
highest-weight product of leading minors. Each generator
is homogeneous in row degrees. Indeed the coefficient
of \(T^h\) in a row substitution changes its two row
degrees by \(+h,-h\). The generators lie in the dual-Weyl
module: generic column invariance and column multidegree
are preserved by row substitutions, as they commute
with the column substitutions; comparing coefficients
of \(T\) proves this for divided powers. Consequently a
nonzero homogeneous generator has first-\(r\)-row degree
\(d\le N\), by the preceding support bound. Multiplication
by \(z^N\) followed by \(D_r(z)\) is exactly multiplication
of this generator by \(z^{N-d}\in O\). It therefore
preserves the Weyl module itself.

In the original signed-weight lattice the determinant
factor is \(\det(g)^{-e}\). Replacing \(g\) by \(gD_r(z)\)
adds the factor \(z^{re}\). Since
\[
                 \sum_{i\le r}\lambda_i=N+re,
 \]
the resulting statement for either model is precisely
\[
 z^{\sum_{i\le r}\lambda_i}L_\lambda(\operatorname{col}_O g)
       \subseteq L_\lambda(\operatorname{col}_O(gD_r(z))).        \tag{GLG.40}
\]
There is no assumption that \(e\) is nonnegative.
If two basis matrices give the same lattice, they differ
by a matrix in \(\mathrm{GL}_q(O)\). The dual-Weyl model
is stable under row changes. The Weyl model is stable too:
elementary unipotent row changes are finite sums of the
defining divided powers, diagonal units act by integral
units on homogeneous generators, and these generate
\(\mathrm{GL}_q(O)\) by elimination over a local ring.
The determinant correction is a unit even for negative
\(e\). Thus (GLG.40) holds for the actual basis-independent
models \(H_0\) and \(\mathrm{Dl0}\), as well as for a
chosen basis.

For an original edge \(\Lambda\to\Lambda'\), the elementary
divisor theorem over its unramified coefficient DVR gives
a basis in which the target is obtained by multiplying
the first \(r\) basis vectors by \(p^{-1}\). Indeed
\(p\Lambda'\subseteq\Lambda\subseteq\Lambda'\) and the
specified quotient dimension is \(r\); the elementary
divisors are therefore \(p\) in exactly \(r\) positions
and \(1\) elsewhere. In the actual comparison ring the
specified element \(z\) satisfies \(z^2=p\). Each Frobenius
twist fixes \(p\), so the extended edge consists of two
successive steps \(D_r(z)\), with the original midpoint
between them. Apply (GLG.40) twice in every tensor factor.
The product of the factors is
\[
 \prod_s z^{2\sum_{i\le r}\lambda_{s,i}}
                         =p^{\sum_s\sum_{i\le r}\lambda_{s,i}}
                         =p^{G.P.\mathrm{ArV}(v)}.               \tag{GLG.41}
\]
This is exactly the normalization of the original \(y_v\).

Finally \(y_vF\) at a lattice tuple is this factor times
the finite sum of \(F\) over incoming neighbors in the
\(v\)-coordinate. Equation (GLG.41) places each summand
in the target tensor lattice; all other coordinates stay
fixed. If a tuple is invalid, each possible incoming
term is zero, since an incoming edge has lattice endpoints
and cannot repair an invalid different coordinate.
The edge relation is equivariant under the given level
group, so this sum preserves its prescribed character
equivariance. These are all three defining conditions of
the actual function lattice, for both coefficient models.
This proves (GLG.39) over the given DVR, without extension
or change of the rank or of \(\mathrm{ArV}\).
\hfill\(\square\)
% END GLG13

% BEGIN GLG14
\paragraph{GLG14. coefficient rationality supplies the literal \(\mathrm{NPDesc}\).}
\label{claim:GLG14}
For the original \(\mathrm{LatCmp}\), retain the original
\(\mathrm{GatTies}\), and assume the original
\(\mathrm{ValRat}\): at each genuine lattice tuple every
coefficient of \(\phi(v)\) is in \(\iota(K_0)\).
Then every original \(y_v\) descends to a \(K_0\)-linear
operator on the same global vector space. Under the
original \(\mathrm{VramDef}\) this proves
\(\mathrm{NPDesc}\) for the original selected lattice.
In particular the original
\(\mathrm{GroupCompatibilityFamily}\) supplies eventual
\(\mathrm{NPDescF}\); its comparison record is unchanged.

\emph{Proof.}
Here is the required rational descent, with no Galois
assumption on \(F/K_0\). Choose a \(K_0\)-linear retraction
\(\pi:F\to K_0\) of \(\iota\), by extending \(1\) to a
\(K_0\)-basis of \(F\). Polynomial coefficient evaluations
at all component indices and lattice tuples separate
the ambient function space. Suppose
\(H=\sum_i a_i\phi(v_i)\) lies in the \(F\)-span of the
comparison image and all those evaluations are rational.
Set \(v=\sum_i\pi(a_i)v_i\). If an evaluation of
\(\phi(v_i)\) is \(\iota(r_i)\), its value on \(\phi(v)\)
is \(\iota(\sum_i\pi(a_i)r_i)\); this is
\(\iota\pi(\sum_i a_i\iota(r_i))\), equal to the
rational value on \(H\). Separation gives \(\phi(v)=H\).
The comparison is injective by its stipulated independent
basis images, so this preimage is unique.

The apparent restriction in \(\mathrm{ValRat}\) to genuine
tuples causes no gap. Every element of \(M\) vanishes
at invalid tuples by definition. So does its \(F\)-span.
Since \(M=O\phi(L_u)\) and \(L_u\) spans the global space,
\(\phi(v)\) belongs to that span and vanishes there as well.
Thus all coefficient evaluations used above are rational.
By GLG13, \(y_v\) preserves the \(F\)-span of \(M\).
Its neighbor-sum formula and its rational scalar
\(p^{\mathrm{ArV}(v)}\) preserve rational coefficients.
The preceding argument therefore gives a unique preimage
of \(y_v\phi(v')\) for every \(v'\). Uniqueness and
linearity of \(y_v,\phi\) make this a \(K_0\)-linear
operator. This proves the entire \(\mathrm{YvRat}\)
predicate, including every local vertex.

Let \(\mathcal V\) be the actual served set and set
\[
 s(B)=\sum_{v\in\mathcal V\setminus B}S_v,\qquad
 T=\max(0,\lceil\tau_p\rceil),\qquad
 J=\sum_{B\subseteq\mathcal V}
           p^{s(B)}\left(\prod_{v\in B}y_v^m\right)M .            \tag{GLG.42}
\]
The \(S_v\) are already nonnegative integers. Products
use the very order in the original finite-set list; no
commutation assertion is necessary for rational descent.
The descended operators show that \(J\) is the
\(O\)-span of the images of rational vectors
\(p^{s(B)}\prod_{v\in B}(y_v^{\rm rat})^m z\), \(z\in L_u\).
Likewise \(p^TM\) is the extension of \(p^TL_u\), and
\(\mathrm{VramDef}\) identifies \(U_F\) with the extension
of the rational ramified subspace.

The original \(\mathrm{GatTies}\) supplies the fraction-field
and local coefficient-ring identifications. Together with
\(\mathrm{LatCmp.embW\_\allowbreak int}\), it gives the injective local
map \(A\to O\) used here. This map is faithfully flat:
it is torsion free over the DVR \(A\), hence flat, and
its maximal ideal remains proper. Under the independent
basis comparison, tensoring therefore commutes with
finite sums, images, and intersections, the last by
the kernel sequence for \(B_1\oplus B_2\to B_1+B_2\).
It also commutes with intersection with a rational
subspace, using its quotient vector space and the same
kernel argument. It follows that
\[
                   P=J\cap p^TM\cap U_F
 \]
is the extension of its rational preimage. This is the
literal \(\mathrm{NPDesc}\), not an identification of
arbitrary commensurable lattices.

Finally \(\mathrm{GroupCompatibilityFamily}\) includes
the actual \(\mathrm{CmpTied}\) polynomial-value record:
at a lattice tuple its \(\mathrm{val\_\allowbreak cmp}\) identity
writes \(\phi(v)\) as coefficient extension of the
specified polynomial over \(K_0\). Thus it supplies
exactly \(\mathrm{ValRat}\). Its \(\mathrm{LatDefs.vram}\)
is the \(\mathrm{VramDef}\) just used. Taking its common
eventual threshold proves \(\mathrm{NPDescF}\) for the
same family. Existence of this compatibility record is
an explicit antecedent, not a consequence of rational
descent.
\hfill\(\square\)
% END GLG14

% BEGIN GLG15
\paragraph{GLG15. finiteness and the exact normalized length.}
\label{claim:GLG15}
Retain the actual comparison, \(\mathrm{GatTies}\),
\(\mathrm{LatDefs}\), \(\mathrm{NuDef}\), and
\(\mathrm{NuUnserved}\), and use GLG14 for \(\mathrm{NPDesc}\).
Then the original \(\mathrm{LatticeLengthsAtPlace}\)
holds at every finite place. In particular at every
served \(w\) with \(p\ne T_u\) there is an integer \(l\ge0\)
such that
\[
 \operatorname{length}_O((M\cap U_F)/P)=l,\qquad
              l=(L.\mathrm{lenMN}(u,w))\,v_O(p).                 \tag{GLG.43}
\]
At an unserved place with \(p\ne T_u\),
\(L.\mathrm{lenMN}(u,w)=0\).

\emph{Proof.}
The original full global column lattice \(L_u\) is finite;
its intersection with the rational ramified space is
finite over the noetherian integer ring. The comparison
and flat intersection calculation identify \(M\cap U_F\)
with its extension, so it is a finite \(O\)-module.
Let \(s=\sum_{v\in\mathcal V}S_v\). The \(B=\varnothing\)
term in (GLG.42) gives \(p^sM\subseteq J\); hence
\[
                  p^{s+T}(M\cap U_F)\subseteq P .                \tag{GLG.44}
\]
The element \(p\) is nonzero and is a nonunit of \(O\).
Indeed its image comes from the rational prime of the
original unramified DVR through an injective local map,
as required by \(\mathrm{GatTies}\). A finite module
annihilated by \(p^{s+T}\) has finite length over a DVR.
This proves finiteness rather than assuming the local
length to be an integer.

Now apply the proved comparison SR2 to these very lattices,
using GLG14 and the served defining identity.
If \(a_w=v_O(\pi_w)>0\), scalar extension multiplies each
Smith exponent by \(a_w\), while \(v_O(p)=a_we_w\).
The source's \(\mathrm{NuDef}\) normalizes
\(L.\mathrm{lenMN}(u,w)\) as the \(A\)-length divided
by \(e_w\). Consequently
\[
 \operatorname{length}_O((M\cap U_F)/P)
   =a_w\operatorname{length}_A((M_{{\rm ram},u})_w/
                                             (\mathrm{NuLat}_u)_w)
   =v_O(p)\,L.\mathrm{lenMN}(u,w),
 \]
which is (GLG.43). The basis and ramified-space
identifications are unchanged, so their ranks agree too.
If \(w\) belongs to the original support, has the original
\(\mathrm{GIdx}\) guard, and its served set is empty,
(GLG.42) has only the empty-subset term:
\(J=M\) and \(\tau_p=0\). Thus \(P=M\cap U_F\), and
the local quotient has length zero. The guard excludes
\(p=T_u\), so the original
\(\mathrm{NuUnserved}\) also gives
\(L.\mathrm{lenMN}(u,w)=0\). This proves the first
clause of \(\mathrm{LatticeLengthsAtPlace}\) in its
empty-served case as well.
At an unserved regular place the original
\(\mathrm{NuUnserved}\) gives equality of the two localized
global lattices, hence length zero. The place \(p=T_u\)
keeps its separate original definition and is not
asserted to satisfy this regular-place zero conclusion.
\hfill\(\square\)
% END GLG15

% BEGIN GLG16
\paragraph{GLG16. the actual selected tops give the full second admissibility field.}
\label{claim:GLG16}
Suppose the original family satisfies \(\mathrm{TopsF}\):
\(\mathrm{ChosenTT}\), and \(\mathrm{ChosenBad}(C_b)\)
for one \(C_b>0\). Then the literal \(\mathrm{Adm2F}\)
holds, on the selected original subset \(TT\). If
\(x_a,x_b\) are the two supplied thresholds, one may take
\[
 x_0=\max\{x_a,x_b,\mathrm{Fm}.x_1+1,
                                     (16C_b+2)^{3/\epsilon}\}.    \tag{GLG.45}
\]

\emph{Proof.}
The chosen-top cardinality is already
\(\delta n/2\le|TT|\le2\delta n\).
Put \(T=x^{\epsilon/3}\). For any set \(\mathcal P\)
of positive integers less than \(\lceil e^T\rceil\),
\[
 \sum_{p\in\mathcal P}\frac{\log p}{p}
 \le2T\sum_{j=1}^{\lceil e^T\rceil}\frac1j
 \le2T\{1+\log\lceil e^T\rceil\}\le2T(T+2).                       \tag{GLG.46}
\]
The first step uses \(\log p\le T\) and
\(1/p\le2/(p+1)\); the harmonic bound follows by
integrating \(1/t\), and
\(\lceil e^T\rceil\le2e^T\) gives the last step.
In particular this applies to the original finite set
of large-order inert primes below \(X'\).
The \(\mathrm{ChosenBad}\) inequality says
\[
 \sum_{p\in\mathcal P}\frac{\log p}{p}
                       |\mathrm{badTops}(p)|
    \le4C_bx^{-2}\delta n
                    \sum_{p\in\mathcal P}\frac{\log p}{p}.
 \]
The actual \(\mathrm{BTT}\) sum is over a subset of
this set and has nonnegative summands. Therefore it is
at most \(8C_bx^{-2}\delta nT(T+2)\).
The threshold makes \(T\ge16C_b+2\), and hence
\(16C_b(T+2)\le T^2\). Thus
\[
 \mathrm{BTT}\le(\delta n/2)x^{-2}T^3
                      \le |TT|x^{-2+\epsilon}.
 \]
This is precisely the other clause of \(\mathrm{adm2}\).
The proof uses no asymptotic prime-counting theorem and
does not resample the selected top set.
\hfill\(\square\)
% END GLG16

\subsection{The original global installation and its remaining inputs}

% BEGIN GLG17
\paragraph{GLG17. the full literal finite-place installer.}
\label{claim:GLG17}
Fix the original \(D,L,X,Y,K\), with \(K\ge0\), and retain:
the actual \(\mathrm{GatTies}\), \(\mathrm{LatTied}\),
\(\mathrm{ParametersFill}\), \(\mathrm{NuDef}\),
the support condition containing every \(w\) with \(p\le X''\),
the actual length identifications of GLG15, and the following
two local-growth inputs:
\[
 \begin{gathered}
 \mathrm{MeanKappaBound}(X,K),\\
 \mathrm{admG1}(D,L,X,Y,K).
 \end{gathered}                                                     \tag{GLG.47}
\]
The latter is the original local-lattice assertion; its
length clause, on the same \(M,P\), is
\[
 \operatorname{length}_O((M\cap U_F)/P)
       \le v_O(p)\,d_u(\tau_p+\mathrm{io}_p).                     \tag{GLG.48}
\]
All its original additional clauses remain in the antecedent.
Assume \(\mathrm{ServedGIdx}\), and the original counting
conclusion at \(C,b_o\), with the ask comparison of GLG9.
Then
\[
                    \mathrm{LinkG\_\allowbreak at}(X,K,4|C|+1,b_o)           \tag{GLG.49}
\]
holds with all five original conjuncts.

\emph{Proof.}
At a served place with \(p\ne T_u\), GLG12 gives the
depth inclusion. GLG15 identifies the integer in
(GLG.48) with \(v_O(p)L.\mathrm{lenMN}(u,w)\).
The local-model rank is exactly \(d_u=D.\mathrm{duS}(u)\)
by \(\mathrm{GatTies}\), and \(v_O(p)>0\) by GLG15.
Division proves
\[
           L.\mathrm{lenMN}(u,w)\le d_u(\tau(w)+\mathrm{io}(w)).
 \]
The filling record identifies \(\tau,\mathrm{io}\)
with the original \(\tau_p,\mathrm{io}_p\).
If the served set is empty these are both zero by their
definitions, and the length is zero by GLG15. The depth
inclusion then has exponent zero and is exactly the
original regular-place inclusion in \(\mathrm{NuDef}\).
This proves the entire first conjunct, including its
empty-served-set implication.

GLG10 proves the second conjunct with constant \(2m\),
and proves the fourth conjunct, the exact coverage.
The fifth conjunct is the unchanged first input in
(GLG.47); GW2 supplies this input whenever its actual
canonical-deficit hypotheses are realized.

For the third conjunct take any signal or inert finite
place with \(p\le X''\). The support condition activates
the filling record's gain identity. Its top-subtype
bijection identifies \(L.\mathrm{TTw}(w)\) with the
original \(TTP(p)\); its coefficient identities
identify \(L.\mathrm{gain}(j,w)\) and \(\tau(w)\)
with \(\mathrm{gain}_j(p)\) and \(\tau_p\).
GLG11 identifies the finite class sum with \(A_j(s)\).
The inequality \(A_j(s)\le A_j^G(s)\), followed by
(GLG.36), makes the required global left side no larger
than the left side of the original counting conclusion.
The nonnegative factor \(m\) preserves the direction
when subtracting \(\operatorname{flG}\).
The gain bounds are exactly (GLG.28), so this proves
the whole third conjunct. Every conjunct concerns the
same supplied member, lattices, filling, degree functions,
and words; their witnesses have not been chosen separately.
\hfill\(\square\)
% END GLG17

% BEGIN GLG18
\paragraph{GLG18. the uniform family consequence and the exact frontier.}
\label{claim:GLG18}
Let \(\mathrm{Fm}\) be the original family. Suppose its
actual \(K=\mathrm{Fm}.GK\), \(X=\mathrm{Fm}.Gx\),
\(Y=\mathrm{Fm}.Gy\) satisfy the full original
\(\mathrm{AdmFam}\) for fixed \(C_{31}>0,C_e,C_m\ge0\);
retain \(\mathrm{GroupCompatibilityFamily}\) (or the
more precise \(\mathrm{NPDescF}\) it supplies), and retain
the exact eventual \(\mathrm{MeanKappaBound}\) on these
same instances. The \(\mathrm{admG1}\) field of the
admitted record supplies the original
\(\mathrm{FamilyFirstGrowthBound}\); when admissibility
is assembled from separate producers this field must
be supplied as well.
Then the complete literal \(\mathrm{LinkG}(\mathrm{Fm})\)
holds. One can use
\[
 C_G=4C_{\rm count}+1,\qquad
                   \delta_G(x)=x^{-\epsilon/2}.                  \tag{GLG.50}
\]
This is a proved implication on the original family;
it does not assert the existence of its retained input records.

\emph{Proof.}
GLG9 supplies one \(C_{\rm count}\ge0\) and the counting
inequalities at \(b^{\delta_G(x)}\), uniformly in the
member and all later choices. In particular they hold
at the actual \(a_L,b_L\), which satisfy (GLG.2).
The mean-deficit hypothesis is exactly the endpoint
proved in GW2 under its recorded original hypotheses:
the canonical \(\kappa_I\) identity in \(\mathrm{Fm.ex}\),
GI8 with its fixed coefficient bounds, and the small-prime
coefficient estimate. Those hypotheses are retained
when GW2 is used to discharge this input. No mean bound
for a new or independently selected deficit is being used.

The original family record already supplies, beyond
its starting threshold, \(\mathrm{GatTies}\), the actual
comparison and \(\mathrm{LatDefs}\), the filling identities,
\(\mathrm{LatTied}\), \(\mathrm{NuDef}\), the unserved
identity and the support condition. GLG14 supplies
\(\mathrm{NPDescF}\) from the stated compatibility
antecedent, and GLG15 supplies the required actual lengths.
GLG10 and GLG12 supply the finite-place conditions.
Take the maximum of their finitely many thresholds,
the admitted-family thresholds used in GLG9, the
stated mean-bound threshold and the admitted
first-growth threshold.
The constant \(C_G\) and function \(\delta_G\) are
already fixed before every later \(x,i,u,w,j\).
Apply GLG17 to every member at this common threshold.
Since \(\epsilon>0\), \(\delta_G(x)\to0\), giving
all quantifiers of the original \(\mathrm{LinkG}\).
\hfill\(\square\)

The finite local projection certificates are not a new
frontier in this statement: TC6 and EA9, specialized by
GI8/GII8 and transported by BC8 and GJ4, already prove them
under their original actual-model hypotheses. GLG13
adds the literal integral-operator supplier, GLG14
its rational descent, GLG16 the exact selected-top
admissibility consumer, and GLG9 the full previously
unexpanded logarithmic counting theorem. What remains
to assert the whole global leaf is the realization,
on this same family, of the retained compatibility,
the full admitted inputs, and the local first-growth
assertion (or its already reconstructed simultaneous
exceptional-direction inputs). No conditional estimate
has been relabeled as a construction of those data.
% END GLG18

\endgroup
\end{document}
